\documentclass[JEP,XML,SOM,Unicode, NoFloatCountersInSection,published]{cedram}
\datereceived{2023-09-18}
\dateaccepted{2024-09-18}
\dateepreuves{2024-09-30}

\usepackage[scr=rsfs,cal=euler]{mathalfa}
\graphicspath{{bellettini-et-al_figures}}
\multlinegap0pt
\hyphenation{solu-tion solu-tions explained respec-ti-vely para-me-tri-za-tion repara-me-tri-zation exam-ple exam-ples}
\newcommand\mto{\mathchoice{\longmapsto}{\mapsto}{\mapsto}{\mapsto}}

\DeclareMathOperator{\sym}{sym}
\DeclareMathOperator{\Sym}{Sym}

\makeatletter
\def\@settitle{%
\vspace*{-8mm}
\raggedleft\includegraphics[scale=.5]{titre-jep}
\vtop to 50 mm{%
\parindent=0pt
{\abstractfont\article@logo\par}

\hrule
\vfil
\begin{center}
\def\baselinestretch{1.2}\large\vfil
{\didottitraille\MakeUppercase\@title\par}
\vfil\vfil
\begin{minipage}{.8\textwidth}\centering
\ifx\@empty\smfbyname\else
{\smf@byfont\smfbyname\ifsmf@byauthor\enspace\else\ \fi}%
\fi {\smf@authorfont \edef\smfandname{{\noexpand\smf@andfont
 \smfandname}} \andify\authors\authors\par}
\end{minipage}
\vfil \vrule height .4pt width .3\textwidth \vfil
\end{center}}%
\par\enlargethispage{.7\baselineskip}%
}
\def\@setthanks{\def\thanks##1{\par##1\@addpunct{{\upshape.}}}\vspace*{-5pt}\thankses}
\makeatother

\usepackage[shortlabels]{enumitem}
\usepackage{yhmath}

\newcommand\C{\mathbb{C}}
\newcommand\NN{\mathbb{N}}
\newcommand{\R}{\mathbb{R}}
\newcommand\Z{\mathcal Z}

\renewcommand\H{\mathcal H}
\renewcommand\L{\mathcal L}
\newcommand{\M}{\mathcal{M}}
\newcommand{\admclass}{\mathcal W}
\newcommand{\largercomp}{\mathcal W}

\let\arc\wideparen
\let\curves\Sigma
\let\eps\varepsilon
\let\grad\nabla
\let\len\ell
\let\Om\Omega
\let\pointinspace P
\let\sigmamin\sigma
\let\psimin\psi

\newcommand{\weakstar}{\stackrel{*}{\rightharpoonup}}

\newcommand\admclassconv{{\mathcal W_\mathrm{conv}}}
\newcommand\compacts{\mathcal K}
\newcommand\convexcurves{\curves_\mathrm{conv}}
\newcommand{\di}{\,\mathrm{d}}
\newcommand{\Div}{\textup{div}}
\newcommand{\dist}{\textup{dist}}
\newcommand{\dOm}{\partial \Om}
\newcommand{\jump}[1]{\textup{\textlbrackdbl}{#1}\textup{\textrbrackdbl}}
\newcommand{\Lip}{\mathrm{Lip}}
\newcommand{\Liploc}{\Lip_\mathrm{loc}}
\newcommand{\nullb}{\partial^0}
\newcommand{\openannulus}{\Sigma_\mathrm{ann}}
\newcommand{\openhalfplane}{H^+}
\newcommand{\immclosedann}{\Phi({\overline\Sigma}_\mathrm{ann})}
\newcommand{\pointinOm}{P'}
\newcommand{\resa}{\mathop{\hbox{\vrule height 7pt width 0.5pt depth 0pt
\vrule height 0.5pt width 6pt depth 0pt}}\nolimits}
\newcommand{\resb}{\mathop{\hbox{\vrule height 4pt width 0.5pt depth 0pt
\vrule height 0.5pt width 3pt depth 0pt}}\nolimits}
\newcommand{\res}{\mathchoice{\resa}{\resa}{\resb}{\resb}}
\newcommand{\Xik}{{\mathcal W}_k}

\DeclareMathOperator{\supp}{supp}

\newtheorem{theorem}{Theorem}[section]
\newtheorem{definition}[theorem]{Definition}
\newtheorem{prop}[theorem]{Proposition}
\newtheorem{cor}[theorem]{Corollary}
\newtheorem{lemma}[theorem]{Lemma}
\newtheorem{ass}[theorem]{Assumption}
\newtheorem*{claim*}{Claim}
\newtheorem*{assertion*}{Assertion}
\newtheorem{assertion}{Assertion}
\renewcommand{\theassertion}{\Alph{assertion}}

\theoremstyle{definition}
\newtheorem{remark}[theorem]{Remark}
\newtheorem{Remark}[theorem]{Remark}
\newtheorem{esempio}[theorem]{Example}

\datepublished{2024-10-08}
\begin{document}
\frontmatter
\title{A non-parametric Plateau problem with\nobreakspace partial free boundary}

\author[\initial{G.} \lastname{Bellettini}]{\firstname{Giovanni} \lastname{Bellettini}}
\address{Dipartimento di Ingegneria dell'Informazione e Scienze Matematiche, Università di Siena,\\
53100 Siena, Italy,\\
\& International Centre for Theoretical Physics ICTP, Mathematics Section,\\ 34151 Trieste, Italy}
\email{giovanni.bellettini@unisi.it}
\urladdr{https://www3.diism.unisi.it/~bellettini/}

\author[\initial{R.} \lastname{Marziani}]{\firstname{Roberta} \lastname{Marziani}}
\address{Dipartimento di Ingegneria dell'Informazione e Scienze Matematiche, Università di Siena,\\
53100 Siena, Italy}
\email{roberta.marziani@unisi.it}
\urladdr{https://sites.google.com/view/roberta-marziani/}

\author[\initial{R.} \lastname{Scala}]{\firstname{Riccardo}\nobreakauthor \lastname{Scala}}
\address{Dipartimento di Ingegneria dell'Informazione e Scienze Matematiche, Università di Siena,\\
53100 Siena, Italy}
\email{riccardo.scala@unisi.it}
\urladdr{https://docenti.unisi.it/en/scala}

\thanks{We acknowledge the financial support of the GNAMPA of INdAM (Italian institute of high mathematics). The work of R. Marziani was partially supported by the Deutsche Forschungsgemeinschaft (DFG, German Research Foundation) under the Germany Excellence Strategy EXC 2044-390685587, Mathematics Münster: Dynamics--Geometry--Structure. GB and RS acknowledge the partial financial support of the PRIN project 2022PJ9EFL "Geometric Measure Theory: Structure of Singular Measures, Regularity Theory and Applications in the Calculus of Variations'', PNRR Italia Domani, funded by the European Union via the program NextGenerationEU, CUP B53D23009400006. RS also acknowledges the partial financial support of F-cur funding of the University of Siena (project number $\textup{2262-2022-SR-CONRICMIUR\_PC-FCUR2022\_002}$)}

\begin{abstract}
We consider a Plateau problem in codimension $1$ in the non-parametric setting, where a Dirichlet boundary datum is assigned only on part of the boundary $\partial \Omega$ of a bounded convex domain $\Omega\subset\mathbb{R}^2$. Where the Dirichlet datum is not prescribed, we~allow a free contact with the horizontal plane. We~show existence of a solution, and prove regularity for the corresponding area-minimizing surface. We~compare these solutions with the classical minimal surfaces of Meeks and Yau, and show that they are equivalent when the Dirichlet boundary datum is assigned on at most $2$ disjoint arcs of $\partial \Omega$.
\end{abstract}

\subjclass{49J45, 49Q05, 49Q15, 28A75}

\keywords{Plateau problem, area functional, minimal surfaces, relaxation, Cartesian currents}

\altkeywords{Problème de Plateau, fonctionnelle d'aire, surfaces minimales, relaxation, courants cartésiens}

\alttitle{Un problème de Plateau non paramétrique avec condition au bord partiellement libre}

\begin{altabstract}
Nous considérons un problème de Plateau en codimension $1$ dans un cadre non paramétrique, où une donnée de Dirichlet n'est assignée que sur une partie de la frontière $\partial\Omega$ d'un domaine convexe borné $\Omega\subset\mathbb{R}^2$. Là où la donnée de Dirichlet n'est pas prescrite, nous autorisons un contact libre avec le plan horizontal. Nous montrons l'existence d'une solution, et prouvons la régularité de la surface minimale correspondante. Nous comparons ces solutions avec les surfaces minimales classiques de Meeks et Yau, et montrons qu'elles sont équivalentes lorsque la donnée de Dirichlet est assignée sur au plus $2$ arcs disjoints de $\partial \Omega$.
\end{altabstract}

\maketitle
\vspace*{-\baselineskip}
\tableofcontents
\mainmatter

\section{Introduction}\label{sec:introduction}

The Plateau problem is a classical problem in the Calculus of Variations modeling configurations of soap films obtained by immersing a wire frame into soapy water. Roughly speaking, it consists in seeking for an area minimizing surface over all surfaces with prescribed boundary a given closed Jordan curve in space.
Over the years several approaches and variants were proposed, each corresponding to a specific choice of the class of admissible surfaces. In the following we list just few of them and we refer for example to \cite{Harrison-Pugh} and references therein for a list of the main approaches available in the literature.
One of the first result is due to Weierstrass and Riemann who studied a non-parametric Plateau problem in $\R^3$
obtained by minimizing the area over all Cartesian surfaces; this gave rise to the theory of minimal surfaces.
Successively Douglas and Radó developed independently \cite{Douglas,Rado} the classical parametric approach for disk type solutions. This method was later generalized by Jost \cite{Jost} to study the Plateau problem for surfaces with higher genus (see also
the paper \cite{MY} by Meeks and Yau).
A more general approach which accounts for a large class of surfaces was instead proposed by Federer and Fleming \cite{FedererFleming},
based on integral currents. Another remarkable work is due to Reifenberg \cite{Reifenberg} which adopts completely different techniques involving the concept of Čech homology. Relevant is also Almgren's contribution with three different approaches, one of these using the notion of varifolds \cite{Almgren65}.
Among all possible variants one might consider a partial free boundary version of the Plateau problem where the boundary datum is partially
fixed and partially free to move within a given surface.
This type of problem has been exhaustively studied (see for instance \cite{DHS}) in the parametric framework but never investigated,
to our best knowledge, with the non-parametric approach. To this aim, in the present paper we will analyze existence and regularity of solutions of a non-parametric partial free boundary Plateau problem.
More precisely,
we look for an area-minimizing surface
which can be written as a graph over a bounded open convex set
$\Omega\subset \R^2$,
and spanning
a Jordan curve $\Gamma_\sigma=\gamma\cup \sigma \subset \R^2 \times [0,+\infty)$ that is partially fixed.
Namely, $\gamma$ is fixed (Dirichlet condition) and
is given by a family
$\{\gamma_i\}_{i=1}^n\subset\partial \Om\times [0,+\infty)$
of $n\in\mathbb{N}$
curves
each joining distinct pairs of points $\{(p_i,q_i)\}_{i=1}^n$
of $\partial \Om$.
Whereas $\sigma$, which represents
the {\it free boundary}, is an unknown and
consists of (the image of)
$n$ curves $\sigma_1,\dots,\sigma_n$
sitting in $\overline\Omega$, and
joining the endpoints of $\gamma$
in order that $\gamma\cup\sigma$ forms a Jordan curve $\Gamma_\sigma$
in~$\R^3$. We~assume that each $\gamma_i$
is Cartesian, \ie it can be expressed as the graph of a given nonnegative function $\varphi$
defined on a corresponding portion of
$\partial \Omega$. This allows to restrict ourselves
to the Cartesian setting, and to assume that the competitors
for the Plateau problem are expressed by graphs of functions $\psi$
defined on a suitable subdomain of $\Omega$ depending on $\sigma$;
see Figure \ref{figura3} when $n=3$.
A peculiarity of our problem is the presence of a free boundary.

The purpose of this paper is twofold. We~start
addressing the question of existence and regularity of solutions.
Our first main result
(Theorems~\ref{teo_main_intro},~\ref{teo:reduction_from_W_to_Wconv} and~\ref{teo:structure_of_minimizers}) asserts that there are always solutions (which can be degenerate, in the sense that they may consist of more than one connected component,
see the example of the catenoid below) and that, under suitable hypotheses on the boundary datum, there is at least one regular solution continuous up to the boundary. Next we compare our solutions with solutions to a parametric Plateau problem when $n=1,2$.
Roughly speaking, our second main result (Theorems~\ref{main-theo-2},~\ref{teo:the_disk_type_Plateau_problem_n=1} and~\ref{teo:the_annulus_type_Plateau_problem_n=2}) shows that any regular solution to our minimization problem is a minimal embedding in the sense of Meeks and Yau~\cite{MY}, and vice-versa.

\subsubsection*{Existence and regularity of solutions}
We describe here our main results with few details, referring to Section \ref{sec:preliminaries}
for the precise description of the mathematical framework.
We
fix $n\in\mathbb N$ and $2n$ distinct points
$p_1,q_1,p_2,q_2,\dots,p_n,q_n \in \partial
\Om$
in clockwise order,
and set $q_{n+1}:=p_1$. The relatively
open arc
of $\partial \Omega$ between the points $p_i$ and $q_i$ is noted by $\partial_i^D\Om$,
and the relatively open
arc between $q_i$ and $p_{i+1}$ by $\partial_i^0\Om$. We~fix a nonnegative
continuous function $\varphi\colon\partial\Om\to[0,+\infty)$
positive on $\partial^D\Om:=\bigcup_{i=1}^n\partial^D_i\Om$ and vanishing
on $\{p_i,q_i\}_{i=1}^n\cup\partial^0\Om$, where $\partial^0\Om:=\bigcup_{i=1}^n\partial^0_i\Om$. For every $i=1,\dots,n$, we~denote by $\gamma_i$ the graph of $\varphi$ over $\overline {\partial^D_i\Om}$
and we consider curves $\sigma_i\colon[0,1]\to\overline \Om$
with the following properties:
\begin{enumerate}[label=(\roman*),wide]
\item{\label{i'}} $\sigma_i$ is injective, $\sigma_i(0)=q_i$ and $\sigma_i(1)=p_{i+1}$, for all $i=1,\dots,n$;
\item{\label{ii'}}
$\textup{int}(E(\sigma_i))\cap \textup{int}(E(\sigma_j))= \emptyset$
for $i,j=1,\dots,n$, $i \neq j$, where int denotes the interior part.
\end{enumerate}
Note carefully that $\sigma_i$ and $\sigma_j$ are allowed
to partially overlap.

We suppose the graph of $\varphi$ over $\partial^D\Om$ to be a Lipschitz curve in $\R^3$ (see Figure \ref{figura3}). Finally we
set
\begin{equation}\label{def:E}
E(\sigma):=\bigcup_{i=1}^n E(\sigma_i),
\end{equation}
and define the two classes
\begin{align}
\curves&:=\big\{\sigma = (\sigma_1,\dots,\sigma_n)
\in (\Lip([0,1]; \overline\Om))^n \text{ satisfies \ref{i'}--\ref{ii'}}\big\},\label{Sigma}
\\
\mathcal X_\varphi&:=\{(\sigma,\psi)\in \Sigma\times W^{1,1}(\Omega):\psi=0\text{ a.e. in }E(\sigma)\text{ and }\psi=\varphi\text{ on }\partial^D\Omega\}.
\label{chi}
\end{align}

If $(\sigma,\psi)\in \mathcal X_\varphi$, then the graph of $\psi$ over $\Om\setminus E(\sigma)$ is a surface spanning the curve~$\Gamma_\sigma$. We~look for a pair
$(\sigma,\psi)$ minimizing the area of such surfaces, that is, we
want to find a solution to the minimum problem
\begin{align}\label{problem1}
\inf_{(\sigma,\psi)\in \mathcal X_\varphi}
\int_{\Omega\setminus E(\sigma)}\sqrt{1+|\nabla \psi|^2}~dx.
\end{align}

\begin{figure}
\begin{center}
\includegraphics[width=0.65\textwidth]{roby1.pdf}
\caption{An example of the setting (in 3D),
when $n=3$. On the boundary of the convex set $\Omega$ fix the
points $p_i$, $q_i$; the arc of $\partial \Omega$ joining~$p_i$ to~$q_i$ is $\partial_i^D\Om$, while the arc joining $q_i$ to $p_{i+1}$
is $\partial^0_i\Om$ ($p_4:=\nobreak p_1$). On~$\partial^D\Omega$ the Dirichlet boundary datum $\varphi$ is imposed, whose graph has been depicted. The dotted arcs are the free planar curves $\sigma_i$ joining~$q_i$ and $p_{i+1}$. }\label{figura3}
\end{center}
\end{figure}

We then prove the following result, accounting for existence and
regularity of solutions to \eqref{problem1}.

\begin{theorem}\label{teo_main_intro}
Let $\Omega$ be strictly convex. Then there exists a solution $(\sigma,\psi)\in \mathcal X_\varphi$ to~\eqref{problem1} such that $\psi$ is continuous on $\overline\Om$, analytic in $\Omega\setminus E(\sigma)$, and $\Om\cap\partial E(\sigma)$ consists of a family of mutually
disjoint analytic curves (joining $p_i$ and $q_j$ in some order). Moreover, each connected component of $E(\sigma)$ is convex.
\end{theorem}

We emphasize that convexity of $\Om$ is necessary (even for the classical non-parametric Plateau problem with no free boundary, existence of regular solutions is not guaranteed if $\Omega$ is not convex).
The proof of existence relies on direct methods; however, since the class $\mathcal X_\varphi$ is not closed under weak* convergence in $BV$, they cannot be applied directly to \eqref{problem1} but
rather to a suitable weak formulation. For this reason we replace $\mathcal X_\varphi$ in \eqref{chi} with a larger class $\mathcal W$ of admissible pairs, and relax accordingly the functional in \eqref{problem1}. We~set
\begin{equation}\label{def:admissible_class_bis}
\largercomp
:=\big\{(\sigma,\psi)\in \curves \times BV(\Omega): \psi=0\text{ a.e. in }E(\sigma)\big\}.
\end{equation}
The weak formulation consists in looking for solutions to the problem
\begin{equation}\label{min-in-W}
\inf_{(\sigma,\psi)\in\mathcal W}\mathcal F(\sigma,\psi),
\end{equation}
where $\mathcal F$ is the functional defined by
\begin{align}\label{def:relaxed_functional}
\mathcal F(\sigma,\psi)
&:=\int_\Om\sqrt{1+|\nabla\psi|^2}~dx+|D^s\psi|(\Om)-|E(\sigma)|+\int_{\partial\Omega}|\psi-\varphi|~d\mathcal H^1\nonumber\\
&\hphantom{:}=\int_{\Om\setminus E(\sigma)}\sqrt{1+|\nabla\psi|^2}~dx+|D^s\psi|(\Om)+\int_{\partial\Omega}|\psi-\varphi|~d\mathcal H^1,
\end{align}
with $D^s\psi$ the singular part of the measure
$D\psi$ and $|E(\sigma)|$ the Lebesgue measure of~$E(\sigma)$.
Observe that $\mathcal F(\sigma,\psi)$
equals the integral in \eqref{problem1} when $\psi\in W^{1,1}(\Om)$ attains the boundary value $\varphi$.
The existence of solutions to \eqref{min-in-W} is shown in two steps. In~the first step we prove existence of minimizers of $\mathcal F$ in a smaller
class $\admclassconv\subset\mathcal W$ of admissible pairs $(\sigma,\psi)$,
where compactness is easier and allows to make use of the direct method.
The class $\admclassconv$
accounts only for specific geometries of the free boundary $\sigma$,
namely, each set $E(\sigma_i)$ is required to be convex (see \eqref{def:admissible_class} for its precise definition).
In the second step we show,
by means of a convexification procedure, that every minimizer $(\sigma,\psi)\in\admclassconv$ is actually a solution to \eqref{min-in-W}.
Eventually we prove that there exists at least a minimizer $(\sigma,\psi)\in\admclassconv$ which satisfies certain regularity properties, and in particular is a solution to \eqref{problem1}.
The fact that, for minimizers, all connected components of $E(\sigma)$ are convex, is somehow a consequence of the maximum principle, \ie every minimal surface is contained in the convex hull of its boundary.
The existence and regularity of a solution to \eqref{min-in-W} are contained in Theorems \ref{teo:reduction_from_W_to_Wconv} and \ref{teo:structure_of_minimizers}
respectively, which in turn imply Theorem \ref{teo_main_intro}. We~stress that Theorems \ref{teo:reduction_from_W_to_Wconv} and \ref{teo:structure_of_minimizers} are actually stated in the more general case of
a convex planar domain $\Omega$. However, if $\Omega$ is convex but not strictly convex it may happen that a solution to \eqref{min-in-W} is ``less regular'', in the sense that $\psi$ may not achieve the boundary condition (as in the next example),
thus failing to be a solution to \eqref{problem1}.

\subsubsection*{The example of the catenoid}
Our prototypical example is given by (half of) the catenoid.
Consider a cylinder in $\R^3$ with basis a circle of radius $r$
and height $\ell$.
Choose Cartesian coordinates for which the $x_1x_2$-plane
contains the cylinder axis, and restrict attention to the half-space
$\{x_3\geq0\}$ as
in Figure \ref{figura1},
where $\Omega=R_\ell := (0,\ell)\times(-r,r)$ and $n=2$.
Write
\[
\partial \Omega=\overline{\partial^D_1\Om}\cup\partial^0_1\Omega\cup \overline{\partial^D_2\Om}\cup \partial^0_2\Omega,
\]
where
\[
\partial^D_1\Om=\{0\}\times(-r,r),\;\; \partial^0_1\Omega=(0,\ell)\times\{r\},\;\; \partial^D_2\Om=\{\ell\}\times(-r,r),\;\; \partial^0_2\Omega=(0,\ell)\times\{-r\}.
\]
On the Dirichlet boundary $\partial^D\Om=\partial^D_1\Om\cup\partial^D_2\Om$ we prescribe the continuous function $\varphi$ whose graph
consists of the two half-circles $\gamma_1$ and $\gamma_2$. The endpoints of $\gamma_1$ and $\gamma_2$ live on the
free boundary plane (the horizontal plane) and are $p_1=(0,-r)$, $q_1=(0,r)$, and $p_2=(\ell,r)$, $q_2=(\ell,-r)$, respectively. The free boundary
$\sigma$ consists of two curves~$\sigma_1$ and~$\sigma_2$
with endpoints $q_1,p_2$, and $q_2,p_1$, respectively,
constrained to stay in $\overline\Omega$. The concatenation of $\gamma=\gamma_1\cup\gamma_2$ and $\sigma$
forms a Jordan curve
\begin{align}\label{gammasigma}
\Gamma_\sigma=\gamma_1\cup\sigma_1\cup\gamma_2\cup\sigma_2 \subset
\R^3.
\end{align}
Therefore we look for an area-minimizer
among all Cartesian surfaces $S$ with boundary~$\Gamma_\sigma$
{\it keeping $\sigma$ free}, \ie we look for a solution to \eqref{problem1} for this specific geometry. In~this case a minimizing sequence $(\sigma_k,\psi_k)\subset\mathcal W$ of the weak formulation \eqref{def:relaxed_functional} tends (in~the sense of Definition \ref{def:conv}) to a minimizer $(\sigma,\psi)\in \admclassconv$ which allows for two different possibilities. If $\ell$ is small, $\sigma_1$ and $\sigma_2$ remain disjoint and $(\sigma,\psi)\in\mathcal X_\varphi$. In~particular, the area-minimizing surface $S$ (given by the graph of $\psi$ over $\Om\setminus E(\sigma)$) is the classical (half) catenoid (namely the intersection between the catenoid and the half-space $\{x_3\geq0\}$). If instead $\ell$ is large, the two curves $\sigma_1$ and $\sigma_2$ merge, the region $\Om\setminus E(\sigma)$ collapses (\ie it reduces to the two segments $\partial_1^D\Omega\cup\partial_2^D\Omega$) and $\psi=0$ and therefore $(\sigma,\psi)\notin\mathcal X_\varphi$. In~particular, the surface $S$ is
the union of two vertical (half) disks. We~emphasize that this example is classical and, due
to the rotational symmetry of the curve $\Gamma$, it can be
reduced to a $1$-dimensional problem (see \cite{Greco,Bevi}).

Let us now quickly describe the second part of the paper.

\begin{figure}
\begin{center}
\includegraphics[width=0.5\textwidth]{disegnoP5.pdf}
\caption{The setting for the catenoid: for $\ell$ large
enough (the basis of the rectangle)
the dotted curves $\sigma_1$ and $\sigma_2$ merge and the (generalized)
graph of $\psi$ reduces to two vertical half-circles on $\partial^D \Om =
\partial_1^D\Om\cup \partial_2^D\Omega$.
In this case $\partial^D \Om \subset \partial E(\sigma_1)
\cup \partial E(\sigma_2)$.
}\label{figura1}
\end{center}
\end{figure}

\subsubsection*{Comparison with embedded minimal surfaces}
We recall that $\gamma_i$ is the graph of the map $\varphi$ on $\overline{\partial_i^D\Om}$. We~consider $\sym(\gamma_i)$, namely the graph of $-\varphi$ on $\overline{\partial_i^D\Om}$, which is symmetric to $\gamma_i$ with respect to the plane containing $\Omega$. Setting $\Gamma_i:=\gamma_i\cup \sym(\gamma_i)$, this turns out to be a simple Jordan curve in $\R^3$, for all $i=1,\dots,n$.
Hence we can consider the classical Plateau problem for the curve $\Gamma:=\bigcup_{i=1}^n\Gamma_i$. In the case $n=1$ a solution is
an area minimizing disk-type surface $S$ spanning $\Gamma=\Gamma_1$. Whereas in the case $n=2$ a solution is either an annulus-type surface spanning $\Gamma=\Gamma_1\cup\Gamma_2$ or the
union of two disjoint
disks spanning $\Gamma_1$ and $\Gamma_2$, respectively.
Then the following result holds true:
\begin{theorem}\label{main-theo-2}
Let $\Om$ be strictly convex.
For $n\in\{1,2\}$ let $(\sigma,\psi)\in \mathcal X_\varphi$ be a minimizer as in Theorem \ref{teo_main_intro}. Let $S^+$ be the graph of $\psi$ over $\Om\setminus E(\sigma)$ and let $S^-$ be the symmetric of $S^+$ with respect to the plane containing $\Om$. Then the set $S=S^+\cup S^-$ is a solution to the classical Plateau problem associated to $\Gamma=\bigcup_{i=1}^n\Gamma_i$.
Vice-versa every solution $S$ to the classical Plateau problem associated to $\Gamma=\bigcup_{i=1}^n\Gamma_i$ is symmetric with respect to the plane containing $\Omega$. Moreover, $S^+:=S\cap \{x_3\geq0\}$ is the graph of~$\psi$ over $\Om\setminus E(\sigma)$ for some $(\sigma,\psi)\in \mathcal X_\varphi$, a minimizer as in Theorem \ref{teo_main_intro}.
\end{theorem}

The above theorem is rigorously stated in Theorems \ref{teo:the_disk_type_Plateau_problem_n=1} ($n=1$) and \ref{teo:the_annulus_type_Plateau_problem_n=2} ($n=2$) in the more general case of $\Om$ convex. In~particular, if $\Om$ is convex, we~prove that there is a correspondence between a regular solution to the weak formulation \eqref{min-in-W} and a solution to the classical Plateau problem (as in the example of the catenoid).
A relevant consequence of this equivalence is
that when the boundary closed curve $\Gamma$ is symmetric with respect to the plane containing $\Om$, and its
upper part
is Cartesian, then the same property holds for the corresponding
Meeks and Yau solution.

The proof of Theorem \ref{main-theo-2} for $n=1$ is not difficult,
whereas for $n=2$ it is considerably more complicated,
and requires several lemmas:
we strongly use the convexity of the domain $\Omega$,
which implies that the cylinder $\Omega\times \R$, whose boundary contains $\Gamma$, is convex, and so the existence results of Meeks and Yau \cite{MY} (see also Theorem \ref{myexistence}) are applicable.

The main steps of the proof are the following: if $S$ is a Meeks-Yau annulus-type minimal surface, we~perform a Steiner symmetrization of the $3$-dimensional finite perimeter set in $\Omega\times\R$ enclosed by $S$ to obtain a set (symmetric with respect to the plane containing $\Om$) whose boundary is an annulus-type minimal surface $\widetilde S$ spanning~$\Gamma$ which is symmetric and such that $\widetilde S^+:=\widetilde S\cap\{x_3\ge0\}$ is Cartesian.
In turn, using standard results on the case
of equality for the perimeter of a set and its symmetrization, we
show that the original surface $S$
was already symmetric with respect to the plane containing $\Omega$, so $S^+$ was already Cartesian, and the conclusion of the proof for $n=2$ is achieved. Note that the aim of Theorem \ref{main-theo-2}
is not to provide new examples of minimal surfaces; rather, it enlightens
(among other things) some interesting qualitative properties of the Meeks-Yau
solutions.
Due to the highly nontrivial arguments, we~have restricted our analysis to the cases $n\in\{1,2\}$, since a generalization to the case $n>2$ probably requires heavy modifications.
Indeed, some lemmas needed to prove Theorem \ref{teo:the_annulus_type_Plateau_problem_n=2} employ crucially the fact that $\partial^0\Om$ consists of just two connected components.
For this reason we leave the case $n>2$ for future investigations.

\begin{figure}
\begin{center}
\includegraphics[width=0.4\textwidth]{disegnoP7.pdf}
\caption{
A possible configuration of the sets $E(\sigma_i)$ in the case $n=3$.
On the (clockwise oriented) arcs $\arc{p_1q_1}=\partial^D_1\Omega$,
$\arc{p_2q_2}=\partial^D_2\Omega$, and $\arc{p_3q_3}=\partial^D_3\Omega$
the function $\varphi$ is prescribed and positive. On $\partial^0\Omega=\arc{q_1p_2}\cup\arc{q_2p_3}\cup\arc{q_3p_1}$ and on $E(\sigma)=E(\sigma_1)\cup E(\sigma_2)\cup E(\sigma_3)$ we prescribe $\psi=0$. The curves $\sigma_i$ joining $q_i$ to $p_{i+1}$ (with the corresponding set $E(\sigma_i)$) are indicated. On the dotted segment $\sigma_1$ and $\sigma_2$ overlaps with opposite orientations. On the dark
region $\Omega\setminus E(\sigma)$, $\psi$ is not necessarily null.}\label{figura2}
\end{center}
\end{figure}

\subsubsection*{Some motivation} The setting of our problem models a cluster of soap films which are constrained to wet a
given system of wires $\gamma$ emanating from a given free
boundary plane (representing a table, or a water surface, on which the soap films can freely move). Our results show that if
the system of wires describes the graphs of functions on $\partial \Om$
as above, then the (Meeks and Yau) solutions of the ``parametric''
Plateau problem are in fact Cartesian, and coincide with the solutions
obtained by the non-parametric approach. This result can be viewed as a generalization of the well-known theorem of Rad\'o stating that any minimal disk spanning a Jordan curve in $\R^3$ whose projection on a plane is a bijection with a convex Jordan curve is the graph of a function defined on the plane \cite{Rado2}.

However, the scope of this article goes beyond this generalization, and the solutions we look for are strongly related with the vertical parts of Cartesian currents arising in the analysis of the relaxation of the
non parametric area functional in dimension $2$ and codimension $2$. We~further comment on this
in Section \ref{sec:final_comments_and_open_problems} where we go more into details.

\subsubsection*{Structure of the paper} The paper is organized as follows.
In Section \ref{sec:preliminaries} we introduce the
setting of the problem in detail.
In Section
\ref{sec:reduction_from_W_to_Wconv}
we show how to reduce the minimum problem
from the wider class $\admclass$ to the class $\admclassconv$
(Theorem \ref{teo:reduction_from_W_to_Wconv}).
Next,
in~Section
\ref{sec:existence_of_minimizers_of_F_in_Wconv}
we prove the existence of minimizers in $\mathcal
W_{\textup{conv}}$. As a consequence, we~gain
the existence of minimizers in class $\mathcal W$ (Corollary \ref{cor:mininmathcalW}).
In Section
\ref{sec:regularity_of_minimizers} we study the regularity of minimizers. Specifically, we~state and prove Theorem \ref{teo:structure_of_minimizers}
which, together with Theorem \ref{teo:reduction_from_W_to_Wconv}, generalize Theorem \ref{teo_main_intro}. Theorem \ref{teo_main_intro}
follows from Theorem \ref{thm:existence}, Corollary \ref{cor:mininmathcalW}, and Theorem \ref{teo:structure_of_minimizers}.
Eventually, in Section \ref{sec:comparison_with_the_parametric_Plateau_problem:the_case_n=1,2}
we compare our solutions with the classical minimal surfaces spanning $\Gamma$. Here, as anticipated, we~restrict the analysis to $n=1,2$,
the case $n=2$ essentially giving rise to either a catenoid-type minimal surface, or two disk-type surfaces spanning $\Gamma_1$ and $\Gamma_2$. The main theorems here are Theorems \ref{teo:the_disk_type_Plateau_problem_n=1} and \ref{teo:the_annulus_type_Plateau_problem_n=2}.
In Section \ref{sec:final_comments_and_open_problems}
we briefly point out our motivations for the present study and some open problems.
The paper concludes with an appendix containing some rather classical
results on convex sets and Hausdorff distance, needed in
Section \ref{sec:regularity_of_minimizers}.

\subsubsection*{Acknowledgements}
We thank the anonymous referees and the editors for suggestions and hints which allowed us to substantially improve the paper.

\section{Preliminaries}\label{sec:preliminaries}

\subsection{Area of the graph of a $BV$ function}
\label{notation}
Let $U\subset \R^2$ be a bounded open set. For any $\psi\in BV(U)$ we
denote by $D\psi$ its distributional gradient, so that
\begin{equation*}
D\psi=\nabla \psi\mathcal{L}^2+D^s\psi,
\end{equation*}
where $\nabla \psi$ is the approximate gradient of $\psi$ and $D^s\psi$ denotes the singular part of $D\psi$. We~recall that the
$L^1$-relaxed area functional reads as \cite{Giusti:84}
\begin{align}\label{def:relaxed_area}
\mathcal A(\psi; U):=\int_U\sqrt{1+|\nabla\psi|^2}~dx+|D^s\psi|(U).
\end{align}
In what follows we denote by $\partial^*A$ the reduced boundary
of a set of finite perimeter $A\subset\R^3$ (see \cite{AFP}).
For any $\psi\in BV(U)$ we denote by $R_\psi\subset U$ the set of regular points of $\psi$, namely the set of points $x\in U$ which are
Lebesgue points for $\psi$, $\psi(x)$
coincides with the
Lebesgue value of $\psi$ at $x$, and $\psi$ is approximately differentiable at~$x$. We~define the subgraph $SG_\psi$ of $\psi$ as
\begin{equation*}
SG_\psi:=\{(x,y)\in R_\psi\times\R\colon y<\psi(x)\},
\end{equation*}
which is a finite perimeter set in $U\times\R$. Its reduced boundary
in $U\times \R$ is the generalized graph $\mathcal G_\psi:=\{(x,\psi(x))\colon x\in R_\psi\}$ of $\psi$,
which turns out to be $2$-rectifiable. If
$\jump{SG_\psi}\in\mathcal D_3(\R^3)$
denotes the integral current given by integration over $SG_\psi$ and
$\partial \jump{SG_\psi}\in\mathcal D_2(\R^3)$ is its boundary in the sense of currents, then
\[
\jump{\mathcal{G}_\psi}=\partial\jump{SG_\psi}\res(U\times\R),
\]
with $\jump{\mathcal{G}_\psi}$ the integer multiplicity $2$-current given by integration over $\mathcal{G}_\psi$ (suitably oriented; see \cite{GiMoSu:98} for more details).

\subsection{Setting of the problem}\label{subsec:setting_of_the_problem}
We fix $\Omega\subset\R^2$ to be an open bounded convex set (strict
convexity is not required) which will be our reference domain. Given two points $p, q \in \partial \Omega$
in clockwise order, $\arc{pq}$ stands for the relatively
open arc on $\partial \Om$ joining~$p$ and~$q$.

Let $n\in\mathbb N$, $n \geq 1$, and let $\{p_i\}_{i=1}^n$ be
distinct points on $\partial\Omega$ chosen in clockwise order;
we set $p_{n+1} := p_1$.
For all $i=1,\dots,n$ let $q_i$ be a point in
$\arc{p_ip_{i+1}}
\subset \dOm$. We~set
\begin{align}
\partial_i^D\Omega:=\arc{p_iq_i},\qquad\nullb_i\Omega:=\arc{q_ip_{i+1}}\qquad \text{ for }i=1,\dots,n,
\end{align}
and
\begin{equation}
\partial^D\Om:=\bigcup_{i=1}^n\partial_i^D\Om,\qquad \nullb\Om:=\bigcup_{i=1}^n\nullb_i\Om.
\end{equation}
Since $\partial_i^D\Om$ and $\nullb_i\Omega$
are relatively open in $\partial \Omega$,
so are $\partial^D\Omega$ and $\nullb\Omega$. It follows that $\partial
\Om$ is the disjoint union
\[
\partial\Omega=\bigcup_{i=1}^n\{p_i,q_i\}\cup \partial^D\Om \cup
\nullb\Om.
\]
We fix a continuous function
$\varphi:\partial \Omega\to [0,+\infty)$ such that
\begin{equation}\label{boundary-value}
\varphi = 0 \text{ on } \nullb\Omega
\qquad \text{and} \qquad
\varphi >0 \text{ on }\partial^D\Omega,
\end{equation}
see Figures \ref{figura1}, \ref{figura3}. We~will make a further regularity assumption on $\varphi$: we require that the graph $\mathcal G_{\varphi\res\partial^D_i\Om}=\{(x,\varphi(x)):x\in \partial^D_i\Om\}$ of $\varphi$ on $\partial^D_i\Om$ is a Lipschitz curve in $\R^3$, for all $i=1,\dots,n$.
\begin{Remark}\rm
The hypothesis $\varphi>0$ on $\partial^D\Om$ excludes from our analysis
the example in Figure \ref{figura1mezzo}. We~will further comment on this later on (see Section \ref{subsec:the_example_of_the_catenoid_containing_a_segment}); the presence of pieces of $\partial^D\Om$ where $\varphi=0$ brings to some additional technical difficulties that we prefer to avoid here. However, the setting in Figure \ref{figura1mezzo} can be
achieved by an approximation argument. Namely, one considers a suitable regularization $\varphi_\eps$ of $\varphi$ on $\partial^D\Om$ such that $\varphi_\eps>0$, and then letting $\eps\to0$ one obtains a solution to the problem with Dirichlet datum $\varphi$.
\end{Remark}

\begin{remark} By definition \eqref{Sigma} any $\sigma\in \Sigma$ satisfies the injectivity property in~\ref{i'} which guarantees that the sets $E(\sigma_i)$
are simply connected (but not necessarily connected).
Assumption \ref{ii'} means essentially that
the curves $\sigma_i$
cannot cross
transversally each other, but might overlap. Notice that $\mathrm{int}(E(\sigma_i))$ might be empty,
the case $\nullb_i\Om=\sigma_i([0,1])$ being not excluded.
\end{remark}
In what follows we will study existence and regularity of solutions
to problem \eqref{min-in-W}. A first step
in this direction is to show in Section \ref{sec:reduction_from_W_to_Wconv}
that
\begin{equation}\label{equivalence-W-Wconv}
\inf_{(s,\zeta)\in \mathcal W}\mathcal F(s,\zeta)=
\inf_{(s,\zeta)\in \admclassconv}\mathcal F(s,\zeta),
\end{equation}
where $\mathcal{F}$ is the functional in \eqref{def:relaxed_functional} and
\begin{equation}\label{def:admissible_class}
\begin{split}
\admclassconv&:=\big\{(\sigma,\psi)\in \convexcurves\times BV(\Omega):
\ \psi=0\text{ a.e. in }E(\sigma)\big\},\\
\convexcurves&:=\big\{\sigma = (\sigma_1,\dots,\sigma_n)
\in \curves: E(\sigma_i) \text{ is convex for all } i=1,\dots,n\big\}.
\end{split}
\end{equation}
{ Notice that, by definition
\begin{equation}\label{eq:W_subset_hatW}
\convexcurves\subset\Sigma\quad\text{and}\quad \admclassconv\subset\mathcal W.
\end{equation}
}
Moreover, we~already know that the sets $\textup{int}(E(\sigma_i))$
might be empty, since from assumption \ref{i'} in \eqref{Sigma} we cannot exclude that $\sigma_i$
overlaps $\nullb_i\Omega$: Recalling that $\Omega$ is convex,
by \ref{ii'} and the convexity of
each $E(\sigma_i)$, this can happen, only if $\arc{q_ip_{i+1}}$
is a straight segment.\footnote{We will show
that for a minimizer, $\sigma_i([0,1])$ cannot intersect
$\partial^D \Omega$ unless $\partial^D \Omega$
is locally a segment (Theorem \ref{teo:structure_of_minimizers}).}
Afterward, in Section \ref{sec:existence_of_minimizers_of_F_in_Wconv},
we prove the existence of $(\sigma,\psi)\in
\admclassconv$ which is a solution to \eqref{min-in-W} by showing that there exists a minimizer to
\begin{equation}\label{min-conv}
\mathcal F(\sigma,\psi)=
\inf_{(s,\zeta)\in \admclassconv}\mathcal F(s,\zeta).
\end{equation}
Eventually in Section \ref{sec:regularity_of_minimizers}
we prove existence of solutions to \eqref{min-conv} which belong to $\mathcal X_\varphi$.

\begin{remark}\label{rem:condition_P}
Exploiting the characterization of the boundaries of convex sets
given in Corollary \ref{lm:charact_convex_sets2} in the appendix, we~see that conditions \ref{i'}, \ref{ii'} and the convexity of $E(\sigma_i)$ for the curves in $\convexcurves$ imply the following:

\begin{enumerate}[label=(P),wide]
\item\label{(P)} Let $\sigma\in \convexcurves$; then for all $i=1,\dots, n$ there are an injective (non-relabeled) reparametrization of $\sigma_i$ in $[0,1]$, and a
nondecreasing function $\theta_i\colon[0,1]\to\R$
with $\theta_i(1)-\theta_i(0)\le2\pi$,
such that, setting $\gamma_i(t):=(\cos(\theta_i(t)),\,\sin(\theta_i(t)))$ for all $t\in[0,1]$, we~have
\begin{equation*}
\sigma_i(t)=q_i+\len(\sigma_i)\int_0^t\gamma_i(s)\,ds\quad \quad\forall t\in[0,1],
\end{equation*}
where $\len(\sigma_i)$ denotes the length of $\sigma_i$.
\end{enumerate}
\end{remark}

\section{ Reduction from \texorpdfstring{$\largercomp$ to $\admclassconv$}{W}}
\label{sec:reduction_from_W_to_Wconv}

The main result of this section is contained in Theorem
\ref{teo:reduction_from_W_to_Wconv} where we prove the equivalence given in \eqref{equivalence-W-Wconv}. The reason being that in minimizing the functional
${\mathcal F}$ on~$\largercomp$ one issue is that the class $\curves$ in \eqref{Sigma}
is not closed under uniform convergence, since
a uniform limit of elements
in $\curves$ needs not be formed by injective curves. {However, we~can always modify a minimizing sequence of curves to curves in $\convexcurves$, since the modification can be done decreasing the energy.}

The fact that the infimum of $\mathcal F$ over $\mathcal W$ coincides with that over $\admclassconv$ is due to the following geometric property: whenever a set $E(\sigma_i)$ is not convex, we~can always convexify it reducing the energy. The procedure of convexification is described in Lemmas \ref{lem:trace_estimate}, \ref{lem:reduction_of_energy_I}, and \ref{lem:reduction_of_energy_II}. Again, the convexification of $E(\sigma_i)$ is still contained in $\overline\Omega$ thanks to the convexity of $\Om$.

\begin{theorem}[Reduction from $\largercomp$ to $\admclassconv$]
\label{teo:reduction_from_W_to_Wconv}
For every $(s,\zeta)\in\largercomp$ there exists $(\sigmamin,\psimin)\in\admclassconv$ such that every connected component of $E(\sigmamin)$ is convex, and
\begin{equation}\label{convexification}
{\mathcal{F}}(\sigmamin,\psimin)\le{\mathcal{F}}(s,\zeta).
\end{equation}
In~particular, \eqref{equivalence-W-Wconv} holds true.
Further, if the connected components of $E(\zeta)$ are not convex, then the
strict inequality holds in \eqref{convexification}.
\end{theorem}

\begin{remark}
Since the $\sigma_i$'s may overlap, the convexity of each
$E(\sigma_i)$ does not imply in general that every connected component of $E(\sigma)=\bigcup_{i=1}^nE(\sigma_i)$ is convex.
\end{remark}
For the reader convenience
we split the proof of Theorem \ref{teo:reduction_from_W_to_Wconv}
into a sequence of intermediate results: Lemmas \ref{lem:trace_estimate}, \ref{lem:reduction_of_energy_I}, \ref{lem:reduction_of_energy_II},
and the conclusion.
First we need to introduce some notation.

Let $(\sigma,\psi)\in \largercomp$. We~fix an extension $\widehat \varphi\in W^{1,1}(B)$
of $\varphi$ on an open ball $B\supset\overline\Omega$, where we recall that $\varphi$ is the boundary datum in \eqref{boundary-value}.
Extending $\psi$ in $B \setminus \overline\Omega$ as $\widehat \varphi$,
and still denoting by $\psi$ such an extension, we~can
rewrite $\mathcal F(\sigma,\psi)$ as
\begin{equation}\label{eq:equivalent_forumulation_of_F}
{\mathcal F}(\sigma,\psi)=\mathcal A(\psi;B)-|E(\sigma)|-\mathcal A(\psi;B\setminus \overline\Omega).
\end{equation}

\begin{lemma}[Trace estimate]\label{lem:trace_estimate}
Let $u\in BV(\R\times (0,+\infty))$ be a nonnegative function with compact support in an open ball $B_r\subset\R^2$. Then
\begin{equation}\label{eq:trace}
\int_{(\R\times\{0\})\cap B_r}u(s)\;d\mathcal H^1(s)\leq \mathcal A(u;B_r\cap (\R\times(0,+\infty)))-|E_{B_r}|,
\end{equation}
where
\[
E_{B_r}:=\{x\in B_{r}\cap (\R\times (0,+\infty)):u(x)=0\}.
\]
Moreover, inequality \eqref{eq:trace} is always strict, unless $u=0$ a.e. on $ \R \times (0,+\infty)$.
\end{lemma}

Notice that the
function $u$ is defined only on the half-plane $\R\times (0,+\infty)$, and
in~\eqref{eq:trace} the symbol $u(s)$ denotes
its trace on the line $\R\times \{0\}$ (which is integrable).

\begin{proof}

We denote by $x=(x_1,x_2)\in \R^2$ the coordinates in $\R^2$.
Set
\[
\openhalfplane := \R \times (0,+\infty),\quad
Z := (B_r\cap \openhalfplane) \times \R \subset \R^3.
\]
Let
\begin{equation*}
L_u:=\{(x,y)\in Z \colon x\in R_u,\ y\in(-u(x),u(x))\} \subset \R^3,
\end{equation*}
where $R_u$ is the set of regular points of $u$. We~have, recalling the notation in Section~\ref{notation},
\begin{equation}\label{eq:claim_step1*}
\begin{split}
2\mathcal{A}(u; B_r\cap\openhalfplane)&=\mathcal{A}(u; B_r\cap\openhalfplane)+\mathcal{A}(-u; B_r\cap\openhalfplane)
\\
&=\mathcal{H}^2(\partial^*(Z \cap SG_{u}))
+\mathcal{H}^2(\partial^*(Z \cap SG_{-u})
)\\
&=\mathcal{H}^2(Z \cap \partial^*L_u)+
2\vert E_{B_r}\vert.
\end{split}
\end{equation}
Write $B_r\cap(\R\times\{0\})=(a,b)\times \{0\}$. Then a slicing argument of the current $\jump{\mathcal{G}_u}$ yields
\begin{equation}\label{eq:claim_step2*}
\begin{split}
\mathcal{H}^2(Z \cap \partial^*L_u)&\ge\int_a^b
\mathcal{H}^1(Z \cap\{x_1=t\}\cap \partial^* L_u)dt\\
&=\int_a^b\mathcal{H}^1\bigl(Z\cap\{x_1=t\}\cap (\textup{spt}(\jump{\mathcal{G}_u}-\jump{\mathcal{G}_{-u}}))\bigr)dt\\
&\ge\int_a^b2u(t,0)dt
= 2\int_{(\R\times\{0\})\cap B_r}u(s)\;d\mathcal H^1(s),
\end{split}
\end{equation}
where the last inequality follows from the following fact: If we denote by $\jump{\mathcal G_u}_t$ the slice of the current $\jump{\mathcal G_{u}}$ on $\{x_1=t\}$, then
\[
\partial \jump{\mathcal G_{u}}_t=\delta_{(t,0,u(t,0))}-\delta_{(t,s_t,0)}
\qquad
\text{for~ a.e.~} t\in(a,b),
\]
where $s_t\ge0 $ is such that $(t,s_t)=B_r\cap (\{t\}\times \R^+)$, and in
writing $\delta_{(t,s_t,0)}$ we are using that $u$ has compact support in $B_r$.
This can be seen, for instance, by approximating\footnote{With respect to the strict convergence of $BV(
B_r \cap (\R\times\{0\}))$, which guarantees the approximation
also of the trace of $u$ on $\partial \big(B_r \cap (\R\times\{0\})\big)$.}~$u$ with a sequence of smooth functions.
Therefore
\begin{equation*}
\partial(\jump{\mathcal{G}_u}_t-\jump{\mathcal{G}_{-u}}_t)=\delta_{(t,0,u(t,0))}-\delta_{(t,0,-u(t,0))}
\qquad \text{for~ a.e.~} t\in(a,b).
\end{equation*}
This justifies the last inequality in \eqref{eq:claim_step2*} and, using \eqref{eq:claim_step1*}, the proof is achieved. Notice that, from the last formula, it follows that the last inequality in \eqref{eq:claim_step2*} is strict if $\jump{\mathcal{G}_u}_t-\jump{\mathcal{G}_{-u}}_t$ is not the straight segment connecting $(t,0,u(t,0))$ and $(t,0,-u(t,0))$ on a set of positive $\mathcal H^1$-measure. This implies that inequality in \eqref{eq:trace} is an equality if and only if $u=0$ a.e. on $H^+$.
\end{proof}

We now turn to two
technical lemmas needed to prove Theorem \ref{teo:reduction_from_W_to_Wconv}. We~introduce a class of sets whose boundaries are regular enough to support the trace of a $BV$ function.
Precisely we say that an open subset of $\R^2$ is \textit{piecewise Lipschitz} if
it can be written as the union of a finite family of (not necessarily disjoint) Lipschitz open sets.
{ Using that, for a Lipschitz set $E \subset \R^2$,
the symmetric difference $(\partial^*E)\Delta\partial E$ has null~$\mathcal H^1$ measure, one
can see\footnote{The conclusion
$\mathcal H^1((\partial^*V)\Delta\partial V)=0$ for a piecewise Lipschitz set
$V=\bigcup_{i=1}^mA_i$,
with $A_i$
Lipschitz open sets,
can be proved by induction on $m$,
using also the following fact: If $B_1$ and $B_2$ are open sets with $\mathcal H^1((\partial^*B_i)\Delta\partial B_i)=0$ for $i=1,2$, then $B=B_1\cup B_2$ satisfies $\mathcal H^1((\partial^*B)\Delta\partial B)=0$. This follows by the identity $\partial (B_1\cup B_2)=((\partial B_1)\setminus \overline{B_2})\cup ((\partial B_2)\setminus \overline{B_1})\cup ((\partial B_1)\cap \partial B_2)$, which shows that $\partial (B_1\cup B_2)$ is a $\mathcal H^1$-measurable subset of $\partial B_1\cup\partial B_2$.} that the same property
holds also for a piecewise Lipschitz set.} In~particular, by \eqref{def:relaxed_area} if $\Subset U$ is a piecewise Lipschitz subset of a
bounded open set $U\subset\R^2$, then
\begin{equation}\label{def:relaxed_area1}
\mathcal{A}(\psi;\overline V)=\mathcal{A}(\psi;V)+\int_{\partial V}|\psi^+-\psi^-|d\mathcal{H}^1,
\end{equation}
where $\psi^+$ (respectively $\psi^-$) denotes the trace of $\psi\res V$ (respectively $\psi\res(U\setminus\overline V)$) on~$\partial V$.

\begin{lemma}[Reduction of energy, I]
\label{lem:reduction_of_energy_I}
For $N\ge1$ let $F_1,\dots,F_N$ be
nonempty connected subsets of $\overline \Om$,
each $F_i$ being the closure of a piecewise Lipschitz
set, with $F_i \cap F_j= \emptyset$ for $i,j \in \{1,\dots,
N\}$, $i\neq j$.
Let $\psi\in BV(B)$ satisfy
\begin{equation}\label{eq:psi_0_union}
\psi=0\quad\text{a.e. in}\quad G:=\bigcup_{i=1}^NF_i\quad\text{and}\quad
\psi=\widehat{\varphi}\quad\text{a.e. in}\quad B\setminus{\Om}.
\end{equation}
Then, for any $i\in\{1,\dots,N\}$,
\begin{align}\label{stretta}
\mathcal A(\psi_i^\star;B)-|G_i^\star|-\mathcal A(\psi_i^\star;B\setminus \overline\Omega)\leq
\mathcal A(\psi;B)-|G|-\mathcal A(\psi;B\setminus \overline\Omega),
\end{align}
where
\begin{equation}\label{def:E_star,Psi_star}
G_i^\star:=\bigcup_{j\neq i}F_j\cup \mathrm{conv}(F_i)\quad\text{and}\quad
\psi_i^\star :=\begin{cases}
0&\text{in } \mathrm{conv}(F_i),\\
\psi&\text{otherwise.}
\end{cases}
\end{equation} Further, inequality in \eqref{stretta} is strict unless $\psi=\psi^\star_i$ a.e.
\end{lemma}

\begin{proof}
Fix $i\in\{1,\dots,N\}$.
By the convexity of $\Omega$, we~have $\psi=\psi_i^\star$ in $B\setminus \overline\Omega$, hence it suffices to show that
\begin{equation*}
\mathcal A(\psi_i^\star;B)-|G_i^\star|\leq\mathcal A(\psi;B)-|G|.
\end{equation*}
We start by observing that we may assume $F_i$ to be simply connected.
Indeed, if not, we~can replace it with the set obtained by filling the
holes of $F_i$, and by setting $\psi$ equal to zero in the holes.\footnote{If $H$ is a hole of $F_i$ and it happens that $F_j\subset H$ for some $j\neq i$, we~redefine $F_i$ as the union of it with $H$, and set $F_j=\emptyset$. This procedure does not invalidate the following argument.}
This procedure reduces the energy since $F_i$ is piecewise Lipschitz, and any hole $H$ of
it has the property that
the external trace of $\psi\res (B\setminus H)$ on $\partial H$ vanishes.

We have that $(\partial \mathrm{conv}(F_i))\setminus\partial F_i$ is a countable union of segments. We~will next modify $\psi$ by iterating at most countably many operations, setting $\psi=0$ in the region between each of these segments and $\partial F_i$.

\subsubsection*{Step 1: Base case} Let $l$ be one of such segments, and
$U$ be the open region enclosed between $\partial F_i$ and $l$. We~define $\psi'\in BV(\Om)$ as
\begin{equation*}
\psi':=\begin{cases}
0&\text{in}\ \ U,\\
\psi &\text{otherwise}.
\end{cases}
\end{equation*}
We claim that
\begin{equation}\label{eq:claim_bis*}
\mathcal{A}(\psi';B)-|G'|\le\mathcal{A}(\psi;B)-|G|,
\end{equation}
with strict inequality unless $\psi'=\psi$ a.e., where $G':= G \cup
\overline U$.
To prove the claim, we~introduce the sets
\[
H:=\mathrm{int}(F_i\cup U),\qquad
V:=U\cap(\bigcup_{j\neq i}F_j).
\]
Note that $H$ is a piecewise Lipschitz set.
By construction
\begin{equation*}
|G'|=|{H}|+ |\bigcup_{j\neq i}F_j|-|V|,
\end{equation*}
and \eqref{eq:claim_bis*} will follow if we show that
\begin{equation*}
\begin{split}
\mathcal{A}(\psi';B)-|{H}| \le \mathcal{A}(\psi;B)-|\bigcup_j F_j|+|\bigcup_{j\neq i}F_j|-|V|=\mathcal{A}(\psi;B)-|F_i\cup V|,
\end{split}
\end{equation*}
with strict inequality unless $\psi'=\psi$ a.e. in $\Om$.
Since $|H|=|F_i\cup V|+|U\setminus V|$, this can also be written as
\begin{equation*}
\mathcal{A}(\psi';B)
\le\mathcal{A}(\psi;B)+|U\setminus V|.
\end{equation*}
In turn $\mathcal{A}(\psi';B)=\mathcal{A}(\psi';\overline U)+\mathcal{A}(\psi';B\setminus \overline U)$ (and similarly for $\psi$), so we have reduced ourselves with proving
\begin{equation}\label{eq:claim_tris**}
\mathcal{A}(\psi';\overline U)
\le\mathcal{A}(\psi;\overline U)+|U\setminus V|.
\end{equation}
In view of the definition of $\psi'$ which is zero
in $U$, we~have\footnote{We use the precise integral formula \eqref{def:relaxed_area1} thanks to the boundary regularity of $U$, where we have $\partial U\setminus l\subset \partial F_i$.}
\[
\mathcal{A}(\psi';\overline U)=\int_l|{\psi}^+|d\mathcal H^1+|U|
\]
($\psi^+$ denoting the trace of $\psi\res (B\setminus U)$ on the segment $l$) implying that \eqref{eq:claim_tris**} is equivalent to
\begin{equation*}
\int_l|{\psi}^+|d\mathcal H^1
\le\mathcal{A}(\psi;\overline U)-|V|.
\end{equation*}
Finally, if $\psi_U$ denotes the trace of $\psi\res U$ on $l$, we~write $\mathcal{A}(\psi;\overline U)=\mathcal{A}(\psi;\overline U\setminus l)+\int_l|\psi^+-\psi_U|d\mathcal H^1$, and the expression above is equivalent to
\begin{equation}\label{eq:claim_tris*}
\int_l|{\psi}^+|d\mathcal H^1
\le\int_l|\psi^+-\psi_U|d\mathcal H^1+\mathcal{A}(\psi;\overline U\setminus l)-|V|.
\end{equation}
We now prove \eqref{eq:claim_tris*}.
Fix a Cartesian coordinate system $(x_1,x_2)$ so that $l$ belongs to
the $x_1$-axis and $U$ belongs to the half-plane $\{x_2>0\}$. Let $u$ be an extension of~$\psi$ in~$\R\times(0,+\infty)$ which vanishes outside $U$.
Lemma \ref{lem:trace_estimate},
applied to $u$ with the ball $B_r=B$, implies
\[
 \int_l|{\psi}_U|d\mathcal H^1=\int_{\{x_2=0\}\cap B}u \;d\mathcal H^1\leq \mathcal A(u;B\cap (\R\times (0,+\infty)))-|E_B|\leq \mathcal A(\psi;\overline U\setminus l)-|V|.
\]
Here the last inequality follows by recalling that $\psi$ (and thus $u$) vanishes on $V$.
From this and the inequality $\int_l|{\psi}^+|d\mathcal H^1\leq \int_l|\psi^+-\psi_U|d\mathcal H^1+\int_l|{\psi}_U|d\mathcal H^1
$ the proof of \eqref{eq:claim_tris*} is achieved, so that \eqref{eq:claim_bis*} follows. Notice that in applying Lemma \ref{lem:trace_estimate} the inequality holds strict when $\psi'$ does not coincide with $\psi$ a.e.

\subsubsection*{Step 2: Iterative case}
We set $\partial(\mathrm{conv}(F_i))
\setminus\partial F_i=\bigcup_{j=1}^\infty l_j$ with $l_j$ mutually disjoint segments.
For every $h\ge1$ we define the pair $(\psi_{h},G_{h})$ as follows:
\begin{itemize}
\item if $h=1$
\begin{equation*}
\psi_{1}:=\begin{cases}
0&\text{in}\ U_1,\\
\psi&\text{otherwise,}
\end{cases}\qquad \text{and}\quad
G_{1}:=G\cup \overline U_1,
\end{equation*}
where $U_1$ is the open region enclosed between $\partial F_i$ and $l_1$. We~also define $H_1:=\mathrm{int}(\overline{F_i\cup U_1})$;
\item if $h\ge2$
\begin{equation*}
\psi_{h}:=\begin{cases}
0&\text{in}\ U_h\\
\psi_{h-1}&\text{otherwise,}
\end{cases}\quad\text{and}\quad
G_{h}:=G_{h-1}\cup\overline U_h ,
\end{equation*}
where $U_h$ is the open region enclosed between $\partial H_{h-1}$ and $l_h$ and
\[
H_h:=\mathrm{int}(\overline{H_{h-1}\cup U_h}).
\]
\end{itemize}
By construction each $H_{h}$ is
simply connected and piecewise Lipschitz, $H_{h}\subset H_{h+1}$, $G_{h}\subset G_{h+1}\subset\overline\Om$ for every $h\ge1$, and moreover
\begin{equation}\label{limit_E*}
\lim_{h\to+\infty} |H_{h}|=|\mathrm{conv}(F_i)|,\qquad
\lim_{h\to+\infty} |G_{h}|=|G^\star_i|,
\end{equation}
where $G^\star_i:=\bigcup_{h=1}^\infty G_h=\bigcup_{j\neq i}F_j\cup \mathrm{conv}(F_i) $. For any $h\geq2$ we apply step 1,
and after~$h$ iterations we get
\begin{equation}\label{iterative_proc*}
\begin{aligned}
\mathcal{A}(\psi_{h};B)-|G_{h}|&\le \mathcal{A}(\psi_{h-1};B)-|G_{h-1}|\le\cdots\\
&\le \mathcal{A}(\psi_{1};B)-|G_{1}|\le \mathcal{A}(\psi;B)-|G|.
\end{aligned}
\end{equation}
In~particular,
\begin{equation*}
|D\psi_{h}|(B)\le\mathcal{A}(\psi_{h};B)\le \mathcal{A}(\psi;B)+|G_{h}\setminus G|\le \mathcal{A}(\psi;B)+|\Om\setminus G|,
\end{equation*}
for all $h\ge1$, and then we easily see that, up to a subsequence, $\psi_{h}\weakstar\psi^\star_i$ in $BV(B)$, where $\psi_i^*$ is defined as in \eqref{def:E_star,Psi_star}.
Now the lower semicontinuity of $\mathcal{A}(\cdot;B)$ yields
\begin{equation}\label{lsc_A*}
\liminf_{h\to+\infty}\mathcal{A}(\psi_{h},B)\ge
\mathcal{A}(\psi_i^\star;B).
\end{equation}
Finally, gathering together \eqref{limit_E*}--\eqref{lsc_A*} we infer
\begin{equation*}
\mathcal{A}(\psi_i^\star;B)-|G^\star_i|\le
\liminf_{h\to+\infty}\mathcal{A}(\psi_{h};B)-\lim_{h\to+\infty}|G_{h}|\le
\mathcal{A}(\psi;B)-|G|.
\end{equation*}
Again we have the strict inequality unless $\psi_h=\psi_{h-1}$ for all $h$ a.e. in $\Om$.
This concludes the proof.
\end{proof}

\begin{lemma}[Reduction of energy, II]
\label{lem:reduction_of_energy_II}
Let $N\ge 1$, $F_1,\dots,F_N,G$ and $\psi$ be
as in Lemma \ref{lem:reduction_of_energy_I}.
Then there exist $\tilde n \in \{1, \dots, N\}$ and mutually disjoint closed convex sets $\widetilde F_1,\dots,\widetilde F_{\tilde n}\subset\overline\Om$ with nonempty interior
such that
\begin{equation}\label{A}
G\subset \bigcup_{i=1}^{\tilde n}\widetilde F_i=:G^\star,
\end{equation}
and
\begin{equation}\label{B}
\mathcal{A}(\psi^\star;B)-|G^\star|-\mathcal{A}(\psi^\star;B\setminus \overline\Om)\le \mathcal A(\psi; B)-|G|-\mathcal{A}(\psi;B\setminus\overline{\Om}),
\end{equation}
where
\begin{equation}\label{C}
\psi^\star:=\begin{cases}
0&\text{in}\ G^\star,\\
\psi&\text{otherwise}.
\end{cases}
\end{equation}
Finally, inequality in \eqref{B} is strict unless $\psi=\psi^\star$ a.e.
\end{lemma}

\skpt
\begin{proof}\vspace*{-\baselineskip}

\subsubsection*{Base case}
If $N=1$ we set $\widetilde F_1:= \mathrm{conv}(F_1)=G^\star$ and the thesis follows by Lemma~\ref{lem:reduction_of_energy_I}. Suppose $N >1$. We~take the sets
\begin{equation}\label{sets_base_case}
\mathrm{conv}(F_1), F_2, \dots, F_N\quad\text{and}\quad G_1^\star:=\bigcup_{i=2}^NF_{i}\cup\mathrm{conv}(F_1),
\end{equation}
and let
\begin{equation*}
\psi_1^\star:=\begin{cases}
0&\text{in}\ G_1^\star,\\
\psi&\text{otherwise}.
\end{cases}
\end{equation*}
Then by Lemma \ref{lem:reduction_of_energy_I},
\begin{align}\label{base_case}
\mathcal A(\psi_1^\star;B)-|G_1^\star|-\mathcal A(\psi_1^\star;B\setminus \overline\Omega)\leq
\mathcal A(\psi;B)-|G|-\mathcal A(\psi;B\setminus \overline\Omega),
\end{align}
with strict inequality unless $\psi_1^\star=\psi$ a.e.

\subsubsection*{Iterative case}
Let $m,k,h$ be natural numbers such that $1\le k\le m\le N$ and $1<h\le 2N-1$, and let $F_{1,h},\dots,F_{m,h}$ be closed subsets of $\overline \Om$
with nonempty interior that satisfy the following property:
\begin{enumerate}[label=($\arabic*$),wide]
\item \label{1}$F_{1,h},\dots,F_{k,h}$ are convex;
\item\label{2} $F_{i,h}\cap F_{j,h}=\emptyset$ for all $i, j\ne k$, $i\ne j$, $i,j=1,\dots,m$.
\end{enumerate}
Notice that for $h=2$ and $m=N$ the sets
\begin{equation*}
F_{1,2}:=\mathrm{conv}(F_1),\ F_{2,2}:=F_2,\dots,\ F_{N,2}:=F_N,
\end{equation*}
satisfy \ref{1}, \ref{2} with $k=1$ by the base case (so the iterative step can be applied to these sets).

We then set $I_{k,h}:=\{1\le i\le m,\,i\ne k\colon F_{i,h}\cap F_{k,h}\ne\emptyset\}$.
If $I_{k,h}=\emptyset$ and $k=m$, we~are done, otherwise we
construct a new family of sets using the following algorithm, distinguishing the two cases \ref{case:a} and \ref{case:b}:
\begin{enumerate}[(a),wide]
\item\label{case:a} if $I_{k,h}=\emptyset$ and $k<m$
we define the sets
\begin{equation*}
F_{i,h+1}:=\begin{cases}
F_{i,h}& \text{for}\ i\neq{k+1},\\
\mathrm{conv}(F_{k+1,h})&\text{for}\ i=k+1,
\end{cases}\quad \text{ for }i=1,\dots,m,
\end{equation*}
and $
G_{h+1}^\star:=\bigcup_{i=1}^mF_{i,h+1}$;
\item\label{case:b}
if $I_{k,h}\ne\emptyset$,
up to relabeling the indices, we~may assume that
\begin{align*}
I_{k,h}=\{k_{h,1}\le i\le k_{h,2}\}\setminus\{k\},
\end{align*}
for some $k_{h,1}\ne k_{h,2}$ with $1\leq k_{h,1}\le k\le k_{h,2}\leq m$,
so that
\[
\{1,\dots,m\}\setminus\{k\}\setminus I_{k,h}=\{1\le i\le k_{h,1}-1\}\cup\{k_{h,2}+1\le i\le m\}.
\]
Note that if $k_{h,1}=1$ then $\{1\le i\le k_{h,1}-1\}=\emptyset$, and similarly if $k_{h,2}=m$ then $\{k_{h,2}+1\le i\le m\}=\emptyset$.
Then we set
\begin{equation*}
F_{i,h+1}:= \begin{cases}
F_{i,h}& \text{for}\ i=1,\dots,k_{h,1}-1,\\
\mathrm{conv}(F_{k,h}\cup(\bigcup_{j\in I_{k,h}}F_{j,h}))& \text{for}\ i=k_{h,1},\\
F_{i+k_{h,2}-k_{h,1},h}& \text{for}\ i=k_{h,1}+1,\dots, m-k_{h,2}+k_{h,1},
\end{cases}
\end{equation*}
and $G_{h+1}^\star:=\bigcup_{i=1}^{m-k_{h,2}+k_{h,1}}F_{i,h+1}$.
\end{enumerate}
In both cases \ref{case:a} and \ref{case:b} a direct check shows that the produced sets satisfy properties~\ref{1} and~\ref{2} with $m$, $k+1$, $h+1$ and $m-k_{h,2}+k_{h,1}$, $k_{h,1}$, $h+1$ respectively.

In both cases we also define the function
\begin{equation*}
\psi_{h+1}^\star:=\begin{cases}
0&\text{in}\ G_{h+1}^\star,\\
\psi_h^\star&\text{otherwise}.
\end{cases}
\end{equation*}
Then, by induction, for all $1<h\le 2N-1$ we use Lemma \ref{lem:reduction_of_energy_I}, and
in view of \eqref{base_case} we infer
\begin{align*}
\mathcal A(\psi_{h+1}^\star;B)-|G_{h+1}^\star|-\mathcal A(\psi_{h+1}^\star;B\setminus \overline\Omega)&\leq 
\mathcal A(\psi_h^\star;B)-|G_h^\star|-\mathcal A(\psi_h^\star;B\setminus \overline\Omega)\\
&\leq 
\mathcal A(\psi;B)-|G|-\mathcal A(\psi;B\setminus \overline\Omega),
\end{align*}
with strict inequality unless $\psi_{h+1}^\star=\psi_h^\star$ for all $h$ a.e. in $\Om$.

\subsubsection*{Conclusion} If $N=1$ it is sufficient to apply
the base case. If instead $N>1$ after a finite number $h^\star\le 2N-1$ of iterations we obtain a collection of mutually disjoint and closed convex sets with nonempty interiors $F_1:=F_{1,h^\star},\dots,F_{\tilde n}:=F_{\tilde n,h^\star}$ with $1\le \tilde n\le N$ such that
\[
G\subset \bigcup_{i=1}^{\tilde n}F_i=:G^\star,
\]
and
\begin{equation*}\mathcal{A}(\psi^\star;B)-|G^\star|-\mathcal{A}(\psi^\star;B\setminus \overline\Om)\le \mathcal A(\psi; B)-|G|-\mathcal{A}(\psi;B\setminus\overline{\Om}),
\end{equation*}
with
\begin{equation*}
\psi^\star:=\psi_{h^\star}^\star=\begin{cases}
0&\text{in}\ G^\star,\\
\psi&\text{otherwise},
\end{cases}
\end{equation*}
with strict inequality unless $\psi^\star=\psi$ a.e.
\end{proof}

\begin{proof}[Proof of Theorem \ref{teo:reduction_from_W_to_Wconv}]
We start by observing that \eqref{equivalence-W-Wconv} readily follows from \eqref{convexification}. Indeed, this implies
\begin{equation*}
\inf_{(\sigma,\psi)\in \admclassconv}
\mathcal{F}(\sigma,\psi)\le \inf_{(\sigma,\psi)
\in\largercomp} {\mathcal{F}}(\sigma,\psi)
.
\end{equation*}
Whereas from
\eqref{eq:W_subset_hatW}
it follows
\begin{equation*}
\inf_{(\sigma,\psi)\in\largercomp}
{\mathcal{F}}(\sigma,\psi)\le\inf_{(\sigma,\psi)\in\admclassconv }\mathcal{F}(\sigma,\psi).
\end{equation*}
Thus, we~only need to show \eqref{convexification}.
Take a pair $(\bar\sigma,\bar\psi)\in\largercomp$; we
suitably modify $(\bar\sigma,\bar\psi)$ into a new pair
$(\sigma,\psi)\in\admclassconv$ such that every connected component of $E(\sigma)$ is convex and
\begin{equation*}
\mathcal{F}(\sigma,\psi)\le{\mathcal F}(\bar\sigma,\bar\psi),
\end{equation*}
and this will conclude the proof. Once again we notice the that strict inequality holds unless $\psi=\bar\psi$ a.e.

Let $E(\bar\sigma_1),\dots,E(\bar\sigma_n)$ be the closed sets with mutually disjoint interiors corresponding to $\bar\sigma$ (as in \ref{ii'}
before \eqref{Sigma}) and let $G:=\bigcup_{i=1}^nE(\bar\sigma_i)$. Let
$F_1,\dots, F_N$ be the
(closure of the) connected components
of $G$, $N\leq n$, which are piecewise Lipschitz.
By~Lemma \ref{lem:reduction_of_energy_II} there exist $1\le \tilde n\le N$ and $\widetilde F_1,\dots,\widetilde F_{\tilde n}\subset\overline\Om$ mutually disjoint closed and convex sets with nonempty interior satisfying \eqref{A}, \eqref{B} and \eqref{C}.
Therefore, by construction, for every $i=1,\dots,n$,
$q_i$ and $p_{i+1}$ belong to $\widetilde F_j$ for a unique $j\in\{1,\dots,\tilde n\}$.
For every $j=1,\dots,\tilde n$ we denote by
\begin{equation*}
q_{j_1},p_{j_1+1},\dots, q_{j_{n_j}},p_{j_{n_j}+1},
\end{equation*}
the ones that belong to $\widetilde F_j$. Then we conclude by
taking $(\sigma,\psi)\in\admclassconv$ with $\sigma:=(\sigma_1,\dots,\sigma_n)$ and
\begin{equation*}
\sigma_{j_k}([0,1])=\begin{cases}
\overline{q_{j_k}p_{j_k+1}}&\text{for}\quad k=1,\dots,n_j-1,\\
\partial\widetilde F_j\setminus\Bigl(\bigcup_{h=1}^{n_j}\partial^0_{j_h}\Om\Bigr)\cup
\Bigl(\bigcup_{h=1}^{n_j-1}\overline{q_{j_h}p_{j_h+1}}\Bigr)&\text{for}\quad k=n_j,
\end{cases}
\end{equation*}
for every $j=1,\dots,\tilde n$ and $\psi:=\psi^\star$.
\end{proof}

\section{Existence of minimizers of \texorpdfstring{$\mathcal F$ in $\admclassconv$}{FW}}

\label{sec:existence_of_minimizers_of_F_in_Wconv}
The main result of this section reads as follows.

\begin{theorem}[Existence of a minimizer of $\mathcal F$ in $\admclassconv$]\label{thm:existence}
Let $\mathcal F$ and $\admclassconv$ be as in \eqref{def:relaxed_functional}
and \eqref{def:admissible_class} respectively. Then there
is $(\sigma, \psi)\in\admclassconv$ such that
\begin{equation}\label{eq:min_F_storto}
\mathcal F(\sigma,\psi)=
\min_{(s,\zeta)\in\admclassconv}\mathcal F(s,\zeta).
\end{equation}
Moreover, every minimizer $(\sigma, \psi)$ of $\mathcal F$ in $\admclassconv$ is such that every connected component of $E(\sigma)$ is convex.
\end{theorem}
As a direct consequence of Theorem \ref{teo:reduction_from_W_to_Wconv} and Theorem \eqref{thm:existence}, we have:

\begin{cor}\label{cor:mininmathcalW}
Let $(\sigmamin,\psimin)\in\admclassconv$ be a minimizer as in Theorem \ref{thm:existence}. Then $(\sigmamin,\psimin)$ is also a minimizer of ${\mathcal F}$ in the class $\largercomp$. Moreover, every minimizer $(\sigma, \psi)$ of $\mathcal F$ in $\mathcal W$ is such that every connected component of $E(\sigma)$ is convex.
\end{cor}

We prove Theorem \ref{thm:existence} using the direct method. To this aim
we need to introduce a notion of convergence in $\admclassconv$.
\begin{definition}[Convergence in $\admclassconv$]
\label{def:conv}
We say that the sequence $((\sigma)_k,\psi_k)_k\subset \admclassconv
$, with $(\sigma)_k=((\sigma_1)_k,\dots,(\sigma_n)_k)$, converges to $(\sigma,\psi)\in \admclassconv$ if:
\begin{enumerate}[{\rm (a)},wide]
\item\label{a}
$((\sigma_{i})_k)_\sharp\jump{[0,1]}$ converges to
$(\sigma_i)_\sharp \jump{[0,1]}$ in the sense of currents in $\mathcal D_1(\R^2)$,
for~all $i=1,\dots,n$;
\item \label{b} $(\psi_k)_k$ converges
to $\psi$ weakly* in $BV(\Omega)$, \ie $\psi_k \to \psi$ in $L^1(\Omega)$ and
$D \psi_k \rightharpoonup D\psi$ weakly* in $\Omega$
as measures as
$k \to +\infty$.
\end{enumerate}
\end{definition}
In Definition \ref{def:conv} $(\sigma_i)_\sharp\jump{[0,1]}$ denotes the push-forward by $\sigma_i$ of the $1$-current given by integration on the segment $[0,1]$, oriented in a standard way (see \cite{Krantz-Parks} for details).

In the next lemma we show a compactness property of $\admclassconv$. In~particular, given $(\sigma)_k\subset\convexcurves$ with equibounded energies, for all $i=1,\dots,n$, up to subsequences, a (not-relabeled) reparametrization of $(\sigma_i)_k$ converges uniformly to some $\widehat\sigma_i$, and there is a parametrization $\sigma_i$ of the support of $(\widehat \sigma_i)_\sharp\jump{[0,1]}$ such that $\sigma=(\sigma_1,\dots,\sigma_n)\in\convexcurves$. This,
together with a uniform bound on the lengths of $(\sigma_i)_k$, implies the convergence of the push-forwards as currents. Notice that $(\sigma_i)_\sharp\jump{[0,1]}$ is invariant under reparametrization of $\sigma_i$.

\begin{lemma}[Compactness of $\admclassconv$]\label{lem:compactness}
Let
$\big((\sigma)_{k},\psi_k\big)_k
\subset\admclassconv$ be a sequence with\\ $\sup_k\mathcal F((\sigma)_k,\psi_k)<+\infty$. Then
$\big((\sigma)_{k},\psi_k\big)_k$ admits
a subsequence converging to an element of $\admclassconv$.
\end{lemma}
\begin{proof}We divide the proof in two steps.\vspace*{-.5\baselineskip}

\subsubsection*{Step 1: Compactness of $(\sigma)_k$}
For simplicity we use the notation $\sigma_{ik} = (\sigma_i)_k$ for every $k\in
\mathbb N$ and $i\in\{1,\dots,n\}$.
By condition (P) in Remark \ref{rem:condition_P}, for every $k \in \mathbb N$ and $i\in\{1,\dots,n\}$
there exists a non-decreasing
function
\[
\theta_{ik}\colon[0,1]\to\R,\quad \theta_{ik}(1)-\theta_{ik}(0)\le2\pi,
\]
such that, for a reparametrization $\widehat \sigma_{ik}$ of $\sigma_{ik}$,
\begin{equation*}
\widehat\sigma_{ik}(t)=q_i+ \len(\sigma_{ik})\int_0^t\gamma_{ik}(s)ds,\quad \gamma_{ik}(t):=(\cos\theta_{ik}(t),\,\sin\theta_{ik}(t))
\quad\forall t\in[0,1],
\end{equation*}
and with $\widehat\sigma_{ik}(1)=p_{i+1}$. We~observe that
\begin{equation}\label{ti-chiamo}
\len(\sigma_{ik})=\int_0^1|\sigma_{ik}'(t)|dt\le\mathcal{H}^1(\partial\Om),
\end{equation}
since the orthogonal projection
\[
\Pi_{ki}\colon\partial\Om\setminus\partial^0_i\Om\to E(\sigma_{ik})
\]
is a contraction and $\mathcal{H}^1(\partial\Om\setminus\nobreak\partial^0_i\Om)\le \mathcal{H}^1(\partial\Om)$.
Hence, up to a (not relabeled) subsequence,
$\len(\sigma_{ik})\to m_i\in \R^+$ as $k \to +\infty$. The number $m_i$ is positive since, for all $k$ and $i$, we~have
$\len(\sigma_{ik})\geq |q_i-p_{i+1}|>0$.
Moreover,
\begin{equation*}
\int_0^1|\theta'_{ik}(t)|dt=\int_0^1\theta_{ik}'(t)dt\le2\pi;
\end{equation*}
hence, up to a not relabeled subsequence, $\theta_{ki}\weakstar\theta_i$ in $BV(0,1)$ and $\theta_i$ is non-decreasing with $\theta_i(1)-\theta_i(0)\le2\pi$. Furthermore $\gamma_{ik}\stackrel{*}{\rightharpoonup}\gamma_i$ in $BV((0,1);\R^2)$ with $\gamma_i(t)=(\cos(\theta_i(t)),\,\sin(\theta_i(t)))$. Thus, arguing as in \eqref{ti-chiamo-per-nome} and using \eqref{ti-chiamo}, we~get $\widehat\sigma_{ik}\to\widehat\sigma_i$ in $W^{1,1}([0,1];\R^2)$, where
\begin{equation}\label{sigmai}
\widehat\sigma_i(t):=q_i+m_i\int_0^t\gamma_i(s)ds=q_i+\len(\sigma_i)\int_0^t\gamma_i(s)ds.
\end{equation}
Thus $\lim_{k \to +\infty}\widehat\sigma_{ik}=\widehat\sigma_i$ uniformly, hence we also conclude that $\widehat\sigma_i$ takes values in $\overline\Om$.
Since by \ref{H3}
\[
d_H(E(\sigma_{ik}),E(\sigma_{ih}))=d_H(\partial E(\sigma_{ik}),\partial E(\sigma_{ih}))\leq \|\sigma_{ik}-\sigma_{ih}\|_{L^\infty}
\]
for all $h,k>0$,
the uniform convergence of $(\widehat\sigma_{ik})$ implies that $(E(\sigma_{ik}))_k$ is a Cauchy sequence with respect to the Hausdorff distance. Hence, by \ref{H2} there is $K_i\in\mathcal K$ such that $ d_H(E(\sigma_{ik}),K_i)\to0$, and $K_i$ is also convex by \ref{H5}.

We now show that $\widehat\sigma_i$ is injective, unless a pathological case that might happen only if $\partial_i^0\Om$
is a straight segment.\footnote{This case corresponds to $E(\sigma_{ik})$ a possibly curvilinear triangle with vertices $p_i$, $q_{i+1}$ and a third point $r_k\in \Om$ converging to a point $r\in \partial \Om$ which is on the same line as $p_i$, $q_{i+1}$, but outside the segment $\overline{p_iq_{i+1}}$.} Notice that, if $\partial_i^0\Om$ is not straight, $K_i$ must have nonempty interior, since it contains the region enclosed between $\overline {q_ip_{i+1}}$ and $\partial_i^0\Om$.

First observe that $\widehat\sigma_i([0,1])\subseteq \partial K_i$. Assume by contradiction that $\widehat\sigma_i(t_1)=\widehat\sigma_i(t_2)$ for some $t_1,t_2\in [0,1]$, $t_1<t_2$. Since $K_i$ is convex, the curve $\widehat\sigma_i\res [t_1,t_2]$ is closed and its image is contained in $\partial K_i$. If $\widehat\sigma_i\res [t_1,t_2]$ is constant and equals to $\widehat\sigma_i(t_1)$ we get a contradiction with \eqref{sigmai} and the fact that $|\gamma_i|=1$ a.e. in $[t_1,t_2]$. Hence there is a point $t_3\in(t_1,t_2)$ such that $\widehat\sigma_i(t_3)\neq \widehat\sigma_i(t_1)$. Let $\ell^k_{13}$ and $\ell^k_{23}$
denote the half-lines in $\R^2$ with endpoint ${\widehat\sigma_{ik}(t_3)}$ and passing through $\widehat\sigma_{ik}(t_1)$ and $\widehat\sigma_{ik}(t_2)$, respectively.
Since $E(\sigma_{ik})$ is convex, we~infer that $\widehat\sigma_{ik}([0,t_1])\cup \widehat\sigma_{ik}([t_2,1])$ is contained in the closed angular sector of $\R^2$ enclosed between $\ell^k_{13}$ and $\ell^k_{23}$.
Since $(\widehat\sigma_{ik})$ converges uniformly to $\widehat\sigma_i$, we~have
$\widehat\sigma_{ik}(t_j)\to \widehat\sigma_{i}(t_j)$ for $j=1,2,3$, and $\widehat\sigma_i(t_3)\neq \widehat\sigma_i(t_1)=\widehat\sigma_i(t_2)$, so we easily conclude that $\widehat\sigma_{ik}([0,t_1])\cup \widehat\sigma_{ik}([t_2,1])$ must be contained in the line passing through $ {\widehat\sigma_{i}(t_1)=\widehat\sigma_{i}(t_2)}$ and ${\widehat\sigma_{i}(t_3)}$. As a consequence also $K_i$, being convex, is a segment contained in such a line, and has empty interior. Hence this leads to a contradiction if $\partial_i^0\Om$ is not a straight segment. In this case we set $\sigma_i:=\widehat\sigma_i$.

If instead $\partial_i^0\Om$
is a straight segment, it might happen that the image of $\widehat\sigma_i$ is contained in a line, which must be the one passing through $q_i$ and $p_{i+1}$. Since uniform convergence of $(\widehat\sigma_{ik})$ and the fact that $\ell(\sigma_{ik})\to\ell(\widehat\sigma_i)$ imply that $(\widehat\sigma_{ik})_\sharp\jump{[0,1]}=(\sigma_{ik})_\sharp\jump{[0,1]}\to (\widehat\sigma_{i})_\sharp\jump{[0,1]}$ as currents, and since $\partial(\sigma_{ik})_\sharp\jump{[0,1]}=\delta_{p_{i+1}}-\delta_{q_i}$ for all $k$, also $\partial(\widehat\sigma_{i})_\sharp\jump{[0,1]}=\delta_{p_{i+1}}-\delta_{q_i}$. We~conclude that $(\widehat\sigma_{i})_\sharp\jump{[0,1]}$ is the integration over the segment $\overline{q_i p_{i+1}}$, and hence there is a Lipschitz injective curve $\sigma_i$ which parametrizes $\overline{q_i p_{i+1}}$ such that
\[
(\sigma_{i})_\sharp\jump{[0,1]}=(\widehat\sigma_{i})_\sharp\jump{[0,1]},\quad\text{and}\quad
(\sigma_{ik})_\sharp\jump{[0,1]}\to (\sigma_{i})_\sharp\jump{[0,1]}.
\]

We next show that $E(\sigma_i)$
is convex for any $i \in \{1,\dots, n\}$. If $\sigma_i$ parametrizes the segment $\overline{q_i p_{i+1}}$ then $E(\sigma_i)$ is that segment, and there is nothing to prove. Assume then that $\sigma_i([0,1])\neq \overline{q_i p_{i+1}}$. As shown above, the uniform limit $\sigma_i$ of $(\widehat\sigma_{ik})$ is injective. We~will show that $K_i=E(\sigma_i)$. Indeed, the uniform convergence of $(\widehat\sigma_{ik})$ yields
\[
\lim_{k\to+\infty}d_H(\partial E(\sigma_{ik}),\partial E(\sigma_{i}))=0.
\]
From \ref{H3} we get
\begin{align*}
d_H(\partial K_i, \partial E(\sigma_i))&\le d_H(\partial E(\sigma_{ik}),\partial K_i)+ d_H(\partial E(\sigma_{ik}),\partial E(\sigma_{i}))\\
&= d_H( E(\sigma_{ik}),K_i)+ d_H(\partial E(\sigma_{ik}),\partial E(\sigma_{i}))\to0\quad\text{as}\quad k\to+\infty.
\end{align*}
Thus $\partial K_i\!=\!\partial E(\sigma_{i})$, so $K_i\!=\! E(\sigma_{i})$ and the convexity is shown. This implies $\sigma\!\in\! \convexcurves$, and since $(\sigma_{ik})_\sharp\jump{[0,1]}\to (\sigma_{i})_\sharp\jump{[0,1]}$ as currents, the compactness of $(\sigma)_k$ is achieved.

\subsubsection*{Step 2: Compactness of $(\psi_k)$}
Setting $F_k=\bigcup_{i=1}^nE(\sigma_{ik})$ we have
\begin{equation*}
|D\psi_k|(\Om)\le \mathcal{A}(\psi_k;\Om)\le \mathcal{F}((\sigma)_k,\psi_k)+|F_k|\le C<+\infty \qquad \forall k>0,
\end{equation*}
where we used that $|F_k|\le |\Om|$. Therefore,
up to a subsequence, $\psi_k\weakstar\psi$ in $BV(\Om)$ and almost everywhere in $\Om$ as $k \to +\infty$.
To conclude it remains to show that \hbox{$\psi=0$} in $E(\sigma)=\bigcup_iE(\sigma_{i})$. If for some $i\in \{1,\dots,n\}$ it happens that $\partial_i^0\Om$ is straight and~$\sigma_i$ is the straight segment $\overline{q_{i}p_{i+1}}$, then $E(\sigma_i)$ has empty interior, and so there is nothing to prove. Otherwise, for the other indices, by $\lim_{k \to +\infty}d_H(E(\sigma_{ik}),E(\sigma_i))=\nobreak0$,
property~\ref{H6} yields
\begin{equation*}
\text{if}\quad x\in \text{int}(E(\sigma_i)) \quad \text{then}\quad x\in E(\sigma_{ik})\quad\text{for $k$ sufficiently large,}
\end{equation*}
and hence, since $\lim_{k\to +\infty} \psi_k= \psi$ a.e. in $\Om$, we~infer $\psi=0$ a.e. in $ E(\sigma)$.
\end{proof}

\begin{remark}\label{remark_uniform}
The previous proof shows a slightly stronger result: under the assumption of Lemma \ref{lem:compactness},
for every $i=1,\dots,n$, we~can find $\sigma_i$ with $\sigma=
(\sigma_1,\dots,\sigma_n)\in \Sigma_\mathrm{conv}$, $\widehat \sigma_i\in\Lip([0,1];\overline\Om)$, and
reparametrizations $\widehat\sigma_{ik}$ of $\sigma_{ik}$ such that
\[
(\widehat \sigma_i)_\sharp\jump{[0,1]}=(\sigma_i)_\sharp\jump{[0,1]},
\]
\[
\widehat\sigma_{ik}\to\widehat \sigma_i\quad\text{uniformly on }\ [0,1].
\]
Moreover, $( \sigma_{ik})_\sharp\jump{[0,1]}$ converges to $( \sigma_i)_\sharp\jump{[0,1]}$ in the sense of currents in $\mathcal D_1(\R^2)$.
Finally $E(\sigma_{ik})=E(\widehat\sigma_{ik})$ converges to $ E(\widehat \sigma_i)= E( \sigma_i)$ in $(\mathcal K,d_H)$, and $\widehat \sigma_i=\sigma_i$ unless $\partial_i^0\Om$ is a straight segment. In the latter case it might happen that $\widehat \sigma_i$ is not injective, but this
happens only if $\widehat \sigma_i([0,1])$ is a segment, $\sigma_i$ is a parametrization of $\overline{q_ip_{i+1}}$, and $E(\sigma_i)=\overline{q_ip_{i+1}}$.
\end{remark}
\begin{remark}\label{rem:conv}
We have also shown that if $(\widehat\sigma_{ik})$ converges uniformly to $\sigma_i\in\convexcurves$ for some $i=1,\dots,n$ then
\[
\lim_{k \to +\infty}
d_H(E(\sigma_{ik}),E(\sigma_i))= 0.
\]
\end{remark}

\begin{lemma}[Lower semicontinuity of $\mathcal F$ in $\admclassconv$]
\label{lem:lower_semicontinuity_of_F}
Let $\big((\sigma)_k,\psi_k\big)_k\subset\admclassconv$ be a sequence converging
to $(\sigma,\psi)\in \admclassconv$. Then
\begin{equation*}
\mathcal{F}(\sigma,\psi)\le\liminf_{k\to+\infty}\mathcal{F}((\sigma)_k,\psi_k).
\end{equation*}
\end{lemma}
\begin{proof}
By a standard argument \cite{Giusti:84}, the functional
\begin{equation*}
\psi\in BV(\Om)\mto \mathcal{A}(\psi;\Om)+\int_{\partial\Om}|\psi-\varphi|d\mathcal H^1
\end{equation*}
is $L^1(\Om)$-lower semicontinuous. We~now show that the map $\sigma \in
\convexcurves
\mto|E(\sigma)|$
is continuous.
Let $(\sigma)_k \subset \convexcurves$,
$\sigma \in\convexcurves$, and suppose that
$((\sigma_{i})_k)_\sharp\jump{[0,1]}$ converges to $(\sigma_i)_\sharp\jump{[0,1]}$ in $\mathcal D_1(\R^2)$
for all $i=1,\dots,n$ as $k \to +\infty$.
Set $F_k:=\bigcup_{i=1}^n E((\sigma_{i})_k)$ and recall that
$E(\sigma)=\bigcup_{i=1}^n E(\sigma_{i})$. { Thanks to Remark \ref{remark_uniform}, we~can always assume that there are reparametrizations $\widehat\sigma_{ik}$ of $\sigma_{ik}$ such that $\widehat\sigma_{ik}$ converges uniformly to $\widehat \sigma_i$ with $(\widehat \sigma_i)_\sharp\jump{[0,1]}=(\sigma_i)_\sharp\jump{[0,1]}$. Let us suppose first that $\widehat \sigma_i$ is injective for all $i=1,\dots,n$, and so $\widehat\sigma_i=\sigma_i$.}
By Remark \ref{rem:conv}
$\lim_{h \to +\infty} d_H(E((\sigma_{i})_k),E(\sigma_i)) = 0$
for all $i=1,\dots,n$ and therefore $d_H(F_k,E(\sigma)) =: \eps_k\to 0^+$.

By invoking \ref{H7} we have $E(\sigma)\subset (F_k)^+_{\eps_k}$.
Moreover, since $d_H((F_k)_{\eps_k}^+,E(\sigma))\leq 2\eps_k$, we~get $(F_k)^+_{\eps_k}\subseteq (E(\sigma))_{2\eps_k}^+$, and so
\begin{equation*}
\vert E(\sigma)\vert\le \vert (F_k)^+_{\eps_k}\vert
\le \vert (E(\sigma))_{2\eps_k}^+\vert.
\end{equation*}
This implies
\[
\limsup_{k\to+\infty} |F_k|\leq \limsup_{k\to+\infty} |(F_k)^+_{\eps_k}|\le\mathcal |E(\sigma)|.
\]
The converse inequality is a consequence of Fatou's lemma and \ref{H6}, indeed
\[
|E(\sigma)|\leq \int_\Om\liminf_{k\to+\infty}{\chi_{F_k}(x)}~dx\leq\liminf_{k\to+\infty}\int_\Om{\chi_{F_k}(x)}~dx =\liminf_{k\to+\infty}|F_k|.
\]
If instead $\widehat\sigma_i$ is not injective for some $i$, we~have $\widehat \sigma_i\in\Lip([0,1];\overline\Om)$ with $(\widehat \sigma_i)_\sharp\jump{[0,1]}=(\sigma_i)_\sharp\jump{[0,1]}$, and we are in the case that $E(\widehat\sigma_i)$ has empty interior (see Remark \ref{remark_uniform}). Thus $ E(\sigma_{ik})=E(\widehat\sigma_{ik})$ converges to a segment $K_i\supsetneq E(\sigma_i)$ in the Hausdorff distance.
Since $|K_i|=0$, the thesis of the lemma follows along the same argument above replacing the symbol $E(\sigma_i)$ by $K_i$.
\end{proof}

\begin{proof}[Proof of Theorem \ref{thm:existence}]
By Lemma \ref{lem:compactness} and Lemma \ref{lem:lower_semicontinuity_of_F} we can apply the direct method and conclude that there exists $(\sigma, \psi)\in\admclassconv$ such that
\eqref{eq:min_F_storto} holds. Moreover, since $\admclassconv\subset\largercomp$ by Theorem \eqref{teo:reduction_from_W_to_Wconv} we can choose $(\sigma,\psi)$ such that every connected component of $E(\sigma)$ is convex.
\end{proof}

\section{Regularity of minimizers}
\label{sec:regularity_of_minimizers}

In this section we investigate regularity properties of
minimizers of $\mathcal F$. We~recall that our boundary datum
$\varphi$
satisfies the conditions
in \eqref{boundary-value}, and $\widehat\varphi\in W^{1,1}(B)$ denotes a
fixed extension of $\varphi$ in the open ball $B\supset\overline\Omega$.
The main result here reads as follows.
\begin{theorem}[Structure of minimizers]
\label{teo:structure_of_minimizers}
Every minimizer $(\sigma,\psi)\in\admclassconv$ of $\mathcal F$ in~$\largercomp$, namely
\begin{equation*}
{\mathcal{F}}(\sigma,\psi)=\min_{(s,\zeta)\in\largercomp}
{\mathcal{F}}(s,\zeta),
\end{equation*}
satisfies the following properties:
\begin{enumerate}[{\rm(1)},wide]
\item\label{1.}
Each connected component of $E(\sigma)$ is convex;
\item\label{2.} $\psi$ is positive and real analytic
in $\Om\setminus E(\sigma)$;
\item\label{3.} If $\partial^D_i\Om$ is not
a segment
for some $i=1,\dots,n$, then $\partial E(\sigma)\cap
\partial_i^D\Om=\emptyset$, $\psi$ is continuous up to $\partial_i^D\Om$, and
$\psi=\varphi$ on $\partial_i^D\Om$;
\item\label{4.} If $\partial^D_i\Om$ is a segment for some $i=1,\dots,n$, then either $\partial E(\sigma)\cap
\partial_i^D\Om=\emptyset$ or $\partial E(\sigma)\cap
\partial_i^D\Om=\partial_i^D\Om$. In the first case $\psi$
is continuous up to $\partial_i^D\Om$ and
$\psi=\varphi$ on~$\partial_i^D\Om$.
\end{enumerate}
Moreover, there is a minimizer $(\sigma,\psi)\in\admclassconv$ such that
\begin{enumerate}[{\rm(5)},wide]
\item\label{5.}
$\Om\cap \partial E(\sigma)$ consists of a finite number of disjoint analytic curves, and $\psi$ is continuous and null on $\partial E(\sigma)\setminus \partial^D\Omega$.
\end{enumerate}
\end{theorem}

\begin{remark}
If $\partial_i^D\Omega$ is a straight segment for some $i=1,\dots,n$,
nothing ensures that $\partial E(\sigma)\cap \partial_i^D\Omega=\emptyset$. However, if this intersection is nonempty, then necessarily
$\partial_i^D\Omega\subset\partial E(\sigma)$. The
prototypical example is given by the classical catenoid, as explained in the introduction (see Figure \ref{figura1}) where, if the basis of the rectangle $\Omega=R_{\ell}$ is large enough, a solution $\psi$ is identically zero, and $\partial^D\Omega\subset\partial E(\sigma)$.
This also explains why in point \ref{5.} of Theorem \ref{teo:structure_of_minimizers} we write $\partial E(\sigma)\setminus \partial^D\Omega$.
\end{remark}

A consequence of Theorem \ref{teo:structure_of_minimizers} is that a regular solution $\psi$ belongs to $W^{1,1}(\Om)$ and, if $\Om$ is strictly convex, it also attains the boundary values. In~particular, Theorem \ref{teo:structure_of_minimizers} implies Theorem \ref{teo_main_intro}.

For the reader convenience we divide the proof in a number of steps.
\begin{lemma}\label{lem:analicity}
Every minimizer $(\sigma,\psi)\in\admclassconv$ of $\mathcal F$ in
$\largercomp$ satisfies \ref{1.}, \ref{2.} and
$\psi=\varphi$ on $\partial^D\Om\setminus\partial E(\sigma)$.
\end{lemma}
\begin{proof} Item \ref{1.} follows by Theorem \ref{teo:reduction_from_W_to_Wconv}.
By \cite[Th.\,14.13]{Giusti:84}
we also have that~$\psi$ is real analytic in $\Om\setminus E(\sigma)$. Together with the strong maximum principle \cite[Th.\,C.4]{Giusti:84}, this implies that,
in $\Om\setminus E(\sigma)$,
either $\psi>0$ or $\psi\equiv0$.
On the other hand,
since $\Omega$ is convex we can apply \cite[Th.\,15.9]{Giusti:84} and get that $\psi$ is continuous up to $\partial^D\Om\setminus\partial E(\sigma)$; in particular
\begin{equation}\label{eq:boundary_value}
\psi=\varphi>0\quad\text{on}\ \partial^D\Omega\setminus\partial E(\sigma),
\end{equation}
which in turn implies $\psi>0$ in $\Omega\setminus E(\sigma)$ .
\end{proof}

\begin{lemma}\label{lem:plateau_solution}
Let $\Gamma \subset \R^3$
be a rectifiable, simple, closed and non-planar
curve satisfying the following properties:
\begin{enumerate}[label=$(\arabic*)$,wide]
\item \label{(1)}$\Gamma\subset\partial(F\times\R)$ for some closed bounded convex set $F\subset\R^2$ with nonempty interior;
\item \label{(2)}$\Gamma$ is symmetric with respect to the horizontal plane $\R^2\times\{0\}$;
\item \label{(3)} There are a nonempty
relatively open arc $\arc{pq}\subset\partial F$ with endpoints $p$ and $q$, and
$f\in C^0( \arc{pq}\cup\{p,q\};[0,+\infty))$ such that $f$ is positive in $\arc{pq}$ and
\begin{equation}\label{ass:sym-Gamma}
\Gamma\cap\{x_3\ge0\}=\mathcal{G}_f\cup (\{p\}\times[0,f(p)])\cup (\{q\}\times[0,f(q)]).
\end{equation}
\end{enumerate}
Let $S$ be a solution to the
classical Plateau problem for $\Gamma$, \ie a disk-type surface
minimizing area among all disk-type surfaces spanning $\Gamma$.
Then:
\begin{enumerate}[label=$(\arabic*')$,wide]
\item \label{1'}$\beta_{p,q}:=S\cap (\R^2\times\{0\})\subset F$
is a simple analytic curve joining $p$ and $q$ with $\beta_{p,q}\cap\partial F=\{p,q\}$;
\item\label{2'} $S$ is symmetric with respect to $\R^2\times\{0\}$;
\item \label{3'}The surface $S^+:=S\cap\{x_3\ge0\}$ is the graph of a
function { $\widetilde \psi\in W^{1,1}(U_{p,q})\cap C^0(\overline{U}_{p,q}\setminus{\{p,q\}})$}, where $U_{p,q}\subset\mathrm{int} (F)$ is the open region enclosed between $\arc{pq}$ and $\beta_{p,q}$. Moreover, $\widetilde\psi$ is analytic in $U_{p,q}$, { and if $f(p)=0$ (\resp $f(q)=0$) then $\psi$ is also continuous at $p$ (\resp at $q$);}
\item \label{4'}The curve $\beta_{p,q}$ is contained in the closed
convex hull of $\Gamma$, and $F\setminus U_{p,q}$ is convex.
\end{enumerate}
\end{lemma}

\begin{remark}
If the function $f$ in \ref{(3)} is such that $f(p)=f(q)=0$ then \eqref{ass:sym-Gamma} becomes $\Gamma\cap\{x_3\ge0\}=\mathcal{G}_f$.
For later convenience we prove Lemma \ref{lem:plateau_solution} under the more general assumption \ref{(3)}.
\end{remark}

\begin{proof}[Proof of Lemma \ref{lem:plateau_solution}]
Even though several arguments are
standard, we~give the proof for completeness.

\subsubsection*{Step 1: $\beta_{p,q}$ is a simple analytic curve joining $p$ and $q$}
Let $B_1\subset\R^2$ be the open unit disk centered at the origin. Let $\Phi=(\Phi_1,\Phi_2,\Phi_3)\colon\overline{B}_1\to S\subset\R^3$
be a parametrization of $S$ with $\Phi(\partial B_1)=\Gamma$,
that is harmonic, conformal, and therefore analytic in $B_1$, continuous up to $\partial B_1$.
Further, by \ref{(1)},
$\Phi$ is an embedding (see \cite{MY} and also \cite[p.\,343]{DHS}).

By assumption \eqref{ass:sym-Gamma} we have $\{w\in\partial B_1\colon\Phi_3(w)=0\}=\{\Phi^{-1}(p,0),\Phi^{-1}(q,0)\}$, so~that $\Phi_3$ changes sign only twice on $\partial B_1$. By applying Rado's lemma (see, e.g., \cite[Lem.\,2, p.\,295]{DHS}) to the harmonic function $\Phi_3$ we deduce that $\nabla\Phi_3\ne0$ in $B_1$ and in particular $\{w\in B_1\colon\Phi_3(w)>0\}$ and $\{w\in B_1\colon\Phi_3(w)<0\}$
are connected, and $\{w\in B_1\colon\Phi_3(w)=0\}$ is a simple smooth curve in $B_1$ joining $\Phi^{-1}(p,0)$ and $\Phi^{-1}(q,0)$.
By the injectivity of $\Phi$ we have that $S\cap(\R^2\times\{0\})=\Phi(\{w\in B_1\colon\Phi_3(w)=0\})$ is a simple analytic curve joining $p$ and $q$.

\subsubsection*{Step 2: $S$ is symmetric with respect to the horizontal plane $\R^2\times\{0\}$}
By step 1 the sets $\{w\in \overline B_1\colon\Phi_3(w)\ge0\}$ and $\{w\in \overline B_1\colon\Phi_3(w)\le0\}$ are simply connected and the two surfaces
\[
S^+:=\Phi(\{w\in \overline B_1\colon\Phi_3(w)\ge0\}),\quad S^-:=\Phi(\{w\in \overline B_1\colon\Phi_3(w)\le0\})
\]
have the topology of the disk. We~assume without loss of generality that $\mathcal{H}^2(S^+)\le \mathcal{H}^2(S^-)$. Let
\[
\Sym(S^+):=\{(x',x_3)\colon (x',-x_3)\in S^+\},\quad \widetilde S:= S^+\cup\Sym(S^+).
\]
Then $\widetilde S$ is a symmetric surface of disk-type with $\partial\widetilde S =\Gamma$ and
\begin{equation*}
\mathcal{H}^2(\widetilde S)=2\mathcal{H}^2(S^+)\le\mathcal{H}^2(S^+)+\mathcal{H}^2(S^-)=\mathcal{H}^2(S).
\end{equation*}
In~particular, $\widetilde S$ is a symmetric solution to the
Plateau problem for $\Gamma$. Further $S=\widetilde S$ on a relatively
open subset of $S$; hence, since they are real
analytic surfaces, they must coincide, $S=\widetilde S$.

\subsubsection*{Step 3: $S^+$
is the graph of a function $\widetilde \psi\in W^{1,1}(U_{p,q})\cap C^0(\overline U_{p,q}\setminus{\{p,q\}})$}
To show this it is enough to
check the validity of the following

\begin{claim*}
Every vertical plane $\Pi$ is tangent to $\mathrm{int}(S)$ at most at one point.
\end{claim*}

We prove the claim arguing by contradiction as in
\cite[p.\,97]{vortex}, that is we assume there is a vertical plane $\Pi$ tangent to $\mathrm{int}(S)$ at $x'$ and $x''$ with $x'\ne x''$. We~define the linear map $d_\nu(x):=(x-x')\cdot\nu $ with $\nu$ a unit normal to $\Pi$, so that clearly $\Pi=\{x\in\R^3\colon d_\nu(x)=0\}$.
Since $F$ is convex, $\Pi\cap (\partial F\times\{0\})$ contains at most two points. By properties \ref{(1)}--\ref{(3)} each of these points is either the projection on the horizontal plane of one or two points of $\Pi\cap\Gamma$, or the projection on the horizontal plane of one of the vertical segments $\{p\}\times[0,f(p)]$ and $\{q\}\times[0,f(q)]$. Hence $\Pi\cap \Gamma$ contains either: (a) at most two points and a segment, (b) two segments, (c) four points.
Without loss of generality we restrict our analysis to the last case (the others are simpler to treat), namely we assume that there are four (clockwise ordered) points $w_1,\dots,w_4\in\partial B_1$ such that
$\Pi\cap\Gamma=\{\Phi(w_1),\dots,\Phi(w_4)\}$, that is $d_\nu\circ\Phi(w_i)=0$ for $i=1,\dots,4$. We~may also assume
$d_\nu\circ\Phi>0$ on $\arc{w_1w_2}\cup\arc{w_3w_4}$ and
$d_\nu\circ\Phi<0$ on $\arc{w_2w_3}\cup\arc{w_4w_1}$. Here $\arc{w_iw_j}$ denotes the relatively open arc in $\partial B_1$ joining $w_i$ and $w_j$ for $i,j\in\{1,\dots,4\}$.
Notice that the function $d_\nu\circ\Phi\colon \overline B_1\to \R$ is harmonic in $B_1$, continuous up to $\partial B_1$ and vanishes at $w_1,\dots,w_4$; hence,
by classical arguments \cite[\S 437]{Nitsche:89} we see that the set $\{w\in B_1\colon d_\nu\circ\Phi=0\}$, in a neighbourhood of $w':=\Phi^{-1}(x')$ (respectively $w'':=\Phi^{-1}(x'')$), is the union of a number $m\geq 2$ of analytic curves crossing at $w'$ (respectively $w''$). Thus near $w'$ and $w''$ the set $\{w\in B_1\colon d_\nu\circ\Phi(w)>0\}$ is the union of at least two disjoint open regions $A_{1,1}$, $A_{1,2}$ and $A_{2,1}$, $A_{2,2}$ respectively such that $\overline A_{1,1}\cap\overline A_{1,2}=\{w'\}$, $\overline A_{2,1}\cap\overline A_{2,2}=\{w''\}$. Moreover, each $A_{i,j}$ belongs either to the connected component of $\{w\in B_1\colon d_\nu\circ\Phi(w)>0\}$ containing $\arc{w_1w_2}$,
or to the one containing $\arc{w_3w_4}$. Up to relabeling the indices we have two possibilities.

\subsubsection*{Case 1: $A_{1,1}$ and $A_{1,2}$ belong to the same connected component containing $\arc{w_1w_2}$}
Then we can find two simple curves $\alpha_1$, $\alpha_2$ contained in $A_{1,1}$ and $A_{1,2}$ respectively, that connect $w'$ to a point in $\arc{w_1w_2}$ and such that the region enclosed by the curve $\alpha_1\cup\alpha_2$ intersects $\{w\in B_1\colon d_\nu\circ\Phi(w)<0\}$. Since $d_\nu\circ \Phi>0$ on $\alpha_1\cup\alpha_2$ by the maximum principle we have a contradiction.

\subsubsection*{Case 2: $A_{1,1}$ and $A_{2,1}$ belong to the connected component containing $\arc{w_1w_2}$ while $A_{1,2}$ and $A_{2,2}$ belong to the connected component containing $\arc{w_3w_4}$}
Then we can find four simple curves $\alpha_{i,j}$ (with $i,j=1,2$) contained respectively in $A_{i,j}$, such that $\alpha_{1,1}$ (respectively $\alpha_{2,1}$) connects $w'$ (respectively $w''$) to a point in $\arc{w_1w_2}$ and $\alpha_{1,2}$ (respectively $\alpha_{2,2}$) connects $w'$ (respectively $w''$) to $\arc{w_3w_4}$. Then the region enclosed by the curve $\bigcup_{i,j}\alpha_{i,j}$ intersects $\{w\in B_1\colon d_\nu\circ\Phi(w)<0\}$, while $d_\nu\circ\Phi>0$ on $\bigcup_{i,j}\alpha_{i,j}$, which again by the maximum principle gives a contradiction.

\smallskip
Thus the claim follows. { Now, by step 2, the claim readily implies that $\mathrm{int}(S^+)$ has no points with vertical tangent plane and hence $\mathrm{int}(S^+)$ is the graph of a function $\widetilde\psi$ defined on ${U_{p,q}}$. Since $\widetilde \psi$ must minimize (locally) the area functional, it is also real analytic in $U_{p,q}$. Moreover, the claim also implies that $\widetilde\psi$ must vanish on $\beta_{p,q}$ and that it must attain the boundary values on $\arc{pq}$. If $f$ vanishes on $p$ or $q$, then also the continuity of $\widetilde\psi$ at these points is achieved. }

\Subsubsection*{Step 4: The curve $\beta_{p,q}$ is contained in the closed
convex hull of $\Gamma$, and the set $F\setminus U_{p,q}$ is convex}
Let $\pi(\Gamma)\subset\partial F$ be the projection of $\Gamma$ onto the plane $\R^2\times\{0\}$. By \cite[Th.\,3, p.\,343]{DHS} the relative interior of $S$ is strictly contained in the convex hull of $\Gamma$, thus in particular
the curve $\beta_{p,q}$ (respectively $\beta_{p,q}\setminus\{p,q\}$) is contained (respectively strictly contained) in the same half-plane (with respect to the line $\overline{pq}$) that contains $\pi(\Gamma)$.

Now, assume by contradiction that $F\setminus U_{p,q}$ is not convex. Then there are $p',q'\in\beta_{p,q}$ with the following properties:
\begin{itemize}
\item The open region $U'$ enclosed by $\beta_{p,q}$ and the segment $\overline{p'q'}$ is nonempty and contained in $U_{p,q}$;
\item the points $p$ and $q$ and the set $U'$ lie on the same side with respect to the line containing $\overline{p'q'}$.
\end{itemize}
Let then $d_W\colon\R^3\to\R$ be an affine function that vanishes on the vertical plane containing $\overline{p'q'}$ and is positive in the half-space $W^+$ containing $p,q$ and $U'$. We~now observe that $\Gamma\cap W^+$ is
the union of two connected subcurves $\widehat\Gamma_1$ and $\widehat\Gamma_2$, containing $p$ and $q$ respectively. As a consequence $\Phi^{-1}(\widehat\Gamma_1)=\arc{w_1w_2}$ and $\Phi^{-1}(\widehat\Gamma_2)=\arc{w_3w_4}$ for some $w_1,w_2,w_3,w_4\in\partial B_1$ (clockwise oriented).

On the other hand since $d_W>0$ on $U'$ we can find $t'\in\partial U'\setminus \overline{p'q'}$ such that
\[
d_W\circ\Phi(\Phi^{-1}(t'))=d_W(t')>0
\]
with $\Phi^{-1}(t')\in B_1$.
Once again by the harmonicity of $d_W\circ\Phi\colon\overline{B}_1\to\R$ we deduce the existence of a curve $\alpha\subset\{w\in B_1\colon d_W\circ\Phi(w)>0\}$ joining $\Phi^{-1}(t')$ either to $\arc{w_1w_2}$ or $\arc{w_3w_4}$. Hence $\Phi(\alpha)\subset \Phi( B_1)$ is a curve joining $t'$ either to $\widehat\Gamma_1$ or $\widehat\Gamma_2$, say $\widehat\Gamma_1$.
This implies that the projection $\pi(\Phi(\alpha))$ of $\Phi(\alpha)$ onto the horizontal plane $\R^2\times\{0\}$ is a curve contained in $U_{p,q}$ that connects $t'$ to $\pi(\widehat\Gamma_1)$. So in particular, the curve $\pi(\Phi(\alpha))$ cannot be included in the half-space $W^+$. But this contradicts the fact that $\alpha\subset\{w\in B_1\colon d_W\circ\Phi(w)>0\}$ (this is because the values of $d_W$ at a point $x$ and~$\pi(x)$ are the same).
\end{proof}

We need also the following technical results on the distance function
$\di_F$ from a convex set $F$.
Recall the definition of $E^+_\eps$ given in \ref{H7} in the appendix,
for $\eps>0$ and $E\subset\R^2$.

\begin{lemma}\label{lem:prop_distanza_da_convesso}
Let $F\subset\R^2$ be bounded, closed and convex.
Then
\[
\Delta\di_F\in L^{\infty}_\mathrm{loc}(\R^2\setminus F)\cap L^1(B\setminus F)
\]
for every ball $B$ with $F\subset\subset B$.
\end{lemma}
\begin{proof}
By \cite[Th.\,3.6.7, p.\,75]{Can-Sin} it follows that $\di_F\in C^{1,1}_\mathrm{loc}(\R^2\setminus F)$, hence $\nabla^2\di_F\in L^{\infty}_\mathrm{loc}(\R^2\setminus F;\R^{2\times 2})$.
Therefore we only have to check that $\Delta\di_F\in L^1(B\setminus F)$.
Let $\eta>0$ be fixed sufficiently small.
Select $(f_k)_{k\in\mathbb{N}}\subset C^1_c(\R^2; \R^2)$ such that
$f_k\to \nabla\di_F$ in $W^{1,1}(B
\setminus F_{\eta/2}^+)$ as $k\to+\infty$.
By the divergence theorem we have
\begin{equation}\label{div-formula}
\int_{B\setminus F^+_\eta}\Div f_k\;dx=\int_{
\partial B\cup\partial (F^+_\eta)}f_k\cdot\nu_\eta \;d\mathcal{H}^1,
\end{equation}
with $\nu_\eta$ the outer unit normal to $\partial B\cup\partial (F^+_\eta)$.
By taking the limit as $k\to \infty$ we get
\begin{equation}\label{lim_k}
\lim_{k\to+\infty} \int_{B\setminus F^+_\eta}\Div f_k\;dx=\int_{B\setminus F^+_\eta}\Delta\di_F \;dx,
\end{equation}
and
\begin{equation}\label{lim_k_1}
\lim_{k\to+\infty} \int_{\partial B\cup\partial (F^+_\eta)}f_k\cdot\nu_\eta\; d\mathcal{H}^1= \int_{\partial B\cup\partial (F^+_\eta)}\nabla\di_F\cdot\nu_\eta\; d\mathcal{H}^1,
\end{equation}
where \eqref{lim_k_1} follows by using that $\partial(F^+_\eta)$ is of class $C^{1,1}$ and hence $f_k\res(\partial B\cup\partial (F^+_\eta))\to\nabla\di_F\res(\partial B\cup\partial (F^+_\eta))$
in $L^1(\partial B\cup\partial (F^+_\eta))$.
Since $\di_F$ is convex we have $\Delta\di_F\ge0$ a.e. in $\R^2\setminus F$, moreover
$|\nabla\di_F|=1$ in $\R^2\setminus F$; then gathering together \eqref{div-formula}, \eqref{lim_k}, \eqref{lim_k_1} we have
\begin{equation*}
\int_{B\setminus F^+_\eta}|\Delta\di_F|\; dx=
\int_{B\setminus F^+_\eta}\Delta\di_F \;dx=
\int_{\partial B\cup\partial (F^+_\eta)}\!\!\!\!\!\!\!\!\nabla\di_F\cdot\nu_\eta \;d\mathcal{H}^1\le \mathcal{H}^1(\partial B\cup\partial(F^+_\eta))\le C,
\end{equation*}
with $C>0$ independent of $\eta$.
By the arbitrariness of $\eta>0$, the thesis follows.
\end{proof}
\begin{cor}\label{cor:integration_by_parts}
Let $U\subset\R^2$ be a bounded open set with Lipschitz boundary. Let $F\subset\R^2$ be closed and convex such that $U\cap F=\emptyset$ and let $\psi\in W^{1,1}(U)\cap L^\infty(U)\cap C^0(U)$. Then the following formula holds:
\begin{equation*}
\begin{split}
- \int_{U}\psi\Delta\di_{F}\,dx
=
\int_{U} \nabla\psi\cdot\nabla\di_{F}\,dx- \int_{\partial U}\psi\,\gamma\,d\mathcal{H}^1,
\end{split}
\end{equation*}
where $\gamma$ denotes the normal trace of $\nabla\di_F$ on $\partial U$.
\end{cor}

\begin{proof}
We have $|\nabla \di_F|=1$ in $\R^2\setminus F$, moreover since $U\cap F=\emptyset$, by
Lemma \ref{lem:prop_distanza_da_convesso} we deduce also $\Delta\di_F\in L^1(U)$. Therefore the thesis readily follows by applying \hbox{\cite[Th.\,1.9]{A}}.
\end{proof}

\begin{remark}\label{rem:traccia}
The normal trace $\gamma$ of $\nabla\di_F$ on $\partial F$ equals $1$ $\mathcal H^1$-a.e. on $\partial F$. Indeed, from Corollary \ref{cor:integration_by_parts} we have that for all $\varphi\in C^1_c(\R^2;\R^2)$ it holds
\begin{align}
- \int_{\R^2\setminus F^+_\eta}\varphi\Delta\di_{F}\,dx
&=
\int_{\R^2\setminus F^+_\eta} \nabla\varphi\cdot\nabla\di_{F}\,dx- \int_{\partial(F_\eta^+)}\varphi\,\gamma\,d\mathcal{H}^1\nonumber\\
&=\int_{\R^2\setminus F^+_\eta} \nabla\varphi\cdot\nabla\di_{F}\,dx- \int_{\partial(F_\eta^+)}\varphi\,d\mathcal{H}^1,\nonumber
\end{align}
where we have used that $\partial (F_\eta^+)$ being a level set of $\di_F$, it results $\nabla\di_F=\nu_\eta$ on it.
Letting $\eta\to 0$ and using that $\Delta \di_F\in L^1(B\setminus F)$ for all balls $B$, we~infer
\begin{align}
-\int_{\R^2\setminus F}\varphi\Delta\di_{F}\,dx
&=
\int_{\R^2\setminus F} \nabla\varphi\cdot\nabla\di_{F}\,dx- \int_{\partial F}\varphi\,d\mathcal{H}^1.\nonumber
\end{align}
By the arbitrariness of $\varphi$ and again by Corollary \ref{cor:integration_by_parts}, the claim follows.
\end{remark}

\begin{lemma}\label{lem:limite_di_traccia}
Let $F\subset\overline\Om$ be closed and convex with nonempty interior,
and let $\delta>0$.
Let $\psi\in W^{1,1}((F^+_{\delta}\setminus F)\cap\Om)\cap L^\infty((F^+_{\delta}\setminus F)\cap\Om)\cap C^0((F^+_{\delta}\setminus F)\cap\Om)$. Then
\begin{equation}\label{eq:lim_tracce}
\lim_{\eps\to0^+} \int_{\Om\cap\partial(F^+_\eps)}\psi\,d\mathcal{H}^1=\int_{\Om\cap\partial F}\psi\,d\mathcal{H}^1.
\end{equation}
\end{lemma}

\begin{proof}
Let $\eps\in(0,\delta)$ and $T_\eps:=(F^+_\eps\setminus F)\cap \Omega$. Since $T_\eps\cap F=\emptyset$, by Corollary \ref{cor:integration_by_parts} we get
\begin{equation}\label{eq:lim_tracce*}
\begin{split}
-\int_{T_\eps}\psi\Delta\di_{F}\,dx
=
\int_{T_\eps} \nabla\psi\cdot\nabla\di_{F}\,dx- \int_{\partial T_\eps}\psi\,\gamma\,d\mathcal{H}^1,
\end{split}
\end{equation}
which, by Remark \ref{rem:traccia}, becomes
\begin{multline}\label{eq:lim-tracce1}
-\int_{T_\eps}\psi\Delta\di_{F}\,dx
= \int_{T_\eps} \nabla\psi\cdot\nabla\di_{F}\,dx\\
+ \int_{\Om\cap\partial F}\psi\,d\mathcal{H}^1
-\int_{\Om\cap\partial(F_\eps^+)}\psi\,d\mathcal{H}^1
-\int_{((F^+_\eps)\setminus F)\cap\partial \Om}\psi\,\gamma\,d\mathcal{H}^1.
\end{multline}
Now
\begin{equation}\label{eq:lim-tracce2}
\lim_{\eps\to0^+} \Big|\int_{T_\eps} \nabla\psi\cdot\nabla\di_{F}\,dx
\Big| \le\lim_{\eps\to0^+} \int_{T_\eps} |\nabla \psi|\,dx=0,
\end{equation}
and
\begin{equation}\label{eq:lim-tracce3}
\lim_{\eps\to0^+} \Big| \int_{(F^+_\eps\setminus F)\cap\partial\Om}\psi\,\gamma\,d\mathcal{H}^1\Big|\le
\lim_{\eps\to0^+} \int_{(F^+_\eps\setminus F)\cap\partial\Om}\psi\,d\mathcal{H}^1=0.
\end{equation}
Moreover, since $\Delta\di_{F}\in L^1(T_\eps)$ by Lemma \ref{lem:prop_distanza_da_convesso}, we~deduce also
\begin{equation}\label{eq:lim-tracce4}
\lim_{\eps\to0^+}\Big| \int_{T_\eps}-\psi\,\Delta\di_{F}\,dx\Big|\le \|\psi\|_{L^\infty}\lim_{\eps\to0^+} \int_{T_\eps}|\Delta \di_{F}|\,dx=0.
\end{equation}
Finally, gathering together \eqref{eq:lim-tracce1}--\eqref{eq:lim-tracce4}, we infer \eqref{eq:lim_tracce}.
\end{proof}
\begin{remark} \label{rem:limite_di_traccia} Let $F$, $\delta$ and $\psi$ be as in Lemma \ref{lem:limite_di_traccia}.
Let $\alpha$ be any connected component of $\Om\cap\partial F$, and for every $0<\eps<\delta$ let $\alpha_\eps$ be the corresponding component of $\Om\cap\partial (F^+_\eps)$; namely, if $\pi_F$ is the orthogonal projection onto the convex closed set $F$, setting
\[
\widehat\alpha_\eps:=\{x\in \partial (F^+_\eps):\pi_F(x)\in \alpha\},
\]
then one has $\alpha_\eps:=\widehat\alpha_\eps\cap \Om$.
Arguing as in Lemma \ref{lem:limite_di_traccia}, we~can show that
\begin{equation*}
\lim_{\eps\to0^+}
\int_{\alpha_\eps}\psi\,d\mathcal{H}^1=\int_{\alpha}\psi\,d\mathcal{H}^1.
\end{equation*}
\end{remark}

\begin{lemma}\label{lem:continuita-bordo-E}
Let $(\sigma,\psi)\in \admclassconv$ be a minimizer of $\mathcal F$ in $\mathcal W$ as in
Theorem \ref{teo:reduction_from_W_to_Wconv}.
Then there is a minimizer $(\widehat\sigma,\widehat\psi)\in \admclassconv$ of $\mathcal F$ in $\mathcal W$
with the following properties:
\begin{enumerate}[{\rm (1)},wide]
\item\label{lem1.}$(\partial E(\widehat\sigma))\cap\partial\Om=(\partial E(\sigma))\cap\partial\Om$;
\item\label{lem2.} $\widehat\psi$ is continuous and null
on $\Om\cap\partial E(\widehat\sigma)$.
\end{enumerate}
\end{lemma}
The second condition means essentially that $\widehat\psi$ vanishes on $\Om\cap\partial E(\widehat\sigma)$ when considering its trace from the side of $\Om\setminus E(\widehat\sigma)$.
\begin{proof}

We know by Lemma \ref{lem:analicity} that $(\sigma,\psi)$, $\sigma=(\sigma_1,\dots,\sigma_n)$, satisfies the following properties:
\begin{itemize}
\item
Each connected component of $E(\sigma)$ is convex;
\item $\psi$ is positive and real analytic
in $\Om\setminus E(\sigma)$;
\item $\psi=\varphi$ on $\partial^D\Om\setminus\partial E(\sigma)$.
\end{itemize}

In what follows we are going to modify $(\sigma,\psi)$ near each arc of $\partial E(\sigma)$ using an iterative argument in order to get a new minimizer $(\widehat\sigma,\widehat\psi)\in\admclassconv$ that satisfies conditions \ref{lem1.} and \ref{lem2.}. To this aim we denote by $F_1,\dots,F_k$ with $1\le k\le n$ the closure of the connected components of $E(\sigma)$ and set
$ \delta_0:=\min_{i\ne j}\dist(F_i,F_j)>0$. Moreover, by the first property we deduce that $\Om\cap\partial E(\sigma)$ is the union of an at most countable family of pairwise disjoint arcs with endpoints in $\partial\Om$, \ie $\Om\cap \partial E(\sigma)=
\bigcup_{i=1}^k\bigcup_{j=1}^\infty \alpha_{i,j}$, where $\alpha_{i,j}$ is a connected component of $\Om\cap\partial F_i$ for $i\in\{1,
\dots,k\}$, $j\ge1$.\footnote{Notice that at this stage we do not have any information about the geometry of the set $(\partial E(\sigma))\cap \partial\Om$, and $\Om\cap\partial F_i$ could a priori be the union of countably many connected components.}

\subsubsection*{Step 1: Base case} Let $\alpha$ be one of the connected components of $\Om\cap\partial F$, with $F:=F_i$ for some $i\in\{1,\dots,k\}$.
In this step we construct a new minimizer $(\sigma^\alpha,\psi^\alpha)\in\admclassconv$ such that $(\partial E(\sigma^\alpha))\cap\partial\Om=(\partial E(\sigma))\cap\partial\Om$
and
$\psi^\alpha$ is continuous and null on $\alpha'$, where
$\alpha'\subset\Om\cap\partial E(\sigma^\alpha)$ is a suitable curve that replaces $\alpha$ and has the same endpoints as~$\alpha$.
For $\eps\in(0,\delta_0/2)$ we define the stripe
\[
\widehat T_\eps(\alpha):=\{x\in\Om\setminus F\colon\dist(x,\alpha)<\eps\}\subset F^+_\eps\setminus F,
\]
and consider the planar curve $\alpha_\eps$ in $\overline\Om$ defined as in Remark \ref{rem:limite_di_traccia}.
Let $T_\eps(\alpha)$ be the connected component of $\widehat T_\eps(\alpha)$ whose boundary contains $\alpha_\eps$. Let $L_\eps$ be defined as
\[
 L_\eps:=(\partial T_\eps(\alpha))\cap \partial\Om,
\]
so that in particular
$\partial T_\eps(\alpha)=\alpha\cup\alpha_\eps\cup L_\eps $.
Let $p,q\in \partial\Om$ be the endpoints of $\alpha$ (and then also the endpoints of $\alpha_\eps\cup L_\eps$, which are independent of $\eps$). We~define the curves
\[
\Gamma_\eps:=\Gamma_\eps^+\cup\Gamma_\eps^-,\quad
\Gamma_\eps^+:=\mathcal{G}_{\psi\res\alpha_\eps}\cup \mathcal{G}_{\varphi\res L_\eps}\cup l^+,\quad \Gamma_\eps^-:=\mathcal{G}_{-\psi\res\alpha_\eps}\cup \mathcal{G}_{-\varphi\res L_\eps}\cup l^-,
\]
where
\begin{equation*}
l^+:=(\{p\}\times[0,\varphi(p)])\cup (\{q\}\times[0,\varphi(q)]),\quad
l^-:=(\{p\}\times[-\varphi(p),0])\cup(\{q\}\times[-\varphi(q),0]).
\end{equation*}
Observing that $L_\eps\subset\partial^D\Om\setminus\partial E(\sigma)$ and recalling that $\psi=\varphi$ on $\partial^D\Om\setminus\partial E(\sigma)$, we~deduce that
$\Gamma_\eps$ is a closed non-planar curve in $\R^3$ that satisfies assumptions \ref{(1)}--\ref{(3)} of Lemma \ref{lem:plateau_solution}.
Therefore, a solution $S_\eps$ to the classical Plateau problem corresponding to $\Gamma_\eps$ is a disk-type surface such that:
\begin{enumerate}[{\rm (1)},wide]
\item $\beta_{p,q}^\eps:=S_\eps\cap(\R^2\times\{0\})$ is a simple analytic curve joining $p$ and $q$;
\item $S_\eps$ is symmetric with respect to the horizontal plane;
\item the surface $S^+_\eps:=S_\eps\cap\{x_3\ge0\}$
is the graph of a function $\psi_{p,q}^\eps\in W^{1,1}( U_{p,q}^\eps)
\cap C^0(\overline U_{p,q}^\eps\setminus\{p,q\})$, where $U_{p,q}^\eps\subset F\cup T_\eps(\alpha)$ is the open region enclosed between $\alpha_\eps\cup L_\eps$ and $\beta_{p,q}^{\eps}$;
\item the curve $\beta_{p,q}^\eps$ is contained in the closed
convex hull of $\Gamma_\eps$ and $( F\cup T_\eps(\alpha))\setminus U_{p,q}^\eps$ is convex.
\end{enumerate}
We would like to compare the area of $S_\eps^+$ with the area of the generalized graph of $\psi$ on $\overline {T_\eps(\alpha)}$. This is not immediate since, due to the fact that $\psi$ is just $BV$, we~cannot, a priori, conclude that its generalized graph is of disk-type.\footnote{This is due to the jump of $\psi$ on $\partial F$ which is, in general, not regular enough.} Hence we proceed as follows. We~fix $\bar\eps\in(0,\delta_0/2)$; we claim that
\begin{equation}\label{eq:comparison-with-plateau}
\mathcal{A}(\psi_{p,q}^{\bar\eps}; U_{p,q}^{\bar\eps})
\le
\mathcal{A}(\psi; T_{\bar\eps}(\alpha))+
\int_{\alpha}\psi\res{T_{\bar\eps}(\alpha)}\,d\mathcal{H}^1.
\end{equation}
Since $\psi$ is analytic in $T_{\bar\eps}(\alpha)\subset\Om\setminus E(\sigma)$, by Lemma \ref{lem:limite_di_traccia} and Remark \ref{rem:limite_di_traccia} it follows that
\begin{equation}\label{claim}
\lim_{\eps\to0^+}
\int_{\alpha_\eps}\psi\res{T_{\bar\eps}(\alpha)}\,d\mathcal{H}^1=\int_{\alpha}\psi\res{T_{\bar\eps}(\alpha)}\,d\mathcal{H}^1.
\end{equation}
We take
\[
T_{\eps}^{\bar\eps}(\alpha):={T_{\bar\eps}(\alpha)\setminus \overline {T_{\eps}(\alpha)}}\quad\text{and}\quad Y_{\bar\eps}:=S_{\eps}\cup\mathcal{G}_{\psi\res T_\eps^{\bar\eps}(\alpha)} \cup\mathcal{G}_{-\psi\res T_\eps^{\bar\eps}(\alpha)}.
\]
Since $S_\eps$ is a disk-type surface and $\psi$ is analytic in $T_\eps^{\bar\eps}(\alpha)$ it turns out that $Y_{\bar\eps}$ is also a disk-type surface satisfying $\partial Y_{\bar\eps}=\Gamma_{\bar\eps}$. Therefore using that $S_{\bar\eps}$ and $S_\eps$ are solutions to the Plateau problems corresponding to $\Gamma_{\bar\eps}$ and $\Gamma_\eps$ respectively, we~have
\begin{equation*}
\begin{split}
\mathcal{H}^2(S_{\bar\eps})\le \mathcal{H}^2(Y_{\bar\eps})&= 2\mathcal{H}^2(\mathcal{G}_{\psi\res T_\eps^{\bar\eps}(\alpha)} )+\mathcal{H}^2(S_\eps)\\
&\le 2\mathcal{H}^2(\mathcal{G}_{\psi\res T_{\bar\eps}(\alpha)}) +
2\int_{\alpha_\eps\cup L_\eps}\psi\res{T_{\bar\eps}(\alpha)}\, d\mathcal{H}^1\\
&= 2\mathcal{H}^2(\mathcal{G}_{\psi\res T_{\bar\eps}(\alpha)}) + 2\int_{\alpha_\eps}\psi\res{T_{\bar\eps}(\alpha)}\, d\mathcal{H}^1+
2\int_{ L_\eps}\psi\res{T_{\bar\eps}(\alpha)}\,d\mathcal{H}^1.
\end{split}
\end{equation*}
Passing to the limit as $\eps\to0^+$, by \eqref{claim} and the fact that $\mathcal{H}^1(L_\eps)\to0$, we~obtain
\begin{equation*}
\mathcal{H}^2(S_{\bar\eps})\le
2\mathcal{H}^2(\mathcal{G}_{\psi\res T_{\bar\eps}(\alpha)}) +
2\int_{\alpha}\psi\res{T_{\bar\eps}(\alpha)} \,d\mathcal{H}^1,
\end{equation*}
which yields
\begin{align*}
\mathcal{A}(\psi_{p,q}^{\bar\eps};U_{p,q}^{\bar\eps})&=
\mathcal{H}^2(S_{\bar\eps}^+)\le
\mathcal{H}^2(\mathcal{G}_{\psi\res T_{\bar\eps}(\alpha)}) +
\int_{\alpha}\psi\res{T_{\bar\eps}(\alpha)}\,d\mathcal{H}^1\\
&=\mathcal{A}(\psi; T_{\bar\eps}(\alpha))+
\int_{\alpha}\psi\res{T_{\bar\eps}(\alpha)}\,d\mathcal{H}^1,
\end{align*}
and \eqref{eq:comparison-with-plateau} is proved.

We now define $E^\alpha:=(E(\sigma)\cup T_{\bar\eps}(\alpha))\setminus U^{\bar\eps}_{p,q}$ and
\begin{equation*}
\psi^\alpha:=\begin{cases}
0&\text{in}\ E^\alpha,\\
\psi^{\bar\eps}_{p,q}& \text{in}\ U^{\bar\eps}_{p,q}, \\
\psi& \text{otherwise}.
\end{cases}
\end{equation*}
By \eqref{eq:comparison-with-plateau} and using that $U^{\bar\eps}_{p,q}\cup E^\alpha=E(\sigma)\cup T_{\bar\eps}(\alpha)$ we derive
\begin{align}
\mathcal{A}(\psi^\alpha;\Om)-|E^\alpha|
&= \mathcal A(\psi^{\bar\eps}_{p,q};U^{\bar\eps}_{p,q})+
\mathcal{A}(\psi; \Om\setminus(U^{\bar\eps}_{p,q}\cup E^\alpha))\notag
\\
&= \mathcal A(\psi^{\bar\eps}_{p,q};U^{\bar\eps}_{p,q})+
\mathcal{A}(\psi; \Om\setminus(T_{\bar\eps}(\alpha)\cup E(\sigma)))\notag\\[-8pt]\label{eq:area_psi'}\\[-8pt]
&\le \mathcal A(\psi; T_{\bar\eps}(\alpha))+
\int_{\alpha}\psi\res{T_{\bar\eps}(\alpha)}\,d\mathcal{H}^1+\mathcal{A}(\psi; \Omega\setminus T_{\bar\eps}(\alpha))-|E(\sigma)|\notag
\\
&=\mathcal{A}(\psi; \Omega)-|E(\sigma)|.\notag
\end{align}
It remains to construct $\sigma^\alpha\in \Sigma_\mathrm{conv}$. Without loss of generality we may assume
\[
\sigma_1([0,1]),\dots,\sigma_h([0,1])\subset F\quad\text{and}\quad \sigma_{h+1}([0,1]),\dots,\sigma_n([0,1])\not\subset F
\]
for some $h\le n$; notice that if $h=n$ the second family of curves is empty. Then we define $\sigma^\alpha:=(\sigma^\alpha_1,\dots,\sigma_h^\alpha,\sigma_{h+1},\dots,\sigma_n)\in\Lip([0,1];\overline\Om)^n$ as follows: if $h>1$
\begin{equation*}
\sigma^\alpha_{i}([0,1])=\begin{cases}
\overline{q_{i}p_{i+1}}&\text{for } i\leq h-1,\\
\partial (F\cup T_{\bar\eps}(\alpha)\setminus U^{\bar{\eps}}_{p,q})\setminus\Bigl((\bigcup_{i=1}^{h}\partial^0_i\Om)\cup
(\bigcup_{i=1}^{h-1}\overline{q_{i}p_{i+1}})\Bigr)&\text{for } i=h,
\end{cases}
\end{equation*}
where $ \overline{q_{i}p_{i+1}}$ is the segment joining $q_i$ to $p_{i+1}$; if instead $h=1$ we simply set
\begin{equation*}
\sigma^\alpha_{1}([0,1])= \partial (F\cup T_{\bar\eps}(\alpha)\setminus U^{\bar{\eps}}_{p,q})\setminus\partial_1^0\Om.
\end{equation*}
Clearly the pair $(\sigma^\alpha,\psi^\alpha)$ belongs to $\admclassconv$, and by \eqref{eq:area_psi'} it satisfies
\begin{equation*}
\begin{split}
{\mathcal F}(\sigma^\alpha,\psi^\alpha)=
{\mathcal F}(\sigma, \psi).
\end{split}
\end{equation*}
Moreover,
$
(\partial E(\sigma^\alpha))\cap\partial\Om=(\partial E(\sigma))\cap\partial\Om$ and $\psi^\alpha$ is continuous and null on $\alpha'$, where
\begin{equation}\label{gamma'}
\alpha':=\beta_{p,q}^{\bar\eps}\subset \Om\cap\partial E(\sigma^\alpha).
\end{equation}
Summarizing, we~have replaced the curve $\alpha$ with $\alpha'$, ensuring that the new function~$\psi^\alpha$ is now continuous and null on $\alpha'$.

\subsubsection*{Step 2: Iterative case} In this step we construct a minimizer $(\widehat\sigma,\widehat\psi)\in\admclassconv$ of~$\mathcal F$ in~$\mathcal W$ that satisfies the thesis by iterating step one at most a countable number of~times.
We first consider $F=F_1$ and apply step 1 for each $\alpha_{1,j}$ with $j\ge1$.
More precisely we define the pair $(\sigma_{1,j},\psi_{1,j})\in\admclassconv$ as follows:
\begin{itemize}
\item if $j=1$ we set\vspace*{-3pt}
\begin{equation*}
(\sigma_{1,1},\psi_{1,1}):=(\sigma^{\alpha_{1,1}},\psi^{\alpha_{1,1}} ),
\end{equation*}
where $(\sigma^{\alpha_{1,1}},\psi^{\alpha_{1,1}} )\in\admclassconv$ is a minimizer constructed as in step 1 with $\alpha=\alpha_{1,1}$;
\item if $j>1$ we set\vspace*{-3pt}
\begin{equation*}
(\sigma_{1,j},\psi_{1,j}):= (\sigma_{1,j-1}^{\alpha_{1,j}},\psi_{1,j-1}^{{\alpha}_{1,j}}),
\end{equation*}
where $(\sigma_{1,j-1}^{\alpha_{1,j}},\psi_{1,j-1}^{{\alpha}_{1,j}})\in\admclassconv$ is a minimizer constructed as in step 1 with $(\sigma,\psi)=(\sigma_{1,j-1},\psi_{1,j-1})$ and $\alpha=\alpha_{1,j}$.
\end{itemize}
Since
$
\mathcal{F}(\sigma_{1,j},\psi_{1,j})=\mathcal F(\sigma,\psi)
$
for all $j\ge1$, by Lemma \ref{lem:compactness} it follows that $(\sigma_{1,j},\psi_{1,j})$ converges to $(\sigma_1,\psi_1)\in\admclassconv$ in the sense of Definition \ref{def:conv}. Moreover, by construction we have that for every $j\ge1$ the pair $(\sigma_{1,j},\psi_{1,j})$ satisfies
\begin{equation*}
(\partial E(\sigma_{1,j}))\cap\partial\Om= (\partial E(\sigma))\cap\partial\Om,
\end{equation*}
and $\psi_{1,j}$ is continuous and null on $\bigcup_{h=1}^j\alpha'_{1,h}\subset \Om\cap(\partial E(\sigma_{1,j}))\cap\partial F_1$,
where $\alpha'_{1,h}$ are defined as in \eqref{gamma'}. As a consequence $(\sigma_1,\psi_1)$ satisfies
\begin{equation*}
(\partial E(\sigma_{1}))\cap\partial\Om= (\partial E(\sigma))\cap\partial\Om,
\end{equation*}
and $\psi_1$ is continuous and null on $\bigcup_{j=1}^\infty\alpha'_{1,j}\subset \Om\cap(\partial E(\sigma_{1}))\cap\partial F_1$.
Moreover,\vspace*{-3pt}
\begin{equation*}
\Om\cap \partial E(\sigma_1)=\biggl(\bigcup_{j=1}^\infty\alpha'_{1,j})\bigcup
(\bigcup_{i=2}^k\bigcup_{j=1}^\infty \alpha_{i,j}\biggr).
\end{equation*}
Now repeating the argument above for the pair $(\sigma_1,\psi_1)$ and $i=2$ we obtain a new minimizer $(\sigma_2,\psi_2)\in\admclassconv$ of $\mathcal F$ in $\mathcal W$ satisfying\vspace*{-3pt}
\begin{equation*}
(\partial E(\sigma_{2}))\cap\partial\Om= (\partial E(\sigma))\cap\partial\Om,
\end{equation*}
with $\psi_2$ continuous and null on $\bigcup_{j=1}^\infty(\alpha'_{1,j}\cup\alpha'_{2,j})\subset \Om \cap (\partial E(\sigma_{1}))\cap((\partial F_1)\cup\partial F_2)$
and\vspace*{-3pt}
\begin{equation*}
\Om \cap (\partial E(\sigma_2))=\biggl(\bigcup_{i=1}^2\bigcup_{j=1}^\infty\alpha'_{i,j})\cup
(\bigcup_{i=3}^k\bigcup_{j=1}^\infty \alpha_{i,j}\biggr).
\end{equation*}
Iterating this process a finite number of times we finally get a minimizer $(\widehat\sigma,\widehat\psi)\in\admclassconv$ of $\mathcal F$ in $\mathcal W$ with the required properties.
\end{proof}

We are finally in the position to conclude the proof of Theorem \ref{teo:structure_of_minimizers}.

\begin{proof}[Proof of Theorem \ref{teo:structure_of_minimizers}]
Let $(\sigma,\psi)\in\admclassconv$ be any minimizer of $\mathcal F$ in $\mathcal W$ as in Theorem \ref{teo:reduction_from_W_to_Wconv}.
By Lemma \ref{lem:analicity} we know that $(\sigma,\psi)$ satisfies properties \ref{1.}, \ref{2.} and the boundary datum is attained, namely
\begin{equation*}
\psi=\varphi\quad\text{on}\quad\partial^D\Om\setminus\partial E(\sigma).
\end{equation*}
Moreover, by Lemma \ref{lem:continuita-bordo-E} there is a minimizer $(\widehat\sigma,\widehat\psi)\in\admclassconv$ such that
\begin{equation}\label{eq:prop1}
\partial E(\widehat{\sigma})\cap\partial\Om= \partial E(\sigma)\cap\partial\Om,
\end{equation}
and $\widehat\psi$ is continuous and null on $\Om\cap\partial E(\widehat\sigma)$.

It remains to show that if $\partial^D_i\Om$ is not straight for some $i=1,
\dots,n$, then
\[
\partial E(\sigma)\cap\partial^D_i\Om=\partial E(\widehat\sigma)\cap\partial^D_i\Om=\emptyset,
\]
and if instead $\partial^D_i\Om$ is straight for some $i=1,\dots,n$, then property \ref{4.} holds.
Eventually we show that there is a minimizer that satisfies property \ref{5.}.
This will be achieved in a number of steps.

\subsubsection*{Step 1} Assuming that there is $i\in\{1,\dots,n\}$ such that $\partial_i^D\Om$ is not straight, we~show that $E(\widehat\sigma)\cap\partial_i^D\Om=\emptyset$. To prove this we proceed by analyzing three different cases.

\begin{figure}
\centering
\def\svgwidth{0.9\textwidth}
\input{bellettini-et-al_figures/caseA-caseB.pdf_tex}
\caption{Case A. $\partial_i^D\Omega\cap\partial E(\hat\sigma)=\arc{ab}$. The
dotted curve upon $\Om$ represents $\Gamma^+$ in \eqref{def:curve}.
Case B. $\partial_i^D\Omega\cap\partial E(\hat\sigma)=\{c\}$. The dotted curve upon $\Om$ represents the curve $\Gamma^+$ in \eqref{eq:gamma}.
}\label{caseA-caseB}
\end{figure}

\subsubsection*{Case A}
Suppose, to the contrary,
that there is a non-straight\footnote{Namely, $\arc{ab}$ is not contained in a line.} arc $\arc{ab}$ (with endpoints $a\ne b$) in $\partial_i^D\Om\cap\partial E(\widehat\sigma)$ (Case A in Figure \ref{caseA-caseB}). Thus in particular $\arc{ab}\subset\bigcup_{j=1}^n\widehat\sigma_j([0,1])$. We~may assume without loss of generality that $\arc{ab}\subset\widehat\sigma_1([0,1])$.
Then we consider the curves
\begin{equation}\label{def:curve}
\Gamma:=\Gamma^+\cup\Gamma^-,\quad
\Gamma^+:=\mathcal{G}_{\varphi\res\arc{ab}}\cup l^+,\quad \Gamma^-:=\mathcal{G}_{-\varphi\res\arc{ab}}\cup l^-,
\end{equation}
where
\begin{equation*}
l^+:=(\{a\}\times[0,\varphi(a)])\cup (\{b\}\times[0,\varphi(b)]),\quad
l^-:=(\{a\}\times[-\varphi(a),0])\cup (\{b\}\times[-\varphi(b),0]).
\end{equation*}
In this way $\Gamma$ satisfies the assumptions of Lemma \ref{lem:plateau_solution} and hence a solution $S$ to the Plateau problem spanning $\Gamma$ is a disk-type surface such that:
\begin{enumerate}[(i),wide]
\item\label{i} $\beta_{a,b}:=S\cap (\R^2\times\{0\})$
is a simple analytic curve joining $a$ and $b$;
\item\label{ii}$S$ is symmetric with respect to $\R^2\times\{0\}$;
\item \label{iii} the surface $S^+:=S\cap\{x_3\ge0\}$
is the graph of a function $\psi_{a,b}\in W^{1,1}( U_{a,b})\cap C^0(\overline U_{a,b}\setminus\{a,b\})$, where $U_{a,b}\subset E(\widehat\sigma_1)$ is the open region enclosed between $\arc{ab}$ and $\beta_{a,b}$;
\item\label{iv} the curve $\beta_{a,b}$ is contained in the closed
convex hull of $\Gamma$ and $E(\widehat\sigma_1)\setminus U_{a,b}$ is convex.
\end{enumerate}
The inclusion $U_{a,b}\subset E(\widehat\sigma_1)$ follows since $\arc{ab}\subset\widehat\sigma_1([0,1])$, $E(\widehat\sigma_1)$ is convex, and $S$ is contained in the convex envelope of $\Gamma$.
Furthermore by the minimality of $S$ one has
\begin{equation}\label{eq:contradiction}
\mathcal{A}(\psi_{a,b};U_{a,b})=\mathcal{H}^2(S^+)< \int_{\arc{ab}}\varphi\,d\mathcal{H}^1= \int_{\arc{ab}}|\widehat\psi-\varphi|\,d\mathcal{H}^1.
\end{equation}
Here the strict inequality follows since the vertical wall spanning $\Gamma$ given by
\[
\{(x',x_3)\colon x'\in\arc{ab},\ x_3\in[-\varphi(x'),\varphi(x')]\}
\]
is a disk-type surface but, since $\arc{ab}$ is not a segment, it cannot be a solution to the Plateau problem. We~now consider the pair $(\widetilde\sigma,\widetilde\psi)\in\admclassconv$ given by
\begin{equation}\label{def:sigma,psi}
\widetilde\sigma:=(\widetilde\sigma_1,\widehat\sigma_2,\dots,\widehat\sigma_n),\qquad
\widetilde\psi:=\begin{cases}
0&\text{in }\widetilde E,\\
\psi_{a,b}&\text{in }U_{a,b},\\
\widehat \psi&\text{otherwise},
\end{cases}
\end{equation}
where $\widetilde\sigma_1$ is such that $\widetilde\sigma_1([0,1])=(\widehat\sigma_1([0,1])\setminus\arc{ab})\cup\beta_{a,b}$ and $\widetilde{E}:= E(\widehat\sigma)\setminus U_{a,b}=E(\widetilde\sigma)$.
Then noticing that $\widehat\psi=0$ in $U_{a,b}$, $E(\widehat\sigma)=E(\widetilde\sigma)\cup U_{a,b}$, and recalling \eqref{eq:contradiction}, we~get
\begin{equation*}
\begin{split}
{\mathcal F}(\widetilde\sigma,\widetilde\psi)&=\mathcal{A}(\widetilde\psi;\Om)-|E(\widetilde\sigma)|+\int_{\partial\Om}|\widetilde\psi-\varphi|\,d\mathcal{H}^1\\
&=\mathcal{A}(\widehat\psi;\Om\setminus U_{a,b})+\mathcal{A}(\psi_{a,b};U_{a,b})-|E(\widetilde\sigma)|+\int_{\partial\Om}|\widetilde\psi-\varphi|\,d\mathcal{H}^1\\
&=\mathcal{A}(\widehat\psi;\Om)+\mathcal{A}(\psi_{a,b};U_{a,b})-|E(\widehat\sigma)|+\int_{\partial\Om}|\widehat\psi-\varphi|\,d\mathcal{H}^1\\
&<\mathcal{A}(\widehat\psi;\Om)-|E(\widehat\sigma)|+\int_{\partial\Om}|\widetilde\psi-\varphi|\,d\mathcal{H}^1+\int_{\arc{ab}}|\widehat\psi-\varphi|\,d\mathcal{H}^1\\
&=\mathcal{A}(\widehat\psi;\Om)-|E(\widehat\sigma)|+\int_{\partial\Om}|\widehat\psi-\varphi|\,d\mathcal{H}^1= {\mathcal F}(\widehat\sigma,\widehat\psi),
\end{split}
\end{equation*}
where the penultimate equality follows from the fact that $\widetilde\psi$ is continuous and equal to $\varphi$ on $\arc{ab}$ while the traces of $\widetilde\psi$ and $\widehat\psi$ coincide on $\partial\Om\setminus\arc{ab}$.
This contradicts the minimality of $(\widehat\sigma,\widehat\psi)$.

\subsubsection*{Case B} Suppose by contradiction that the set $\partial_i^D\Om\cap\partial E(\widehat\sigma)$ contains an isolated point $c$ or has a straight segment $\overline{cc'}$ as isolated connected component (Case B in Figure \ref{caseA-caseB}).
Then there are two arcs $\arc{ab}\subset\partial_i^D\Om$ and $\arc{a'b'}\subset\partial E(\widehat\sigma)$ with either $a\neq a'$ or $b\neq b'$ (and with endpoints $a\ne b$ and $a'\ne b'$) such that $\overline{aa'}\cap\overline{bb'}=\emptyset$ and $\arc{ab}\cap\arc{a'b'}=\{c\}$ (respectively $\arc{ab}\cap\arc{a'b'}=\overline{cc'}$). Notice also that, since $\partial_i^D\Om$ is not straight, the segment $\overline{cc'}$ does not coincide with $\partial_i^D\Om$ and hence the arc $\arc{ab}$ can be chosen so that it properly contains the segment $\overline{cc'}$. We~consider the curves $\Gamma:=\Gamma^+\cup\Gamma^-$ with
\begin{equation}\label{eq:gamma}
\Gamma^+:=\mathcal{G}_{\varphi\res\arc{ab}}
\cup\mathcal{G}_{\widehat\psi\res\overline{aa'}}\cup\mathcal{G}_{\widehat\psi\res\overline{bb'}}
,\quad
\Gamma^-:=\mathcal{G}_{-\varphi\res\arc{ab}}
\cup\mathcal{G}_{-\widehat\psi\res\overline{aa'}}
\cup\mathcal{G}_{-\widehat\psi\res\overline{bb'}}.
\end{equation}
Notice that $\Gamma^\pm$ connect $a'$ to $b'$.
By applying again Lemma \ref{lem:plateau_solution} to the nonplanar curve~$\Gamma$ and arguing as in case A we obtain the contradiction also in this case.

\subsubsection*{Case C} More generally, assume by contradiction that both the sets $\partial_i^D\Om\cap\partial E(\widehat\sigma)$ and $\partial_i^D\Om\setminus\partial E(\widehat\sigma)$ are nonempty.
Then we can find a not flat arc $\arc{ab}\subset\partial_i^D\Om$ such that the following holds:\footnote{This is a consequence of the fact that $\arc{ab}\setminus\partial E(\widehat\sigma)$ is relatively open in $\arc{ab}$, so it is an at most countable union of disjoint relatively open arcs.} there are pairs of points $\{c_j,d_j\}_{j\in \mathbb N}\subset\partial_i^D\Om\cap\partial E(\widehat\sigma)$ such that the arcs $\arc{ad_0}$, $\arc{c_0b}$, and $\{\arc{c_jd_j}\}_{j=1}^\infty$ are mutually disjoint and
\begin{equation*}
\arc{ab}\setminus\partial E(\widehat\sigma)=\arc{ad_0}\cup\Bigl(\bigcup_{j=1}^\infty\arc{c_jd_j}\Bigr)\cup\arc{c_0b}.
\end{equation*}
Without loss of generality, we~might assume that all the points $c_j,d_j\in\widehat\sigma_1([0,1])$.
For all $j\ge1$ we denote by $V_j$ the region enclosed by $\arc{c_jd_j}$ and $\partial E(\widehat\sigma)$.\footnote{These regions are simply connected since $c_j,d_j\in\widehat\sigma_1([0,1])$.} We~now argue as in case B and choose $a',b'\in \widehat\sigma_1([0,1])$. Additionally, let $V_0=V_0^a\cup V_0^b$, with~$V_0^a$ (respectively $V_0^b$) be the region enclosed between $\partial E(\widehat\sigma)$ and $\overline{aa'}\cup\arc{ad_0}$ ($\partial E(\widehat\sigma)$ and $\overline{bb'}\cup\arc{c_0b}$, respectively). We~finally define $\Gamma$ correspondingly, as in \eqref{eq:gamma}.
Again by Lemma \ref{lem:plateau_solution} the solution $S$ to the Plateau problem corresponding to $\Gamma$ satisfies the properties \ref{i}--\ref{iv} considered in case A, with $a'$ and $b'$ in place of $a$ and $b$ respectively.
Moreover, by the minimality of $S$ for every $N\ge1$ there holds\footnote{The right-hand side is the area of the surface given by the (positive) subgraph of $\varphi$ on $\arc{ab}\setminus
\bigcup_{j=1}^N\arc{c_jd_j}$ and the graph of $\widehat\psi$ on the region $\bigcup_{j=0}^NV_j$, which is of disc-type.
To see this we use that the trace of $\widehat\psi$ on the
subarcs of $\partial E(\widehat\sigma)$ between the points $c_j$ and $d_j$
is zero (and between $a'$ and $d_0$, and $d_0$ and $b'$).}
\begin{multline}\label{eq:contradiction_bis}
\mathcal{A}(\psi_{a',b'};U_{a',b'})=\mathcal{H}^2(S^+)\\
\le \int_{\arc{ab}}\varphi\,d\mathcal{H}^1-\int_{\arc{ad_0}\cup\arc{c_0b}}\!\varphi\,d\mathcal{H}^1-\sum_{j=1}^N\int_{\arc{c_jd_j}}\!\varphi\,d\mathcal{H}^1+\sum_{j=0}^N \mathcal{A}(\psi;V_j).
\end{multline}
In~particular, by taking the limit as $N\to+\infty$ in \eqref{eq:contradiction_bis} we get
\begin{equation}\label{eq:contradiction_tris}
\mathcal{A}(\psi_{a',b'};U_{a',b'})=\mathcal{H}^2(S^+)\le \int_{\arc{ab}\setminus\partial E(\widehat\sigma)}\varphi\,d\mathcal{H}^1+\textstyle\mathcal{A}(\widehat\psi;\bigcup_{j=0}^\infty V_j).
\end{equation}
Let $(\widetilde\sigma,\widetilde\psi)\in\admclassconv$ be defined as in \eqref{def:sigma,psi}, then observing that $\widehat\psi=0$ in $U_{a',b'}\setminus(\bigcup_{j=0}^\infty V_j)$, $E(\widehat\sigma)=E(\widetilde\sigma)\cup (U_{a',b'}\setminus\bigcup_{j=0}^\infty V_j)$ and using \eqref{eq:contradiction_tris} we deduce
\begin{equation*}
\begin{split}
{\mathcal F}(\widetilde\sigma,\widetilde\psi)
&=\mathcal{A}(\widehat\psi;\Om\setminus U_{a',b'})+\mathcal{A}(\psi_{a',b'};U_{a',b'})-|E(\widetilde\sigma)|+\int_{\partial\Om}|\widetilde\psi-\varphi|\,d\mathcal{H}^1\\
&=\mathcal{A}(\widehat\psi;\Om\setminus({\textstyle\bigcup_{j=0}^\infty V_j)})+\mathcal{A}(\psi_{a',b'};U_{a',b'})-|E(\widehat\sigma)|+\int_{\partial\Om}|\widetilde\psi-\varphi|\,d\mathcal{H}^1\\
&\le\mathcal{A}(\widehat\psi;\Om\setminus({\textstyle\bigcup_{j=0}^\infty V_j)})-|E(\widehat\sigma)|\\
&\qquad +\int_{\partial\Om}|\widetilde\psi-\varphi|\,d\mathcal{H}^1+\int_{\arc{ab}\cap\partial E(\widehat\sigma)}\varphi\,d\mathcal{H}^1+{\textstyle\mathcal{A}(\widehat\psi;\bigcup_{j=0}^\infty V_j)}\\
&=\mathcal{A}(\widehat\psi;\Om)-|E(\widehat\sigma)|+\int_{\partial\Om}|\widehat\psi-\varphi|\,d\mathcal{H}^1= {\mathcal F}(\widehat\sigma,\widehat\psi),
\end{split}
\end{equation*}
which in turn implies\vspace*{-3pt}
\begin{equation}\label{eq:contradiction_four}
{\mathcal F}(\widetilde\sigma,\widetilde\psi)\le
{\mathcal F}(\widehat\sigma,\widehat\psi).
\end{equation}
To conclude we need to show that the inequality in \eqref{eq:contradiction_four} is strict. To this aim, we~choose $c\in\{c_j\}_{j=1}^\infty$.
Consider the curves $\Gamma_1$ and $\Gamma_2$ defined as follows
\begin{equation*}
\Gamma_1:=\Gamma_1^+\cup\Gamma_1^-,\quad \Gamma_1^+:=\mathcal{G}_{\varphi\res\arc{ac}}
\cup\mathcal{G}_{\widehat\psi\res\overline{aa'}}\cup l^+
,\quad
\Gamma_1^-:=\mathcal{G}_{-\varphi\res\arc{ac}}
\cup\mathcal{G}_{-\widehat\psi\res\overline{aa'}}\cup l^-,
\end{equation*}
\begin{equation*}
\Gamma_2:=\Gamma_2^+\cup\Gamma_2^-,\quad \Gamma_2^+:=\mathcal{G}_{\varphi\res\arc{cb}}
\cup\mathcal{G}_{\widehat\psi\res\overline{bb'}}\cup l^+
,\quad
\Gamma_2^-:=\mathcal{G}_{-\varphi\res\arc{cb}}
\cup\mathcal{G}_{-\widehat\psi\res\overline{bb'}}\cup l^-,
\end{equation*}
where\vspace*{-3pt}
\begin{equation*}
l^+:=(\{c\}\times[0,\varphi(c)]),\quad
l^-:=(\{c\}\times[-\varphi(c),0]).
\end{equation*}
Let $S_1$ and $S_2$ be the solutions to the Plateau problem corresponding to~$\Gamma_1$ and~$\Gamma_2$ respectively, so that properties \ref{i}--\ref{iv} are satisfied with $c$ in place of $b'$ and $a'$ respectively.
By the minimality of $S$ we have
\begin{equation}\label{eq:ab}
\mathcal{A}(\psi_{a',b'};U_{a',b'})<\mathcal{A}(\psi_{a',c};U_{a',c})+\mathcal{A}(\psi_{c,b'};U_{c,b'}).
\end{equation}\label{eq:ac}
On the other hand by arguing as above,\footnote{With the arc $\arc{ac}$ ($\arc{cb}$, respectively) in place of $\arc{ab}$.} we conclude
\begin{equation}
\mathcal{A}(\psi_{a',c};U_{a',c})\le\int_{\arc{ac}\cup\partial E(\widehat\sigma)}\varphi\,d\mathcal{H}^1+\mathcal{A}(\widehat\psi;\bigcup_{j\in I_1} V_j\cup V^a_0),
\end{equation}
and\vspace*{-3pt}
\begin{equation}\label{eq:cb}
\mathcal{A}(\psi_{c,b'};U_{c,b'})\le\int_{\arc{cb}\cup\partial E(\widehat\sigma)}\varphi\,d\mathcal{H}^1+\mathcal{A}(\widehat\psi;\bigcup_{j\in I_2} V_i\cup V^b_0),
\end{equation}
where $I_1:=\{j\colon \arc{c_jd_j}\subset\arc{ac}\}$ and $I_2:=\{j\colon \arc{c_jd_j}\subset\arc{cb}\}$. Gathering together \eqref{eq:ab}--\eqref{eq:cb} we derive\vspace*{-3pt}
\begin{equation*}
\mathcal{A}(\psi_{a',b'};U_{a',b'})<
\int_{\arc{ab}\cup\partial E(\widehat\sigma)}\varphi\,d\mathcal{H}^1+\mathcal{A}(\widehat\psi;\bigcup_{j=0}^\infty V_j),
\end{equation*}
which in turn implies\vspace*{-3pt}
\begin{equation*}
{\mathcal F}(\widetilde\sigma,\widetilde\psi)<
{\mathcal F}(\widehat\sigma,\widehat\psi),
\end{equation*}
and thus the contradiction.

\subsubsection*{Step  2} Assuming there is $i\in\{1,\dots,n\}$ such that $\partial_i^D\Om$ is a straight segment, we~show that either $(\partial E(\widehat\sigma))\cap\partial_i^D\Om=\emptyset$ or $(\partial E(\widehat\sigma))\cap\partial_i^D\Om=\partial_i^D\Om$.
Suppose by contradiction that $(\partial E(\widehat\sigma))\cap\partial_i^D\Om\ne\emptyset$ and also $\partial_i^D\Om\setminus \partial E(\widehat\sigma)\neq \emptyset$. Without loss of generality we can restrict to the case $(\partial E(\widehat\sigma))\cap\partial_i^D\Om=(\partial F)\cap\partial_i^D\Om$ with $F$ any connected component of $E(\widehat\sigma)$. Since $F$ is convex and $\partial_i^D\Om$ is a segment $(\partial F)\cap\partial_i^D\Om$ has to be connected, \ie it is either a single point $a$ or a segment $\overline{aa'}\neq \partial_i^D\Om$.
In both cases we then consider a (small enough) ball $B$ centered at $a$ such that $B\cap E(\widehat\sigma)=B\cap F$ (in the second case we also require that the radius of $B$ is smaller than $\overline{aa'}$).

If $(\partial F)\cap\partial_i^D\Om=\{a\}$ we let $\{p,q\}:=(\partial B)\cap\partial F$ and $\{b,c\}:=(\partial B)\cap\partial_i^D\Om$ (with $b,p$ and $c,q$ lying on the same side with respect to $a$). Then
we define the curves
\begin{equation*}
\Gamma:=\Gamma^+\cup\Gamma^-,\quad \Gamma^+:=\mathcal{G}_{\varphi\res\overline{bc}}\cup \mathcal{G}_{\psi\res\arc{bp}} \cup \mathcal{G}_{\psi\res\arc{cq}},\quad \Gamma^-:=\mathcal{G}_{-\varphi\res\overline{bc}}\cup \mathcal{G}_{-\psi\res\arc{bp}}
\cup \mathcal{G}_{-\psi\res\arc{cq}},
\end{equation*}
where $\arc{bp}$, $\arc{cq}$ denote the arcs in $\partial B$ joining $b$ to $p$ and $c$ to $q$ respectively.

If $(\partial F)\cap\partial_i^D\Om=\overline{aa'}$ we let $\{p,q\}:=(\partial B)\cap\partial F$ and $\{b,c\}:=(\partial B)\cap\partial_i^D\Om$, where we identify $q$ and $c$. Then we consider the curves $\Gamma:=\Gamma^+\cup\Gamma^-$ with
\begin{equation*}
\Gamma^+:=\mathcal{G}_{\varphi\res\overline{bc}}\cup \mathcal{G}_{\psi\res\arc{bp}} \cup (\{c\}\times[0,\varphi(c)]),\quad\! \Gamma^-:=\mathcal{G}_{-\varphi\res\overline{bc}}\cup \mathcal{G}_{-\psi\res\arc{bp}}
\cup (\{c\}\times[-\varphi(c),0]).
\end{equation*}
By applying again Lemma \ref{lem:plateau_solution} to $\Gamma$ and arguing as above we get the contradiction.

\subsubsection*{Step 3} We show that there is a minimizer $(\widetilde\sigma,\widetilde\psi)$ that satisfies property \ref{5.}. We~first notice that $\widehat\psi$ is continuous and null on $\partial E(\widehat\sigma)\setminus\partial^D\Om$.
Moreover, by steps 1 and 2 it follows that $\Om\cap \partial E(\widehat\sigma)$ is the union of a finite number of pairwise disjoint Lipschitz curves each of them joining each $p_i$ for $i=1,\dots,n$ to each of the $q_j$ for some $j=1,\dots,n$.
To conclude it is enough to replace each curve, without increasing the energy, with an analytic one having the same endpoints.
More precisely, let $\gamma$ be any of such curves. Reasoning as in the proof of Lemma \ref{lem:continuita-bordo-E} step 1, we~can replace $(\widehat\sigma,\widehat\psi)$ with a new minimizer $(\sigma^\gamma,\psi^\gamma)\in\admclassconv$ such that $(\partial E(\sigma^\gamma))\cap\partial\Om=(\partial E(\sigma))\cap\partial\Om$
and
$\psi^\gamma=0$ on $\gamma'$, where
$\gamma'\subset(\partial E(\sigma^\gamma))\cap\Om$ is a suitable analytic curve that replaces $\gamma$ and has the same endpoints of $\gamma$.
In~particular, $\psi^\gamma$ is continuous and null on $\partial E(\sigma^\gamma)\setminus\partial^DR_{2\ell}$.
Eventually iterating this procedure for each curve in $\partial E(\widehat\sigma)\setminus\partial\Om$ we can construct a new minimizer
$(\widetilde\sigma,\widetilde\psi)$ with the required properties.
\end{proof}

\subsection{The example of the catenoid containing a segment}
\label{subsec:the_example_of_the_catenoid_containing_a_segment}
Consider the setting of
Figure \ref{figura1mezzo}. Recall that
$\Omega=R_{2\ell}=(0,2\ell)\times (-1,1)$,
$n=1$, $\partial^D\Om=(\{0,2\ell\}\times (-1,1))\cup ((0,2\ell)\times \{-1\})$ and $\partial^0\Om=(0,2\ell)\times \{1\}$, $p=(0,1)$, $q=(2\ell,1)$.
The map $\varphi$ given in \eqref{vortex} is $\varphi(z_1,z_2)=\sqrt{1-z_2^2}$ on $\partial^D
\Om$, and thus
vanishes on $[0,2\ell]\times \{-1\}$; for this reason this case is not covered by our analysis. However we can find a solution as in Theorem~\ref{teo_main_intro} also in this case, by an approximation procedure.
Precisely, for $\eps>0$ consider an approximating sequence $(\varphi_\eps)$ of continuous Dirichlet data, with~$\mathcal G_{\varphi_\eps}$ Lipschitz, which tends to $\varphi$ uniformly and satisfies $\varphi_\eps=0$ on $\partial^0\Om$, $\varphi_\eps>0$ on $\partial^D\Om$. Let $(\sigma_\eps,\psi_\eps)$ be a solution as in Theorem \ref{thm:existence} corresponding to the boundary datum $\varphi_\eps$; since $\mathcal F(\sigma_\eps,\psi_\eps)$ is equibounded,\footnote{We can bound it from above by $|\Om|+\int_{\partial_D\Om}|\varphi_\eps|d\mathcal H^1$.} arguing as in the proof of Lemma \ref{lem:compactness}, we~can see that, up to
a subsequence, $((\sigma_\eps,\psi_\eps))$ tends to some $(\sigma,\psi)\in \admclassconv$, which minimizes the functional $\mathcal F$ with Dirichlet condition $\varphi$.
In this case however we cannot guarantee that $\sigma$ does not touch $\partial^D\Om$, even if this is not a straight segment.
This is essentially due to the presence of the
portion $[0,2\ell]\times \{-1\}$ of $\partial \Om$
where $\varphi$ is zero, which does not allow to apply the arguments used in the proof of Theorem \ref{teo:structure_of_minimizers}.\enlargethispage{2pt}

In~particular, it can be seen that if $\ell$ is large enough,
the solution $(\sigma,\psi)$ splits and becomes degenerate,
being $\psi\equiv0$ and the value of $\mathcal F$ is
just the area of two vertical half-disks of radius $1$. For $\ell$ under a certain threshold, instead, the solution satisfies the regularity properties stated in Theorem \ref{teo:structure_of_minimizers}, and in particular $\psi=\varphi$ on $\partial^D\Omega$, and~$\sigma$ is the graph of a smooth convex function passing through $p$ and $q$. We~refer to~\cite{BES2} for details and comprehensive proofs of these facts;
we also notice that in this special case further regularity of solutions
can be obtained.

\section{Comparison with the parametric Plateau problem: the case \texorpdfstring{$n=1,2$}{12}}
\label{sec:comparison_with_the_parametric_Plateau_problem:the_case_n=1,2}

In this section we compare the solutions of
Theorems \ref{teo:reduction_from_W_to_Wconv} and \ref{teo:structure_of_minimizers} with the solutions to the classical Plateau problem in parametric form.
Specifically, motivated by the example of the catenoid, we~restrict our analysis to the classical
disk-type and annulus-type Plateau problem.
These configurations correspond to the cases $n=1$ and $n=2$ respectively, \ie the Dirichlet boundary $\partial^D\Om$ is either an open arc or
the union of two open arcs of $\partial\Om$ with disjoint closure.
Due to the highly involved geometric arguments,
we do not discuss the case $n>2$, which requires further investigation.
Thus, in this section we assume $n=1,2$. We~first discuss the case $n=1$, which is a consequence of
Lemma \ref{lem:plateau_solution}, and next the case $n=2$.

\subsection{The case $n=1$}\label{sec:the_case_n=1}

Let $n=1$. Let $p_1,q_1\in\partial\Om$, $\partial^D\Om=\partial^D_1\Om$, $\varphi$ be as in Section~\ref{subsec:setting_of_the_problem} and
consider the space curve $\gamma_1:=\mathcal{G}_{\varphi\res\partial_1^D\Om}$ joining $p_1$ to $q_1$. We~define the curve
\[
\Gamma:=\gamma_1\cup \Sym(\gamma_1),
\]
where $\Sym(\gamma_1):=\mathcal{G}_{-\varphi\res\partial_1^D\Om}$,
and
consider the classical Plateau problem in parametric form
spanning $\Gamma$. More precisely we look for a solution to
\begin{equation}\label{plateau_1}
m_1(\Gamma):=\inf_{\Phi\in\mathcal{P}_1(\Gamma)}\int_{B_1}|\partial_{w_1}\Phi\wedge\partial_{w_2}\Phi|dw,
\end{equation}
where
\begin{multline}\label{def:adm_Phi}
\mathcal{P}_1(\Gamma):=\bigl\{\Phi\in H^1( B_1;\R^3)\cap C^0(\overline B_1;\R^3)\text{ such that }\\\Phi\res\partial B_1\colon\partial B_1\to\Gamma \text{ is a weakly monotonic parametrization of $\Gamma$}
\bigr\}.
\end{multline}
By classical arguments, it is well-known that
every solution to \eqref{plateau_1} is a harmonic and conformal parametrization of an area-minimizing surface spanning $\Gamma$.

\begin{theorem}[The disk-type Plateau problem ($n=1$)]
\label{teo:the_disk_type_Plateau_problem_n=1}
Assume $\Gamma$ is not planar, let $\Phi\in\mathcal{P}_1(\Gamma)$ be a solution to \eqref{plateau_1} and let
\[
S^+:=\Phi(\overline B_1)\cap \{x_3\ge0\},
\qquad \qquad S^-:=\Phi(\overline B_1)\cap \{x_3\le0\}.
\]
Then there exists a minimizer $(\sigma,\psi)\in \admclassconv$ of $\mathcal F$ in
$\mathcal W$
satisfying properties \ref{1.}--\ref{5.} of Theorem \ref{teo:structure_of_minimizers} and such that
\begin{equation}\label{eq:graph}
S^{\pm}=\mathcal G_{\pm\psi\res (\overline{\Om\setminus E(\sigma)})}.
\end{equation}
Conversely let $(\sigma,\psi)\in \admclassconv$ be a minimizer of $\mathcal F$ in $\mathcal W$ satisfying properties \ref{1.}--\ref{5.} of Theorem \ref{teo:structure_of_minimizers}. Then the disk-type surface
\[
S:=\mathcal G_{\psi\res (\overline{\Om\setminus E(\sigma)})}\cup\mathcal G_{-\psi\res (\overline{\Om\setminus E(\sigma)})}
\]
is a solution to the classical Plateau problem associated to $\Gamma$, \ie there is a harmonic and conformal map $\Phi\in\mathcal{P}_1(\Gamma)$
solving \eqref{plateau_1} and such that $\Phi(\overline{B_1})=S$.
\end{theorem}
{ We have assumed $\Gamma$ is not planar, otherwise the classical solution is flat, and any solution to Theorem \ref{teo:structure_of_minimizers} satisfies $(\partial E(\sigma))\cap \partial^D\Om=\partial^D\Om$.}

\subsection{The case $n=2$}\label{sec:the_case_n=2}
Let $n=2$. Let $\Om$, $p_1,q_1,p_2,q_2\in\partial\Om$, $\partial^D\Om$, $\partial_1^D\Om$, $\partial_2^D\Om$, $\varphi$ be as in Section \ref{subsec:setting_of_the_problem} and
consider the space curve $\gamma_i:=\mathcal{G}_{\varphi\res\partial_i^D\Om}$ joining $p_i$ to $q_i$ for $i=1,2$. We~define the curves
\begin{equation*}
\Gamma_1:=\gamma_1\cup\Sym(\gamma_1),\qquad \Gamma_2:=\gamma_2\cup\Sym(\gamma_2),
\end{equation*}
where $\Sym(\gamma_i):=\mathcal{G}_{-\varphi\res\partial_i^D\Om}$ for $i=1,2$. We~consider the
classical Plateau problem in parametric form spanning the curve
\[
\Gamma:=\Gamma_1\cup\Gamma_2.
\]
Precisely we set $\openannulus\subset\R^2$ to be an open annulus
enclosed between two concentric circles { $C_1:=\partial B_1(0)$ and $C_2:=\partial B_2(0)$},
and we look for a solution to
\begin{equation}\label{catenoid_plateau}
m_2(\Gamma):=\inf_{\Phi\in\mathcal{P}_2(\Gamma)}\int_{\openannulus}|\partial_{w_1}\Phi\wedge\partial_{w_2}\Phi|dw,
\end{equation}
where
\begin{multline*}
\mathcal{P}_2(\Gamma):=\big\{\Phi\in H^1(\openannulus;\R^3)\cap C^0(\overline\Sigma_\mathrm{ann};\R^3) \text{ such that }\Phi(\partial \openannulus)=\Gamma\text{ and }\\
\Phi\res C_j:C_j\to \Gamma_j\text{ is a weakly monotonic parametrization of $\Gamma_j$, $j=1,2$}\big\}.
\end{multline*}

Here the crucial assumption that we require is that the curves $\Gamma_j$
have the orientation inherited by the orientation\footnote{Once we fix
an orientation of $\partial \Omega$, the orientation of the graph
$\mathcal G_\varphi$
of $\varphi$ is inherited, since~$\mathcal G_\varphi$
is defined in a standard way as the push-forward of the current of integration on $\partial_D\Om$ by the map $x\mto(x,\varphi(x))$.} of the graph of $\varphi$ on $\partial_j^D\Om$.

Due to the specific geometry of $\Gamma$ we can appeal to Theorem \ref{myexistence} below (which is a consequence of \cite[Th.\,1 \& Th.\,5]{MY}) to deduce the existence of a minimizer. This might not be true for a more general $\Gamma$. To this purpose for $j=1,2$ we consider
the minimization problem defined in \eqref{plateau_1} for the curve $\Gamma_j$, namely
\begin{align}\label{minimizationSigma}
m_1(\Gamma_j)=\inf_{\Phi\in\mathcal{P}_1(\Gamma_j)}\int_{B_1}|\partial_{w_1}\Phi\wedge\partial_{w_2}\Phi|dw,
\end{align}
with $\mathcal{P}_1(\Gamma_j)$ defined as in \eqref{def:adm_Phi}.

By standard arguments one sees that $m_2(\Gamma)\le m_1(\Gamma_1)+m_1(\Gamma_2)$. Indeed, two disk-type surfaces can be joined by a thin tube (with arbitrarily small area) in order to change the topology of the two disks into an annulus-type surface.

\begin{definition}[$\mathcal{MY}$ solution]\label{def:MY_solution}
Let $\Phi\in\mathcal{P}_2(\Gamma)$ be a solution to \eqref{catenoid_plateau}. We~say that~$\Phi$ is a $\mathcal{MY}$ solution to \eqref{catenoid_plateau} if
$\Phi$ is harmonic, conformal, and it is an embedding. In~particular, in such a case, $m_2(\Gamma)=\mathcal{H}^2(\immclosedann)$.
\end{definition}
\begin{theorem}[Meeks and Yau]\label{myexistence}
Suppose $m_2(\Gamma)<m_1(\Gamma_1)+m_1(\Gamma_2)$. Then there exists a $\mathcal{MY}$ solution $\Phi\in\mathcal{P}_2(\Gamma)$ to \eqref{catenoid_plateau}. Furthermore, every minimizer of \eqref{catenoid_plateau} is a $\mathcal{MY}$ solution.
\end{theorem}

\begin{proof}
See \cite{MY}.
\end{proof}
This result allows us to prove the following:
\begin{theorem}[The annulus-type Plateau problem ($n=2$)]
\label{teo:the_annulus_type_Plateau_problem_n=2}
The following holds:
\begin{enumerate}[{\rm(i)},wide]
\item \label{1plateau} Suppose $m_2(\Gamma)<m_1(\Gamma_1)+m_1(\Gamma_2)$.
Let $\Phi\in\mathcal{P}_2(\Gamma)$
be a $\mathcal{MY}$ solution to \eqref{catenoid_plateau} and let
\begin{equation*}
S:=\immclosedann, \qquad
S^+:=S\cap\{x_3\ge0\},\qquad
S^-:=S\cap\{x_3\le0\}.
\end{equation*}
Then there exists a minimizer $(\sigma,\psi)\in \admclassconv$
of $\mathcal F$ in $\mathcal W$ satisfying properties \ref{1.}--\ref{5.} of Theorem \ref{teo:structure_of_minimizers} and such that
\begin{equation}\label{param_to_nonparam}
S^{\pm}=\mathcal G_{\pm\psi\res (\overline{\Om\setminus E(\sigma)})}.
\end{equation}
\item\label{2plateau} Suppose $m_2(\Gamma)= m_1(\Gamma_1)+m_1(\Gamma_2)$, { and assume that both $\Gamma_1$ and $\Gamma_2$ are not planar}. For $j=1,2$ let $\Phi_j\in\mathcal{P}_1(\Gamma_j)$ be a solution to \eqref{minimizationSigma} and let $S_j:=\Phi_j(\overline B_1)$. Let also
\begin{equation*}
S^+:=(S_1\cup S_2)\cap\{x_3\ge0\}\qquad\text{and}\qquad
S^-:=(S_1\cup S_2)\cap\{x_3\le0\}.
\end{equation*}
Then $S_1\cap S_2=\emptyset$ and there exists a minimizer $(\sigma,\psi)\in \admclassconv$ of $\mathcal F$ in $\mathcal W$ satisfying properties \ref{1.}--\ref{5.} of Theorem \ref{teo:structure_of_minimizers} and such that \eqref{param_to_nonparam} holds.
\item \label{3plateau}
Conversely, let $(\sigma,\psi)\in \admclassconv$ be a minimizer of $\mathcal F$ in $\mathcal W$ satisfying properties \ref{1.}--\ref{5.} of Theorem \ref{teo:structure_of_minimizers}. Then the surface
\[
S:=\mathcal G_{\psi\res (\overline{\Om\setminus E(\sigma)})}\cup\mathcal G_{-\psi\res (\overline{\Om\setminus E(\sigma)})}
\]
is either an annulus-type surface or the union of two disjoint disk-type surfaces, and is a solution to the classical Plateau problem associated to $\Gamma$. More precisely, either there is a $\mathcal{MY}$ solution $\Phi\in\mathcal{P}_2(\Gamma)$ to \eqref{catenoid_plateau} with
$S=\immclosedann$, or there are $\Phi_j\in\mathcal{P}_1(\Gamma_j)$ solutions to \eqref{minimizationSigma} for $j=1,2$, such that
\[
S=\Phi_1(\overline B_1)\cup \Phi_2(\overline B_1)\quad\text{and}\quad\Phi_1(\overline B_1)\cap \Phi_2(\overline B_1)=\emptyset.
\]
\end{enumerate}
\end{theorem}

\subsection{Toward the proof of Theorems \ref{teo:the_disk_type_Plateau_problem_n=1}
and \ref{teo:the_annulus_type_Plateau_problem_n=2}: preliminary lemmas}

In order to prove Theorems \ref{teo:the_disk_type_Plateau_problem_n=1} and
\ref{teo:the_annulus_type_Plateau_problem_n=2}, we~collect some technical lemmas.

\begin{lemma}[Graphicality of minimizers for $n=2$]\label{lem_tec1}
Let $n=2$, and $(\sigma,\psi)\in\admclassconv$ be a minimizer
of $\mathcal F$ in
$\admclass$
satisfying properties \ref{1.}--\ref{5.} of Theorem \ref{teo:structure_of_minimizers}.
\begin{enumerate}[{\rm(a)},wide]
\item\label{lem_tec1a}
Suppose that $\overline{\Om\setminus E(\sigma)}$ is connected.
Then there exists an injective map $\Phi\in W^{1,1}(\openannulus;\R^3)\cap C^0(\overline{\Sigma}_{\mathrm{ann}};\R^3)$ such that
\[
\immclosedann=\mathcal G_{\psi\res (\overline{\Om\setminus E(\sigma)})}\cup\mathcal G_{-\psi\res (\overline{\Om\setminus E(\sigma)})},
\]
and $\Phi\res C_j\colon C_j\to \Gamma_j$ is a weakly monotonic parametrization of $\Gamma_j$ for $j=1,2$.
\item\label{lem_tec1b}
Suppose that $\Om\setminus E(\sigma)$
consists of two connected components, whose closures~$F_1$ and~$F_2$ are disjoint,
with
$F_j \supseteq \partial^D_j\Om$ for $j=1,2$.
Then there exist two injective maps $\Phi_1,\Phi_2\in W^{1,1}(B_1;\R^3)\cap C^0(\overline{B_1};\R^3)$ such that
\[
\Phi_j(\overline{B_1})=\mathcal G_{\psi\res {F_j}}
\cup\mathcal G_{-\psi\res{F_j}},\qquad j=1,2,
\]
and $\Phi_j\res \partial B_1\colon\partial B_1\to \Gamma_j$ is a weakly monotonic parametrization of $\Gamma_j$ for $j=1,2$.
\end{enumerate}
\end{lemma}
Supposing that $\Om\setminus E(\sigma)$
has two connected components as in \ref{lem_tec1b}, it readily follows that $\Gamma_1$
and $\Gamma_2$
cannot be planar
(otherwise the solution will be flat on $\partial^D_j\Om$ and $F_j=\emptyset$
for $j=1,2$).

\skpt
\begin{proof}
\ref{lem_tec1a}
Since $\overline{\Om\setminus E(\sigma)}$ is simply connected,\footnote{ This is the region enclosed by $\partial^D\Om\cup\sigma_1([0,1])\cup\sigma_2([0,1])$.} the
maps
\begin{equation}\label{eq:Psi_tilde}
\widetilde \Psi^\pm\in W^{1,1}({\Om\setminus E(\sigma)};\R^3)
\cap C^0(\overline{\Om\setminus E(\sigma)};\R^3),
\qquad \widetilde \Psi^\pm(p):=(p,\pm\psi(p)),
\end{equation}
are disk-type parametrizations of $\mathcal G_{\pm\psi\res (\overline{\Om\setminus E(\sigma)})}$, thanks to properties \ref{1.}--\ref{5.} of Theorem \ref{teo:structure_of_minimizers}.

Now, by using a homeomorphism of class $H^1$ between $\overline{\Om\setminus E(\sigma)}$ and a disk, we~can parametrize\footnote{For instance,
we can consider a (flat) disk-type Plateau solution
spanning $\partial(\Om\setminus E(\sigma))$.
Then we can employ a Lipschitz homeomorphism between the disk and the
half-annulus.}
$\overline{\Om\setminus E(\sigma)}$ with a half-annulus, obtained as the region enclosed between two concentric half-circles
with endpoints $A_1,A_2,A_3,A_4$ (in the order) on the same diameter, and the two segments $\overline{A_1A_2}$ and $\overline{A_3A_4}$.
Then we construct a parametrization $\Psi^+$
of $\mathcal G_{\psi\res (\overline{\Om\setminus E(\sigma)})}$ as in \eqref{eq:Psi_tilde} from the
half-annulus, such that $\Psi^+(A_1)=(q_1,0)$, $\Psi^+(A_2)=(p_2,0)$, $\Psi^+(A_3)=(q_2,0)$, $\Psi^+(A_4)=(p_1,0)$, and mapping weakly monotonically the two
half-circles
into $\gamma_1$ and $\gamma_2$, and the two segments into $\sigma_1([0,1])$ and
$\sigma_2([0,1])$, respectively.
Similarly, we~construct a parametrization $\Psi^-$ of $\mathcal G_{-\psi\res (\overline{\Om\setminus E(\sigma)})}$ from
another copy of a half-annulus, just
setting $\Psi^-:=\Sym(\Psi^+)$,
the symmetric of~$\Psi^+$ with respect to the plane containing $\Omega$.

Eventually, gluing the two half-annuli
along the two segments, we~get a parametrization $\Phi$
of $\mathcal G_{\psi\res (\overline{\Om\setminus E(\sigma)})}
\cup \mathcal G_{-\psi\res (\overline{\Om\setminus E(\sigma)})}$
defined on $\overline\Sigma_\mathrm{ann}$.
By the continuity of $\psi$ on $\partial^D\Omega$ we have that $ \Phi$ parametrizes $\Gamma_i$ on $C_i$, $i=1,2$.

\ref{lem_tec1b}
It is sufficient to argue as in case \ref{lem_tec1a},
by replacing $\Omega\setminus E(\sigma)$ in turn with $F_1$ and $F_2$ and
$\openannulus$ with $B_1$ to find $\Phi_1$ and $\Phi_2$, respectively.
\end{proof}

\begin{lemma}\label{lem_tec2}
Let $n=2$, and $(\sigma,\psi)\in\admclassconv$ be a minimizer of
$\mathcal F$ in
$\mathcal W$
satisfying properties \ref{1.}--\ref{5.} of Theorem \ref{teo:structure_of_minimizers}.
\begin{enumerate}[{\rm(a)},wide]
\item\label{lem_tec2a}
Suppose that $\overline{\Om\setminus E(\sigma)}$ is connected and
\begin{equation}\label{eq:lem_tec2}
\mathcal H^2(\mathcal G_{\psi\res (\overline{\Om\setminus E(\sigma)})}\cup\mathcal G_{-\psi\res (\overline{\Om\setminus E(\sigma)})})\leq m_2(\Gamma).
\end{equation}
Let $\Phi$ be the parametrization given by Lemma \ref{lem_tec1}\ref{lem_tec1a}.
Then there exists a reparametrization of the annulus $\openannulus$
such that, using it to reparametrize $\Phi$, the corresponding map (still denoted by $\Phi$)
belongs to $\mathcal P_2(\Gamma)$ and solves \eqref{catenoid_plateau}.
\item\label{lem_tec2b}
Suppose that $\Om\setminus E(\sigma)$
consists of two connected components whose closures~$F_1$ and~$F_2$ are disjoint,
$F_j \supseteq \partial_j^D \Om$ for $j=1,2$, and
\[
\mathcal H^2(\mathcal G_{\psi\res {F_j}}
\cup\mathcal G_{-\psi\res {F_j}})\leq m_1(\Gamma_j),
\qquad j=1,2.
\]
Let $\Phi_1,\Phi_2$ be the maps given by Lemma \ref{lem_tec1}\ref{lem_tec1b}. Then, for $j=1,2$, there is a reparametrization
of $\Phi_j$ belonging to $\mathcal P_1(\Gamma_j)$ and solving \eqref{minimizationSigma}.
\end{enumerate}
\end{lemma}

\skpt
\begin{proof}
\ref{lem_tec2a}
Fix a point $\widetilde p\in \Om\setminus E(\sigma)$ and set $\widetilde \Psi^+_k:=\widetilde\Psi^+\res H_k$,
where $\widetilde \Psi$ is defined in \eqref{eq:Psi_tilde}
and, for $k\in \mathbb N$ sufficiently large, $H_k$ is the
connected component of
\[
\widetilde H_k:=\{p\in \overline{\Om\setminus E(\sigma)}:\textup{dist}(p,\partial(\overline{\Om\setminus E(\sigma)}) )\geq 1/k\}
\]
containing $\widetilde p$.
For $k \in \mathbb N$ large enough $H_k$ is
simply connected with rectifiable boundary, thanks to the simply-connectedness of $ \Om\setminus E(\sigma)$.
In~particular, $\widetilde \Psi^+_k$ parametrizes a disk-type surface, and
using the regularity of $ \psi$ in ${\Om\setminus E(\sigma)}$,
it follows that $\widetilde \Psi^+_k$ is Lipschitz continuous.
Furthermore,
$\widetilde \Psi^+_k\res \partial H_k$
parametrizes a Jordan curve,
and these curves, suitably parametrized,
converge in the sense of Fréchet (see \cite[Th.\,4, \S 4.3]{DHS})
as $k\to +\infty$, to the curve having image
$\widetilde \Psi^+(\partial ({\Om\setminus E(\sigma)}) ) )=:\lambda$.
Notice that
\begin{align}\label{lambda}
\lambda=\sigma_1([0,1])\cup\sigma_2([0,1])\cup\gamma_1\cup \gamma_2.
\end{align}
Call $\lambda_k$ the image of the curve given by
$\widetilde \Psi^+_k\res \partial H_k$.
Let $\mathcal{P}_1(\lambda_k)$, $\mathcal{P}_1(\lambda)$, $m_1(\lambda_k)$, $m_1(\lambda)$ be defined as in \eqref{def:adm_Phi} and \eqref{plateau_1} with $\lambda_k$ and $\lambda$ in place of $\Gamma$ respectively.
Up to reparametrizing $\overline B_1$ (see footnote 15),
$\widetilde\Psi^+_k$ belongs to $\mathcal{P}_1(\lambda_k)$, therefore
\[
\mathcal H^2( \mathcal G_{\psi\res H_{k}})=\int_{H_k}|\partial_{w_1}\widetilde \Psi_k^+\wedge \partial_{w_2}\widetilde \Psi_k^+ |dw\geq m_1(\lambda_k)\qquad\forall k\geq1.
\]
We claim that equality holds in the previous expression, namely
\begin{align}\label{claim_speranza}
\mathcal H^2( \mathcal G_{\psi\res H_k})=m_1(\lambda_k)\qquad \forall k\geq1.
\end{align}
Indeed, assume by contradiction that $\mathcal H^2( \mathcal G_{\psi\res H_{k_0}})>m_1(\lambda_{k_0})$ for some $k_0\geq1$,
and pick $\delta>0$ with
\begin{equation}\label{eq:absurd}
\mathcal H^2( \mathcal G_{\psi\res H_{k_0}})\geq \delta+
m_1(\lambda_{k_0}).
\end{equation}
Take $\Phi_{k_0}\in \mathcal{P}_1(\lambda_{k_0})$
a solution to $m_1(\lambda_{k_0})$.
For $k> k_0$, as $H_{k_0}\subset H_k$, by a
gluing argument,\footnote{This is done, for instance, by gluing an external annulus to a disk, and using $\Phi_{k_0}$ from the disk, and a reparametrization of $\mathcal G_{\psi\res(H_k\setminus H_{k_0})}$ from the annulus.} we~can find $\Phi_k\in\mathcal{P}_1(\lambda_k)$ such that $\Phi_k(\overline{B_1})=\Phi_{k_0}(\overline{B_1})\cup \mathcal G_{\psi\res(H_k\setminus H_{k_0})}$. Thus by \eqref{eq:absurd} we have
\[
\begin{aligned}
\mathcal H^2( \mathcal G_{\psi\res H_{k}})
&\geq
\delta + m_1(\lambda_{k_0})
+
\mathcal H^2( \mathcal G_{\psi\res (H_k\setminus H_{k_0})})
\\
&=
\delta + \mathcal H^2(\Phi_{k_0}(\overline B_1))
+
\mathcal H^2( \mathcal G_{\psi\res (H_k\setminus H_{k_0})})
\geq
\delta+ m_1(\lambda_k)
\qquad \forall k> k_0.
\end{aligned}
\]
Letting $k\to +\infty$, since $\lambda_k\to\lambda$ in the sense of Fréchet, we~have $m_1(\lambda_k)\to m_1(\lambda)$ \cite[Th.\,4, \S 4.3]{DHS}.
In~particular, from the previous inequality we infer
\begin{equation*}
\mathcal{F}(\sigma,\psi)=\mathcal{H}^2(\mathcal G_{\psi\res(\overline{\Om\setminus E(\sigma)})})\ge\delta+m_1(\lambda).
\end{equation*}
Hence we conclude
\[
\mathcal{H}^2(\mathcal G_{\psi\res(\overline{\Om\setminus E(\sigma)})}\cup \mathcal G_{-\psi\res(\overline{\Om\setminus E(\sigma)})})\geq 2\delta+2m_1(\lambda)\geq 2\delta+m_2(\Gamma),
\]
which contradicts \eqref{eq:lem_tec2}.
In the last inequality we have used that $2m_1(\lambda)\geq m_2(\Gamma)$; this follows from the fact that a disk-type parametrization of a minimizer for $m_1(\lambda)$ can be reparametrized on a half-annulus (as in the proof of Lemma \ref{lem_tec1}), and glued with another reparametrization of it on the other half-annulus, so to obtain a parametrization of an annulus-type surface spanning $\Gamma$ which is admissible for \eqref{catenoid_plateau}.
Hence claim \eqref{claim_speranza} follows.

Now, since $\psi$ is Lipschitz continuous on $\overline H_k$, for
all $k\in \mathbb N$ sufficiently large there exists a
map $\Psi_k\in H^1(B_1;\R^3)\cap C^0(\overline {B_1};\R^3)$ with $\Psi_{k}(\partial B_1)=\lambda_{k}$ monotonically which solves the classical disk-type
Plateau problem {spanning} $\lambda_k$ and such that
\[
\Psi_k(B_1)=\mathcal G_{\psi\res H_k}.
\]
Letting $k\to +\infty$
and using that the Dirichlet energy of $\Psi_k$ equals the area of $\mathcal G_{\psi\res H_k}$, we
conclude that $(\Psi_k)$ tends to a map $\Psi\in H^1(B_1;\R^3)\cap C^0(\overline {B_1};\R^3)$
with $\Psi(\partial B_1)=\lambda$ {weakly} monotonically, and that is a solution of the classical disk-type Plateau problem with
\[
\Psi(\overline B_1)=\mathcal{G}_{\psi\res(\overline{\Om\setminus E(\sigma)})}.
\]
Arguing as in the proof of Lemma \ref{lem_tec1} we
finally get a map $\Phi:
\overline\Sigma_\mathrm{ann}\to \R^3$ which belongs to $\mathcal P_2(\Gamma)$ and parametrizes $\mathcal{G}_{\psi\res(\overline{\Om\setminus E(\sigma)})}\cup \mathcal{G}_{-\psi\res(\overline{\Om\setminus E(\sigma)})}$.
This concludes the proof of \ref{lem_tec2a}.

\ref{lem_tec2b}
It is sufficient to argue as in case \ref{lem_tec2a},
by replacing $\Omega\setminus E(\sigma)$ in turn with $F_1$ and $F_2$ and
$\openannulus$ with $B_1$ to find $\Phi_1$ and $\Phi_2$, respectively.
\end{proof}

Using the arguments above to show conditions \ref{lem_tec1b} of Lemma \ref{lem_tec1} and \ref{lem_tec2b} of Lem\-ma~\ref{lem_tec2}, we~deduce the following:

\begin{cor}\label{lem_tec3}
Let $n=1$, assume that $\Gamma$ is not planar, and let $(\sigma,\psi)\in\admclassconv$ be a minimizer of
$\mathcal F$ in
$\mathcal W$
satisfying properties \ref{1.}--\ref{5.} of Theorem \ref{teo:structure_of_minimizers}. Then there exists an injective map $\Phi\in W^{1,1}(B_1;\R^3)\cap C^0(\overline{B_1};\R^3)$ such that
\[
\Phi(\overline{B_1})=\mathcal G_{\psi\res {(\overline{\Om\setminus E(\sigma)})}}
\cup\mathcal G_{-\psi\res{\overline{(\Om\setminus E(\sigma))}}},
\]
and $\Phi\res \partial B_1\colon\partial B_1\to \Gamma$ is a weakly monotonic parametrization of $\Gamma$. Moreover, if $\mathcal H^2(\mathcal G_{\psi\res (\overline{\Om\setminus E(\sigma)})}
\cup\mathcal G_{-\psi\res(\overline{\Om\setminus E(\sigma)})})\leq m_1(\Gamma)$ then there is a reparametrization
of $\Phi$ belonging to $\mathcal P_1(\Gamma)$ and solving \eqref{minimizationSigma}.
\end{cor}

Now we can start the proof of Theorems \ref{teo:the_disk_type_Plateau_problem_n=1} and
\ref{teo:the_annulus_type_Plateau_problem_n=2}.

\Subsection{Proof of Theorem \ref{teo:the_disk_type_Plateau_problem_n=1}}

\begin{proof}[Proof of Theorem \ref{teo:the_disk_type_Plateau_problem_n=1}] Let $\Phi\in\mathcal{P}_1(\Gamma)$ be a solution to \eqref{plateau_1}.
The curve $\Gamma$
satisfies the assumptions
of Lemma \ref{lem:plateau_solution} (notice that in this case we have $f(p_1) = f(q_1)=0$), hence the minimal disk-type surface $S:=\Phi(\overline{B_1})$ satisfies the following properties:
\begin{itemize}
\item $\beta_{p_1,q_1}:=S\cap (\R^2\times\{0\})\subset \overline \Om$
is a simple analytic curve
joining $p_1$ and $q_1$ and such that $\beta_{p_1,q_1}\cap\partial\Om=\{p_1,q_1\}$;
\item $S$ is symmetric with respect to $\R^2\times\{0\}$;
\item the surface $S^+=S\cap\{x_3\ge0\}$ is the graph of a
function $\widetilde\psi\in W^{1,1}(U_{p_1,q_1})\cap C^0(\overline U_{p_1,q_1})$, where $U_{p_1,q_1}\subset \Om$ is the open region enclosed between $\partial_1^D\Om$ and $\beta_{p_1,q_1}$. Moreover, $\widetilde\psi$ is analytic in $U_{p_1,q_1}$;
\item the curve $\beta_{p_1,q_1}$ is contained in the closed
convex hull of $\Gamma$, and $\Om\setminus U_{p_1,q_1}$ is convex.
\end{itemize}

Let $(\sigma,\psi)\in\admclassconv$ be given by
\[
\sigma:=\sigma_1\qquad\text{and}\qquad
\psi:=\begin{cases}
0& \text{in } \Om\setminus U_{p_1,q_1},\\
\widetilde\psi & \text{in } U_{p_1,q_1},
\end{cases}
\]
where $\sigma_1([0,1])=\beta_{p_1,q_1}$.
Clearly \eqref{eq:graph} holds, and $\mathcal H^2(S)=2\mathcal F(\sigma,\psi)=m_1(\Gamma)$. It~remains to show that $(\sigma,\psi)$ is a minimizer of $\mathcal F$.
Let $(\sigma',\psi')\in\admclassconv$ be a minimizer of $\mathcal F$ that satisfies properties \ref{1.}--\ref{5.} of Theorem \ref{teo:structure_of_minimizers} and consider the disk-type surface with boundary $\Gamma$ given by $S':=\mathcal{G}_{\psi'\res(\overline{\Om\setminus E(\sigma')})}
\cup \mathcal{G}_{-\psi'\res(\overline{\Om\setminus E(\sigma')})}$.
Since $(\sigma,\psi)$ is admissible for $\mathcal F$, we~deduce
\[
\mathcal{H}^2(S')=2\mathcal F(\sigma',\psi')\leq m_1(\Gamma).
\]
Thus we are in the hypotheses of Corollary \ref{lem_tec3} and so there is a map
$\Phi'\in \mathcal P_1(\Gamma) $ with $\Phi'(\overline B_1)=S'$.
By minimality of $(\sigma',\psi')$ and of $S$ we have
\begin{equation}
\mathcal H^2(S)\le \mathcal{H}^2(S')=2\mathcal{F}(\sigma',\psi')\le 2\mathcal{F}(\sigma,\psi)=\mathcal H^2(S).
\end{equation}
Hence $(\sigma,\psi)$ is a minimizer of $\mathcal F$ in $\mathcal W$ and $\Phi'$ is a solution to \eqref{plateau_1}.

Conversely, let $(\sigma,\psi)\in \admclassconv$ be a solution that
satisfies properties \ref{1.}--\ref{5.} of Theorem \ref{teo:structure_of_minimizers} and let $S:=\mathcal{G}_{\psi\res(\overline{\Om\setminus E(\sigma)})}
\cup \mathcal{G}_{-\psi\res(\overline{\Om\setminus E(\sigma)})}$.
Let $\widetilde\Phi$ be a solution to \eqref{plateau_1}; then we can find $(\widetilde\sigma,\widetilde\psi)\in \mathcal W$ whose
doubled graph $\widetilde S=\mathcal G_{\widetilde\psi\res (\overline{\Om\setminus E(\widetilde\sigma)})}\cup\mathcal G_{-\widetilde\psi\res (\overline{\Om\setminus E(\widetilde\sigma)})}$ satisfies
\[
 \mathcal H^2(S)=2\mathcal F(\sigma,\psi)\leq2\mathcal F(\widetilde\sigma,\widetilde\psi)=\mathcal H^2(\widetilde S)=m_1(\Gamma).
\]
Arguing as before we find a map $\Phi\in \mathcal P_1(\Gamma)$ parametrizing $S$. We~conclude that $\Phi$ is a solution to \eqref{plateau_1}, and the theorem is proved.
\end{proof}

\subsection{Proof of Theorem \ref{teo:the_annulus_type_Plateau_problem_n=2}}
The proof of
Theorem \ref{teo:the_annulus_type_Plateau_problem_n=2}
is much more involved, so we divide it in a number of steps. We~start with a result (which can be seen as the counterpart of
Lemma \ref{lem:plateau_solution} for the Plateau problem defined in \eqref{catenoid_plateau}) that will be crucial to prove \ref{1plateau}.
In what follows we denote by $\pi\colon\R^3\to\R^2\times\{0\}$ the orthogonal projection.
\begin{theorem}\label{crucial_teo}
Suppose $m_2(\Gamma)<m_1(\Gamma_1)+m_1(\Gamma_2)$ and let
$\Phi\in\mathcal{P}_2(\Gamma)$ be a $\mathcal{MY}$ solution to \eqref{catenoid_plateau}.
Then the minimal surface $\immclosedann$ satisfies the following properties:
\begin{enumerate}[label=$(\arabic*)$,wide]
\item \label{crucial1} The set $\pi(\immclosedann)$
is simply connected in $\overline\Omega$; $\Om\cap\partial \pi(\immclosedann)$
consists of two disjoint embedded analytic curves $\beta_1$ and $\beta_2$ joining $q_1$ to $p_2$, and $q_2$ to $p_1$, respectively.
Moreover, for
$i=1,2$,
the closed region $E_i$
enclosed between $\partial_i^0\Om$ and $\beta_i$
is convex;
\item\label{crucial2} $\immclosedann$ is symmetric with respect to the plane $\R^2\times\{0\}$;
\item\label{crucial3} $\immclosedann\cap(\R^2\times\{0\})=\beta_1\cup\beta_2$;
\item\label{crucial4} $S^+:=\immclosedann\cap \{x_3\ge0\}$ is Cartesian. Precisely, it is
the graph of a function $\widetilde\psi\in W^{1,1}(\mathrm{int}(\pi(\immclosedann)))\cap C^0({\pi(\immclosedann)})$.
\end{enumerate}
\end{theorem}
The proof of Theorem \ref{crucial_teo} is a consequence of Lemmas \ref{lem:simply_connectedness}, \ref{lem:trace_on_the_horizontal_plane}, \ref{lem:region_enclosed_by_Phi(Sigma_min)}, \ref{lem:graphicality_of_boundary_and_continuity_up_to_the_boundary} , \ref{lem:construction-of-parametrization}, and \ref{lemma_step5} below.
\begin{lemma}[Simply connectedness]
\label{lem:simply_connectedness}
Suppose $m_2(\Gamma)<m_1(\Gamma_1)+m_1(\Gamma_2)$ and let
$\Phi\in\mathcal{P}_2(\Gamma)$ be a $\mathcal{MY}$ solution to \eqref{catenoid_plateau}.
Then $\pi(\immclosedann)$ is a simply connected region in $\overline\Omega$ and { $\pi(\immclosedann)\cap\partial\Om=\partial_1^D\Om\cup\partial_2^D\Om$}.
\end{lemma}
\begin{proof}
We recall that $\Phi:\overline\Sigma_\mathrm{ann}\to \R^3$ is an embedding. The fact that $\pi(\immclosedann)$ is a subset of $\overline\Omega$ and $\pi(\immclosedann)\cap\partial\Om=\partial_1^D\Om\cup\partial_2^D\Om$ follows from the fact that the interior of $\immclosedann$ is contained in the convex hull of $\Gamma$.
So it remains to show that $\pi(\immclosedann)$ is simply connected.

Suppose by contradiction that $\pi(\immclosedann)$ is not simply connected.
Let $H$ be a hole of it, namely a region in $\Om$ surrounded by a loop contained in $\pi(\immclosedann)$ and such that $H\cap
\pi(\immclosedann)=\emptyset$; choose a point $P\in H$. We~will look for a contradiction by exploiting that $\openannulus$ is an annulus and using that the map $\Phi$ is analytic and harmonic.

Let $\theta$ be the angular coordinate of a cylindrical coordinate system $(\rho,\theta,z)$ in $\R^3$ centered at $P$ and with $z$-axis the vertical line $\pi^{-1}(P)$. For $\theta \in [0,2\pi)$ we consider the
half-plane orthogonal to $\R^2\times\{0\}$ defined by
\[
\Pi_\theta:=\{(\rho,\theta,z)\colon\rho>0,z\in\R\}.
\]
\begin{figure}
\begin{center}
\def\svgwidth{0.40\textwidth}
\input{bellettini-et-al_figures/fig1_bis.pdf_tex}
\caption{The horizontal section of
two half-planes $\Pi_{\theta_1}$ and $\Pi_{\theta_2}$ intersecting $\partial^0_1\Om$ and $\partial^0_2\Om$, respectively.
}\label{fig-costruzione}
\end{center}
\end{figure}
Now we fix two values $\theta_1$ and $\theta_2$ so that $\Pi_{\theta_1}$ and $\Pi_{\theta_2}$ intersect (the interior of) $\partial_1^0\Om$ and $\partial_2^0\Om$ respectively. The half-planes\footnote{The angles are considered $(\textup{mod}\; 2\pi)$.} $\Pi_{\theta_1+\pi}$ and $\Pi_{\theta_2+\pi}$ might intersect $\partial^D\Om$ (see Figure \ref{fig-costruzione}). However, since the points $p_1$, $q_1$, $p_2$, $q_2$, are in clockwise order on $\partial \Om$, and~$\Omega$ is convex, it is not difficult to conclude the following assertion:

The half-planes $\Pi_{\theta_1+\pi}$ and $\Pi_{\theta_2+\pi}$ cannot intersect the two components $\partial_1^D\Om$ and $\partial_2^D\Om$ of $\partial^D\Om$ at the same time.

In other words: if, for instance, $\Pi_{\theta_1+\pi}$ intersects $\partial_1^D\Om$, then $\Pi_{\theta_2+\pi}$ does not intersect $\partial_2^D\Om$. Let us prove the assertion in the form of the last statement, being the other
cases similar. This is trivial since, if $\Pi_{\theta_1}$ intersects $\partial_1^0\Om$ and $\Pi_{\theta_1+\pi}$ intersects $\partial_1^D\Om$ (as~in Figure \ref{fig-costruzione}), we~have that $\Pi_{\theta}$ intersects $\partial_1^D\Om\cup\partial_1^0\Om$ for all $\theta\in [\theta_1,\theta_1+\pi]$. As either $\theta_2$ or $\theta_2+\pi$ belongs to $[\theta_1,\theta_1+\pi]$, we~have that $\Pi_{\theta_2}\cup\Pi_{\theta_2+\pi}$ intersects $\partial_1^D\Om\cup\partial_1^0\Om$. Since by hypothesis $\Pi_{\theta_2}$ intersects $\partial_2^0\Om$, it follows that $\Pi_{\theta_2+\pi}$ does not intersect $\partial_2^D\Om$, and the statement follows.

Moreover, since $\Pi_{\theta_1}$ intersects $\partial^0_1\Om$ and $\Pi_{\theta_2}$ intersects $\partial^0_2\Om$, it is straightforward that, if $\Pi_{\theta_1+\pi}$ intersects $\partial_1^0\Om$ then also $\Pi_{\theta_2+\pi}$ intersects $\partial_1^0\Om$.

We are now ready to conclude the proof of the lemma. We~have to discuss the following cases:
\begin{enumerate}[(1),wide]
\item\label{case:new1}
$\Pi_{\theta_1+\pi}$ intersects $\partial^0\Om$;
\item\label{case:new2}
$\Pi_{\theta_1+\pi}$ intersects $\partial_1^D\Om$;
\item\label{case:new3}
$\Pi_{\theta_1+\pi}$ intersects $\partial_2^D\Om$.
\end{enumerate}
By hypothesis on $P$, for all $\theta\in[0,2\pi)$ the intersection between $\immclosedann$ and $\Pi_\theta$ consists of a
family of smooth simple curves, either closed or with endpoints on $\Gamma$. Correspondingly, $\Phi^{-1}(\immclosedann\cap \Pi_\theta)$ is a family of closed curves in
$\overline\Sigma_\mathrm{ann}$, possibly with endpoints on $C_1\cup C_2$.
In~particular, since $\Pi_{\theta_1} \cap \partial^0_1\Om \neq \emptyset$, the set\footnote{Since $\Pi_{\theta_1} \cap \partial^D\Om = \emptyset$ these curves must be closed in $\openannulus$.} $\Phi^{-1}
(\immclosedann\cap \Pi_{\theta_1})$ is a family of closed curves in $\openannulus$.

In case \ref{case:new1} also, $\Phi^{-1}(\immclosedann\cap \Pi_{\theta_1+\pi})$ consists of closed curves in $\openannulus$.
Take two loops $\alpha$ and $\alpha'$ in $ \Phi^{-1}(
\immclosedann\cap \Pi_{\theta_1})$ and in $ \Phi^{-1}(
\immclosedann\cap \Pi_{\theta_1+\pi})$ respectively.
Let~$d_1$ be the signed distance function from the plane $\overline\Pi_{\theta_1}\cup\Pi_{\theta_1+\pi}$, positive on $\partial_2^D\Om$. Since $d_1\circ\Phi$ changes its sign when one crosses
transversally $\alpha$ and $\alpha'$, we~easily see that both $\alpha$ and $\alpha'$ cannot be homotopically trivial in $\openannulus$ (by harmonicity of $d_1\circ\Phi$, if for instance $\alpha $ is homotopically trivial in $\openannulus$, by the maximum principle $d_1\circ\Phi=0$ in the region enclosed by $\alpha$, \ie the image of $\Phi$ is locally flat, contradicting the analyticity of $\Phi$). Hence, since $\Phi$ is an embedding, they run exactly one time around $C_1$; as a consequence, they must be homotopically equivalent to each other in $\openannulus$. On the other hand, they do not intersect each other
($\Phi$ is an embedding), so they bound an annulus-type region in $\openannulus$, and by harmonicity $d_1\circ\Phi$ is constantly null in this region.
This would imply again that the image by $\Phi$ of this annulus is contained in $\overline \Pi_{\theta_1}\cup\Pi_{\theta_1+\pi}$, a contradiction.

In case \ref{case:new2}, from our assertion, we
deduce that $\Pi_{\theta_2+\pi}$ might intersect either $\partial^0\Om$ or $\partial^D_1\Om$.
Further we can exclude that $\Pi_{\theta_2+\pi}$ intersects $\partial^0\Om$ (otherwise, we~repeat the argument for case \ref{case:new1} switching the role of $\theta_1$ and $\theta_2$).
Therefore the only remaining possibility is that $\Pi_{\theta_2+\pi}$ intersects $\partial_1^D\Om$ (see Figure \ref{fig-costruzione}).
Let $d_2$ be the signed distance function from the plane $\overline{\Pi_{\theta_2}}\cup\Pi_{\theta_2+\pi}$ positive on $\partial_2^D\Om$. In~particular, $d_i\circ\Phi$, $i=1,2$, is positive on the circle $C_2$ of
$\overline\Sigma_\mathrm{ann}$. By hypothesis on $d_i$, $i=1,2$, we~see that $d_1$ is positive on $\Pi_{\theta_2}$, and $d_2$ is positive on $\Pi_{\theta_1}$.

As in case \ref{case:new1}, let $\alpha\subseteq \Phi^{-1}(\immclosedann\cap \Pi_{\theta_1})$ and
$\beta\subseteq \Phi^{-1}(\immclosedann\cap \Pi_{\theta_2})$
be two loops. We~know that $\alpha$ and~$\beta$ are closed in $\openannulus$.
Again, we~conclude that $\alpha$ and~$\beta$ are homotopically equivalent in $\openannulus$, and both run one time around $C_1$. Assume without loss of generality that $\beta$ encloses $\alpha$, which in turn encloses $C_1$. Since $d_2\circ\Phi$ is positive on both $\alpha$ and $C_2$, $d_2\circ\Phi$ must be positive in the region enclosed between them, contradicting the fact that it
vanishes on $\beta$.

If instead we are in case \ref{case:new3} we can argue as in case \ref{case:new2} and get a contradiction. In~all cases \ref{case:new1}, \ref{case:new2}, and \ref{case:new3}, we~reach a contradiction which derives by assuming that $\pi(\immclosedann)$ is not simply connected. The proof is achieved.
\end{proof}

We next proceed to characterize the geometry of $\Om\cap\partial
\pi(\immclosedann)$.

\begin{lemma}[Trace on the horizontal plane]
\label{lem:trace_on_the_horizontal_plane}
Suppose $m_2(\Gamma)<m_1(\Gamma_1)+m_1(\Gamma_2)$ and let
$\Phi\in\mathcal{P}_2(\Gamma)$ be a $\mathcal{MY}$ solution to \eqref{catenoid_plateau}.
Then $\Om\cap\partial \pi(\immclosedann)$
consists of two disjoint Lipschitz
embedded curves $\beta_1$ and $\beta_2$ joining $q_1$ to $p_2$, and $q_2$ to $p_1$, respectively.
Moreover, the closed regions $E_i$
enclosed between $\partial^0_1\Om$ and $\beta_i$
are convex for $i=1,2$.
\end{lemma}

\begin{proof}
By Lemma \ref{lem:simply_connectedness}, $\pi(\immclosedann)$ is
simply connected in $\overline\Omega$, and $\pi(\immclosedann)\cap \partial\Om=\partial^D\Om$.
Therefore $\overline\Omega\setminus \pi(\immclosedann)$ consists of two simply connected components, one containing $\partial_1^0\Omega$ and the other containing $\partial_2^0\Om$.
Let $E_1$ and $E_2$ be the closures of these two components,\footnote{The sets $E_1$ and $E_2$ have nonempty interior, since $\Phi(\openannulus)$ is contained in the interior of the convex hull of $\Phi(\partial\openannulus)$, hence contained in the cylinder $\Omega\times\R$.}
so that in particular the boundary of $E_i$ is a
simple Jordan curve
of the form $ \beta_i\cup\partial_i^0\Om$ for some
embedded curve $\beta_i\subset\overline\Om$ joining the endpoints of $\partial_i^0\Om$. We~will prove that $E_i$ is convex for $i=1,2$.
This will also imply that $\beta_i$ are Lipschitz.

Take $i=1$, and assume by contradiction that $E_1$ is not convex.
Thus we can find a line $l$ in $\R^2$ and three different points $A_1$, $A_2$, $A_3$ on $l$, with $A_2\in\overline{A_1A_3}$, so that $A_2$ is contained in $\Om\setminus E_1$, and $A_1$ and $A_3$ belong to the interior of $E_1$.

Consider the region $\pi(\immclosedann)\setminus l$, which consists in several (open) connected components. There is one of
these connected components, say $U$,
which does not intersect $\partial^D\Om$ and
whose boundary contains $A_2$. In addition,
$\overline U \cap \partial^D\Om = \emptyset$.
Indeed, $\partial U$ is the union of a segment $L$ (containing $A_2$) and a curve $\gamma$ (contained in $\beta_1\subseteq\partial (\pi(\immclosedann)$) joining its endpoints.
Hence, $\overline U\setminus U=\gamma\cup L$, and $L$ cannot intersect $\partial^D\Om$ by the hypothesis on $A_1$, $A_2$, and $A_3$.

Let $\Pi_l\subset\R^3$ be the
plane containing $l$ and
orthogonal to the plane containing $\Om$.
As usual, $\Pi_l\cap
\immclosedann$ is a family of closed curves, possibly with endpoints on $\Gamma\cap \Pi_l$.
Now, pick a point $P$ on $\partial U\setminus L$, and let $Q$ be a point on $\immclosedann$ so that $\pi(Q)=P$.
Let $d_l:\R^3\to \R$ be the signed distance from $\Pi_l$,
with $d_l(Q)=d_l(P)>0$. We~claim that, if $D$ is the connected component of $\{w\in
\overline\Sigma_\mathrm{ann}:d_l\circ\Phi(w)>0\}$ containing the point $\Phi^{-1}(Q)$, then $D\cap \partial\openannulus=\emptyset$.
This would contradict the harmonicity of $d_l\circ\Phi$, since $d_l\circ\Phi$ would be zero
on $D$, but $d_l(Q)>0$, in contrast with the maximum principle.

Assume by contradiction that the converse holds. Then there is an arc $\alpha:[0,1]\to D\cup\partial\openannulus$ joining $\Phi^{-1}(Q)$ to $\partial\openannulus$ .
The image of the
map $\pi\circ\Phi\circ \alpha$ is an arc in $\overline \Omega$
joining $P$ to $\partial^D\Om$ and such that $d_l\geq0$ on it.
Clearly this arc is a subset of $\pi(\immclosedann)$. Since $\pi\circ\Phi\circ \alpha(0)=P$, it follows that the image of $\pi\circ\Phi\circ \alpha$ is contained~in~$\overline U$. Now,~$\overline U$~does not intersect $\partial^D\Om$, contradicting that $\pi\circ\Phi\circ \alpha(1)\in \partial^D\Om$. This concludes the proof.
\end{proof}

In the next step we show
that there exists a set $E\subset\R^3$ of finite perimeter such that
\[
\partial E=\partial^*E=\Phi(\openannulus)\cup\overline\Delta_1\cup\overline\Delta_2,
\]
where $\partial^*$ denotes the reduced boundary, and
\begin{multline}\label{delta}
\Delta_i:=\{P=(P',P_3)\in \R^3:P'=(P_1,P_2)\in \partial_i^D\Om, \;P_3\in (-\varphi(P'),\varphi(P'))\},
\\
i=1,2.
\end{multline}
In~particular, $\overline\Delta_1\cup\overline\Delta_2\subset(\partial\Om)\times\R$ and
$(\Om\times\R)\cap\partial E=\Phi(\openannulus)$.

We first fix some notation. We~let $\jump{E}\in\mathcal{D}_3(\R^3)$ be the 3-current given by integration over $E$ with $E\subset\R^3$ a set of finite perimeter.
To every $\mathcal{MY}$ solution
$\Phi\in\mathcal{P}_2(\Gamma)$
to~\eqref{catenoid_plateau} we associate the push-forward 2-current $\Phi_\sharp\jump{\openannulus}\in\mathcal D_2(\R^3)$ given by integration over the (suitably oriented) surface $\Phi(\openannulus)$
\cite[\S 7.4.2]{Krantz-Parks}.
Finally, if $\mathcal{T}\in\mathcal{D}_k(U)$ with $U\subset\R^3$ open and $k=2,3$, we~denote by $|\mathcal{T}|$ the mass of $\mathcal{T}$ in $U$
\cite[p.\,358]{Federer}.

\begin{lemma}[Region enclosed by $\Phi(\openannulus)$]
\label{lem:region_enclosed_by_Phi(Sigma_min)}
Suppose $m_2(\Gamma)<m_1(\Gamma_1)+m_1(\Gamma_2)$ and let $\Phi\in\mathcal{P}_2(\Gamma)$ be a $\mathcal{MY}$ solution to \eqref{catenoid_plateau}.
Then there is a closed
finite perimeter set $E\subset\overline\Om\times\R$ such that $(\Om\times\R)\cap\partial E=\Phi(\openannulus)$.
\end{lemma}

\begin{proof}
As $\Phi_\sharp\jump{\openannulus}$ is a boundaryless integral $2$-current in $\Omega\times\R$, there exists (see, e.g., \cite[Th.\,7.9.1]{Krantz-Parks}) an integral $3$-current $\mathcal E\in \mathcal D_3(\Omega\times\R)$ with $\partial\mathcal E=\Phi_\sharp\jump{\openannulus}$, and we might also assume that the support of $\mathcal E$ is compact in $\overline\Omega\times\R$. We~claim that, up to switching the orientation of $\Phi_\sharp\jump{\openannulus}$, $\mathcal E$ has multiplicity in $\{0,1\}$, and hence is the integration $\jump{E}$ over a bounded measurable set $E$.
Since $\partial\mathcal E=\Phi_\sharp\jump{\openannulus}$, this will be a
finite perimeter set, and $\jump{(\Om\times\R)\cap\partial^*E}=\Phi_\sharp\jump{\openannulus}$.

By Federer decomposition theorem \cite[\S 4.2.25, p.\,420]{Federer} (see also \cite[\S 4.5.9]{Federer} and \cite[Th.\,7.5.5]{Krantz-Parks}) there is a sequence
$(E_k)_{k\in \mathbb N}$ of finite perimeter subsets of $\Om\times\R$ such that
\begin{equation}\label{6.9a}
\mathcal E=\sum_{k=1}^{+\infty}\sigma_k\jump{E_k},\qquad \sigma_k\in \{-1,1\},
\end{equation}
and
\begin{align}\label{6.9}
|\mathcal E|=\sum_{k=1}^{+\infty}|E_k|\quad\text{and} \quad|\partial\mathcal E|=\mathcal H^2(\Phi(\openannulus))=\sum_{k=1}^{+\infty}\mathcal H^2(\partial^*E_k).
\end{align}
We start by observing that
\begin{equation}\label{inclusione}
\partial^*E_k\subseteq \Phi(\openannulus)\qquad \forall k\in \mathbb N.
\end{equation}
Indeed, fixing $k\in \mathbb N$, by the second equation in \eqref{6.9}, we~have that $\partial^*E_k$ is contained in the support of $\partial\mathcal E$, which in turn is $\Phi(\openannulus)$.
As a consequence, if $P=(P_1,P_2,P_3)\in(\Om\times\R)\cap \overline{\partial^*E_k}$, then $P\in \Phi(\openannulus)$. Around $P$ we can find suitable coordinates and a cube $U=(P_1-\eps,P_1+\eps)\times (P_2-\eps,P_2+\eps)\times(P_3-\eps,P_3+\eps)$ such that $\Phi(\openannulus)\cap U$ is the graph $\mathcal G_h$ of a smooth function $h:(P_1-\eps,P_1+\eps)\times (P_2-\eps,P_2+\eps)\to (P_3-\eps,P_3+\eps)$. Moreover, $\Phi_\sharp\jump{\openannulus}=\jump{\mathcal G_h}$ in $U$.

We claim that
\[
\forall k\quad\text{either}\quad E_k\cap U=U \cap SG_h
\quad\text{or} \quad E_k\cap U=U\setminus SG_h.
\]
Indeed, assume for instance that $|E_k\cap U \cap
SG_h|>0$ and $|(SG_h\setminus E_k)\cap U|>0$; by the constancy lemma \cite{Krantz-Parks} it follows that $\partial\jump{E_k}$ is nonzero in the simply connected open set $SG_h$, contradicting \eqref{inclusione}.
As a consequence of the preceding claim, we~have $U\cap \partial^*E_k=U\cap \Phi(\openannulus)$. Since this argument holds for any choice of $P\in (\Omega\times\R)\cap \overline{\partial^*E_k}$, we~have proved that $(\Omega\times\R)\cap \overline{\partial^*E_k}$ is relatively open
(and relatively closed at the same time) in $\Phi(\openannulus)$, which in turn being a connected open set, implies
\[
\immclosedann=\overline{\partial^*E_k}\qquad \forall k\in \mathbb N.
\]

Denote by $\mathcal I^{\pm}:=\{k\in \mathbb N:\sigma_k=\pm1\}$, with
$\sigma_k$ as in \eqref{6.9a}.
Going back to the local behaviour around $P\in \Phi(\openannulus)$, if $U$ is a neighbourhood as above, we~see that for all $k\in \mathcal I^+$ either $E_k\cap U=SG_h$ or $E_k=U\setminus SG_h$ (namely, all the $E_k$'s coincide in $U$), since otherwise, there will be cancellations in the series $\sum_{k\in \mathcal I^+}\partial \jump{E_k}$, in contradiction with the second formula in \eqref{6.9}. Assume without loss of generality that for all $k\in \mathcal I^+$ we have $E_k\cap U=SG_h$; thus, arguing as before, for all $k\in \mathcal I^-$ we must have $E_k\cap U=U\setminus SG_h$.

We obtain that $\mathcal E\res U=m\jump{SG_h}-n\jump{U\setminus SG_h}$
for some nonnegative integers $n,m$. Since $(\partial \mathcal E)\res U=(m+n)\jump{\mathcal G_h}$
and also $(\partial\mathcal E)\res U=\Phi_\sharp\jump{\openannulus}=\jump{\mathcal G_h}$ in $U$, we~conclude $m+n=1$. Hence either $m=1$
and $n=0$, or $m=0$ and $n=1$. On the other hand, we~know that
$\mathcal E\res U=\sum_{k\in \mathcal I^+}\jump{E_k\cap U}-\sum_{k\in \mathcal I^-}\jump{E_k\cap U}$, from which it follows that~$\mathcal I^+$ has cardinality $m$ and $\mathcal I^-$ has cardinality $n$.
Namely, one of the sets~$\mathcal I^\pm$ is empty, and the other contains
one index only.

We conclude that the sum in \eqref{6.9a} involves one index only, that is, there is only one compact set $E$ in $\overline\Omega\times\R$ such that (up to switching the orientation)
\[
\mathcal E=\jump{E}.
\]
This concludes the proof.
\end{proof}
For later convenience, from now on we denote by $E$ the closure of a
precise representative of the set found in Lemma
\ref{lem:region_enclosed_by_Phi(Sigma_min)}.
\begin{remark}\label{rem:steiner-sym}
From the fact that $ (\overline \Omega\times\R)\cap \partial E=\immclosedann\cup\overline\Delta_1\cup\overline\Delta_2$, we
easily see that $\pi( E)=\pi(\immclosedann)$ which, by Lemma \ref{lem:simply_connectedness}, is simply connected.
\end{remark}

We denote by $\sym_\mathrm{st}(E)$ the set (symmetric with respect to the horizontal plane $\R^2\times\{0\}$) obtained applying to $E$ the Steiner symmetrization with respect to $\R^2\times\{0\}$. \\
Clearly $\sym_\mathrm{st}(E)\cap(\partial_i^D\Om\times\R)=\overline{\Delta_i}$ with $\Delta_i$ defined as in \eqref{delta}. We~define
\begin{equation}\label{surf}
S:= \partial(\sym_\mathrm{st} (E))\setminus(\Delta_1\cup\Delta_2),\quad
S^+:= S\cap\{x_3\ge0\},\quad S^-:= S\cap\{x_3\le0\}.
\end{equation}
Since $P(\sym_\mathrm{st}(E))\le P(E)$ (here $P(\cdot)$ is the perimeter in $\R^3$ \cite{AFP}) we have $\mathcal{H}^2(S)\le\mathcal{H}^2(
\immclosedann)$.

\begin{lemma}[Graphicality
of $\partial(\sym_\mathrm{st} (E))$ and continuity up to the boundary]
\label{lem:graphicality_of_boundary_and_continuity_up_to_the_boundary}
Suppose that $m_2(\Gamma)<m_1(\Gamma_1)+m_1(\Gamma_2)$ and let $\Phi\in\mathcal{P}_2(\Gamma)$ be a $\mathcal{MY}$ solution to~\eqref{catenoid_plateau}.
Let $E$ be the finite perimeter set given by Lemma \ref{lem:region_enclosed_by_Phi(Sigma_min)} and $S^\pm$ be as in \eqref{surf}.
Then there is $\widetilde \psi\in BV(\mathrm{int}(\pi(E)))\cap C^0(\pi(E))$ such that
$
S^\pm=
\mathcal {G}_{\pm\widetilde\psi}$.
In~particular, $S^\pm\cap(\R^2\times\{0\})=\overline{\Om\cap\partial(\pi(E))}$.
\end{lemma}

\begin{proof}
Since $E$ has finite perimeter,
there exists a function $\widetilde \psi\in BV(\mathrm{int}(\pi(E)))$
such that
$
S^\pm=\mathcal G_{\pm\widetilde \psi}$ \cite{CCF}. So,
we only need to show that
$\widetilde\psi$ is continuous (note that $\pi(E)$ is a closed set).
Take a point $\pointinOm$ in the interior of $\pi(E)$; if $\pointinOm
=\pi(\Phi(w))$ for some~$w$, then $w\in \openannulus$, since
$\pi(\Phi(C_i)) \subset \partial \Om$
for $i=1,2$ (recall that $C_1$ and $C_2$ form the boundary of $\openannulus$).
If at none of the points of $\pi^{-1}(\pointinOm)\cap
\immclosedann$ the tangent plane to $\immclosedann$ is vertical,
then $\widetilde\psi$ is $C^\infty$ in a neighbourhood of $\pointinOm$,
since it is the linear combination of smooth functions
(see the discussion after formula \eqref{6.18} below, where details are given).
Therefore we only have to check continuity of $\widetilde\psi$
at those points $\pointinOm$ for which there is $\pointinspace\in
\pi^{-1}(\pointinOm)\cap\immclosedann$ such that
$\immclosedann$ has a vertical tangent plane $\Pi$ at $\pointinspace$.

Consider a system of Cartesian coordinates
centered at $ \pointinspace$, with the $(x,y)$-plane
coinciding with $\Pi$, the $x$-axis coinciding with the line $\pi^{-1}( \pointinOm)$,
and let $z=z(x,y)$ (defined at least in a neighbourhood of $0$) be the analytic function whose graph coincides with $\immclosedann$.
This map, restricted to the $x$-axis, is analytic and
vanishes at $x=0$; hence it is
either identically zero or it
has a discrete set of zeroes (in the neighbourhood where it exists). We~now exclude the former case: If $z(\cdot,0)$ is identically zero,
it means that around $\pointinspace$ there is a vertical open
segment included in $\pi^{-1}(\pointinOm)$,
which is
contained in $\immclosedann$. Let $Q$ be an extremal point of this segment, and
let $\Pi_Q$ be the tangent plane to $\immclosedann$ at $Q$.
This plane must
contain as tangent vector the above segment, hence $\Pi_Q$ is vertical and
contains $\pi^{-1}(\pointinOm)$.
Choosing again a suitable Cartesian coordinate system centered at $Q$
we can express locally
the surface $\immclosedann$ as the graph of an analytic function defined in a
neighbourhood of $Q$ in $\Pi_Q$, and so
the restriction of this map to
$\pi^{-1}(\pointinOm)$ is analytic in a neighbourhood of $Q$,
hence it must be identically zero since it is zero in a
left (or right) neighbourhood of $Q$.
What we found is that we can properly extend the segment $\overline{\pointinspace Q}$
on the $Q$ side to a segment $\overline{PR}$ contained in
$\immclosedann$.
This proves that $\immclosedann\cap \pi^{-1}(\pointinOm)$ is relatively open in $\pi^{-1}(\pointinOm)$. Since it is also relatively closed, it coincides with the whole line $\pi^{-1}(\pointinOm)$, which
is impossible since ${\Phi}(\overline\Sigma_\mathrm{ann})$ is bounded.

Hence the zeroes of the function $z(\cdot,0)$ are isolated, so we have
shown:

\begin{assertion}\label{assertionA}
Let $\pointinspace\in \pi^{-1}(\pointinOm)\cap
\immclosedann$. Then in a neighbourhood of $\pointinspace$
the only intersection between $\immclosedann$
and $\pi^{-1}(\pointinOm)$ is $\pointinspace$ itself.
\end{assertion}

Now, we~can conclude the proof of
the continuity of the function $\widetilde\psi$.
Write
\[
\pi^{-1}(\pointinOm)\cap
\immclosedann=\{Q_1,Q_2,\dots,Q_{m}\}\subset\Omega\times\R.
\]
It follows that
\begin{align}\label{verticalslice}
2\widetilde\psi(\pointinOm)=\mathcal H^1(\pi^{-1}(\pointinOm)\cap E)=\sum_{j=1}^m\sigma_j (Q_j)_3,
\end{align}
where $(Q_j)_3$ is the vertical coordinate of $Q_j$ and
\begin{align}
\sigma_j=\begin{cases}
-1&\text{if }\overline{Q_{j-1}Q_j}\subset\overline{ \R^3\setminus E}\text{ and }\overline{Q_{j}Q_{j+1}}\subset E,\\
1 &\text{if }\overline{Q_{j-1}Q_j}\subset E\text{ and }\overline{Q_{j}Q_{j+1}}\subset \overline{\R^3\setminus E},\\
0&\text{otherwise,}
\end{cases}
\qquad \qquad j=1,\dots,m.
\end{align}
Let $(\pointinOm_k)
\subset \textup{int}(\pi(E))$ be a sequence converging to $\pointinOm$, and write $\pi^{-1}
(\pointinOm_k )\cap \immclosedann=\{Q^k_1,Q^k_2,\dots,Q^k_{m_k}\}\subset\Omega\times\R$. With a similar
notation as above, we~have
\begin{equation}\label{verticalslice2}
2\widetilde\psi(\pointinOm_k)=\mathcal H^1(\pi^{-1}(\pointinOm_k)
\cap E)=\sum_{j=1}^{m_k}\sigma^k_j (Q^k_j)_3.
\end{equation}
Now, if at
every point $Q_j$ the tangent plane to $\immclosedann$
is not vertical, then $\immclosedann$ is a smooth Cartesian surface in a neighbourhood of
$Q_j$, and so it is clear that, for $k$ large enough,
\begin{align}\label{6.18}
m=m_k,\qquad
Q_j^k\to Q_j,\qquad \sigma^k_j\to\sigma_j\qquad\text{for all }j=1,\dots,m,
\end{align}
and the continuity of $\widetilde \psi$ at $P'$ follows.
Therefore it remains to check continuity in the case that the tangent plane to some $Q_j$ is vertical.

Let $\widetilde Q$
be one of these points, with associated sign $\widetilde\sigma \in \{0,1\}$.
By assertion \ref{assertionA} there is $\delta>0$
so that $\widetilde Q$ is the unique intersection between $\pi^{-1}(\pointinOm)$ and
$\immclosedann$ with vertical coordinate in $[\widetilde Q_3-\delta,\widetilde Q_3+\delta]$.
This means that the segments
\[
\pi^{-1}(\pointinOm)\cap\{\widetilde Q_3-\delta<x_3<\widetilde Q_3\}\quad\text{and}\quad \pi^{-1}(\pointinOm)\cap\{\widetilde Q_3<x_3<\widetilde Q_3+\delta\}
\]
are contained in either $\textup{int}(E)$ or
$\R^3\setminus E$. In~particular,
there is a neighbourhood $U\subset\Omega$ of $\pointinOm$ such that $U\times \{x_3=\widetilde Q_3-\delta\}$ and $U\times \{x_3=\widetilde Q_3+\delta\}$ are
subsets of $\textup{int}(E)$ or of $\R^3\setminus E$. Suppose
without loss of generality that both are inside $\R^3\setminus E$ (the other cases being
similar),
so that $\widetilde \sigma=0$. We~infer that, for $k$ large enough so that $\pointinOm_k\in U$,
there is a finite subfamily $\{Q^k_j:j\in J\}$ of $\{Q^k_1,Q^k_2,\dots,Q^k_{m_k}\}$ contained in $\{ \widetilde Q_3<x_3<\widetilde Q_3+\delta\}$ and which satisfies the following: The sum in \eqref{verticalslice2} restricted to such subfamily reads as:
\[
 \sum_{j\in J}\sigma^k_j (Q^k_j)_3=(Q^k_{j_l})_3-(Q^k_{j_{l-1}})_3+\dots+(Q^k_{j_2})_3-(Q^k_{j_1})_3,
\]
where
$J=\{j_1,j_2,\dots,j_l : j_1<j_2<\dots<j_l\}$ and $(Q^k_{j_l})_3>(Q^k_{j_{l-1}})_3>\dots>(Q^k_{j_2})_3>(Q^k_{j_1})_3$ (if $j_l=1$ necessarily $\sigma_{j_1}^k=0$ and the sum is zero). We~have to show that this sum tends to $\widetilde\sigma \widetilde Q_3=0$ as $k\to
+\infty$, which is true, since each $Q^k_j$ tends to $\widetilde Q$.
Repeating this argument for each point $\widetilde Q$
with a vertical tangent plane to $\immclosedann$,
the proof of continuity of $\widetilde\psi$ in the interior of $\pi(E)$ follows.

Now, let
$\pointinOm\in\partial(\pi(E))$. If $\pointinOm\in \Om
\cap \partial(\pi(E))$
then every point in $\pi^{-1}(\pointinOm)\cap \immclosedann$ has vertical tangent plane
and we can argue as in the previous case.
It remains to show continuity of $\widetilde \psi$
on $\partial\pi(E)\cap \partial\Omega$. In this case we
exploit the fact that the interior of
$\immclosedann$ is contained in $\Om\times \R$. We~sketch the proof without details since it is very similar to the previous argument. Let
$\pointinOm\in \partial_1^D\Om$,
thus $\pi^{-1}(\pointinOm)\cap\Gamma_1$ consists of two distinct points
$Q_1$ and $Q_2$. Let $(\pointinOm_k)$
be a sequence of points in ${\pi(E)}$
converging to~$P$.
For $\pointinOm_k \in \partial^D_1\Om$ it follows $\pi^{-1}(\pointinOm_k)
\cap\Gamma_1=\{Q^k_1,Q^k_2\}$ and the continuity of $\widetilde\psi$
follows from the continuity of $\varphi$ on $\partial^D_1\Om$,
whereas if $\pointinOm_k$
is in the interior of $\pi(E)$ there holds $\pi^{-1}(\pointinOm_k)
\cap\Gamma_1=\{Q^k_1,Q^k_2,\dots, Q^k_{m_k}\}$.
Using the continuity of $\Phi$ up to $C_1$, it is easily seen that all such points must converge,
as $k\to +\infty$, either to $Q_1$ or to $Q_2$.
Hence we can repeat an argument similar to the one used before.
\end{proof}

\begin{lemma}\label{lem:construction-of-parametrization}
Suppose $m_2(\Gamma)<m_1(\Gamma_1)+m_1(\Gamma_2)$ and let $\Phi\in\mathcal{P}_2(\Gamma)$ be a $\mathcal{MY}$ solution to \eqref{catenoid_plateau}. Let $E$ be the finite perimeter set given in Lemma \ref{lem:region_enclosed_by_Phi(Sigma_min)} and let
$S$ be defined as in \eqref{surf}.
Then there is an injective map $\widetilde\Phi\in H^1
(\openannulus;\R^3)\cap C^0(\overline \Sigma_\mathrm{ann};\R^3)$
which maps $\partial \openannulus$
weakly monotonically to $\Gamma$ and such that $\widetilde\Phi (\overline \Sigma_\mathrm{ann}) =S$, and
furthermore
\begin{align}\label{eq:area-symmetrization}
\mathcal H^2(S)= \int_{\Sigma_{\mathrm{ann}}}|\partial_{w_1}\widetilde \Phi\wedge \partial_{w_2}\widetilde \Phi |dw=
\int_{\Sigma_{\mathrm{ann}}}|\partial_{w_1}\Phi\wedge\partial_{w_2}\Phi|dw=m_2(\Gamma).
\end{align}
In~particular, $\widetilde \Phi$ is a solution of \eqref{catenoid_plateau}.
\end{lemma}

\begin{proof}
By Lemma \ref{lem:graphicality_of_boundary_and_continuity_up_to_the_boundary}
there is $\widetilde \psi\in BV(\mathrm{int}(\pi(E)))\cap C^0(\pi(E))$ such that $S^\pm=
\mathcal {G}_{\pm\widetilde\psi}$.
As a consequence, for $p\in \partial^D\Om$
we have $\widetilde\psi(p)=\varphi(p)$ and for $p\in\Omega\cap \partial (\pi(E))$
we have $\widetilde\psi(p)=0$.

By Lemma \ref{lem:simply_connectedness} $\pi(E)$ is simply connected, and so
the maps $\widetilde \Psi^\pm:\pi(E)\to \R^3$ given by
$\widetilde \Psi^\pm(p):=(p,\pm\widetilde\psi(p))$ are disk-type parametrizations of $S^\pm$.
Moreover, $S^+$ and $S^-$ glue to each other along
$(\R^2\times\{0\})\cap\partial(\sym_\mathrm{st}(E))=
\beta_1\cup\beta_2$, where $\beta_1$ and $\beta_2$
are the curves given by Lemma \ref{lem:trace_on_the_horizontal_plane} .

Let $(\sigma,\psi)\in\admclassconv$ be a minimizer of $\mathcal F$ which satisfies
properties \ref{1.}--\ref{5.} of Theorem \ref{teo:structure_of_minimizers}. Setting $\widetilde\sigma:=(\beta_1,\beta_2)$ and
extending $\widetilde\psi$ to zero in $\overline\Om\setminus \pi(E)$
(still calling $\widetilde \psi$ such an extension),
by minimality we get
\[
2\mathcal F(\sigma,\psi)\leq 2\mathcal F(\widetilde\sigma,\widetilde\psi)=\mathcal H^2(S),
\]
whence
\begin{align}\label{last_eq}
2\mathcal F(\sigma,\psi)\leq \mathcal{H}^2(S)\le\mathcal{H}^2(\immclosedann)
=\int_{\Sigma_{\mathrm{ann}}}|\partial_{w_1}\Phi\wedge\partial_{w_2}\Phi|dw=m_2(\Gamma).
\end{align}
We are in the hypotheses of Lemma \ref{lem_tec2}\ref{lem_tec2a},
therefore there exists a map $\widehat\Phi\in P_2(\Gamma)$ parametrizing $\mathcal G_{\psi\res(\overline{\Om\setminus E(\sigma)})}\cup \mathcal G_{-\psi\res(\overline{\Om\setminus E(\sigma)})}$ which is a minimizer
of \eqref{catenoid_plateau}. In~particular,
$2\mathcal F(\sigma,\psi)=m_2(\Gamma)$, and all inequalities in \eqref{last_eq} are equalities. We
deduce also that
$(\widetilde \sigma,\widetilde\psi)$ is a minimizer of $\mathcal F$ in $\admclassconv$, so that by Theorem \ref{teo:structure_of_minimizers} $\widetilde \psi$ is analytic in $\textup{int}(\pi(E))$. As a consequence it belongs to $W^{1,1}(\textup{int}(\pi(E));\R^3)$.
Applying Lemma \ref{lem_tec1}\ref{lem_tec1a} and Lemma \ref{lem_tec2}\ref{lem_tec2a}, we~get the existence of $\widetilde \Phi\in P_2(\Gamma)$ as in the statement, and we have concluded.
\end{proof}

\begin{lemma}\label{lemma_step5}
Suppose $m_2(\Gamma)<m_1(\Gamma_1)+m_1(\Gamma_2)$ and let $\Phi\in\mathcal{P}_2(\Gamma)$ be a $\mathcal{MY}$ solution to \eqref{catenoid_plateau}. Let $E$ be the finite perimeter set given in Lemma \ref{lem:region_enclosed_by_Phi(Sigma_min)} and let $S$ be defined as in \eqref{surf}. Then
$\immclosedann=S$
and in particular
\[
E=\sym_\mathrm{st}(E).
\]
\end{lemma}

\begin{proof} By Lemma \ref{lem:construction-of-parametrization} we have $\mathcal H^2(S)=m_2(\Gamma)$, from which it follows that
$P(\sym_\mathrm{st}(E))= P(E)$. Then we can apply \cite[Th.\,1.1]{CCF}
to deduce the existence of two functions
$f,g:\pi(E)\to \R$ of bounded variation, such that $\partial^* E=\mathcal G_f\cup\mathcal G_{g}$ (up~to $\mathcal H^2$-negligible sets). We~will show that $f=\widetilde\psi$ and $g=-\widetilde\psi$. To this aim,
thanks again to \cite[Th.\,1.1]{CCF}, we~know that for a.e. $p\in \pi(E)$, the two unit (external to $E$) normal
vectors $\nu^f=(\nu^f_1,\nu^f_2,\nu^f_3)$ and $\nu_g=(\nu^g_1,\nu^g_2,\nu^g_3)$ to $\mathcal G_f$ and $\mathcal G_g$ at the points $(p,f(p))$ and $(p,g(p))$, respectively, satisfy\enlargethispage{2pt}
\begin{align}\label{normals}
(\nu^f_1,\nu^f_2,\nu^f_3)=(\nu^g_1,\nu^g_2,-\nu^g_3).
\end{align}
To conclude the proof it is then sufficient to show that $f=-g$ a.e. on $\pi(E)$: indeed this would readily imply $E=\sym_\mathrm{st}(E)$ and hence $f=\widetilde\psi$.
Let $p\in \textup{int}(\pi(E))$; if
\begin{equation}
\pi^{-1}(p)\cap S=\{P_1,P_2,\dots,P_k\},
\end{equation}
then for a.e. $p\in \textup{int}(\pi(E))$ it is $k\leq2$.
Now we show that, for all $p\in \textup{int}(\pi(E))$, if $k>1$, none of the points $\{P_1,P_2,\dots,P_k\}$ has vertical tangent plane.
Assume by contradiction that $P_1$ has vertical tangent plane $\Pi_1$. In this case $\Pi_1\cap S$ consists, in a neighbourhood $U$ of $P_1$, of at least $2$ curves crossing {transversally} (see \cite[\S 373]{Nitsche:89}) at $P_1$. These curves, by assertion A in the proof of Lemma \ref{lem:graphicality_of_boundary_and_continuity_up_to_the_boundary} , intersect $\pi^{-1}(p)$ only at $P_1$. Moreover, in a neighbourhood $V$ of $P_2$, with $U\cap V=\emptyset$, $\Pi_1\cap S$ consists of (at least) one
curve passing through $P_2$. This curve is locally Cartesian if $\pi^{-1}(p)$ crosses $S$ transversally in $P_2$, otherwise it is
locally the union of two curves ending at $P_2$, with vertical tangent plane, which lie on the same side of $\Pi_1$ with respect to $\pi^{-1}(p)$. In both cases, we~deduce that there is a point $q\in \Pi_1\cap (\Om\times\{0\})$ for which $\pi^{-1}(q)$ intersects transversally $S$ in at least three points. As a consequence, for all $q'$ in a neighbourhood of $q$ in
$\Om$, the line $\pi^{-1}(q')$ intersects $S$ at more than two points, which is a contradiction. We~have proved the following:

\begin{assertion*}
For all $p\in\textup{int}(\pi(E))$ the line $\pi^{-1}(p)$
intersects $S$ either
transversally at two points $P_1,P_2$, or at only one point $P_1$.
\end{assertion*}

Now we see that the latter case cannot happen.
Indeed, first one checks
that in this case the intersection cannot be transversal,\footnote{This is a consequence of the fact that the line $\pi^{-1}(p)$ must lie outside the set $E$, with the only exception of the point $P_1$. Indeed, otherwise, there must be some other point in $\pi^{-1}(p)\cap S$, $E$~being compact in $\R^3$.} and that $\pi^{-1}(p)$ must be tangent to $S$ at $P_1$. Let $\Pi_1$ be the
vertical tangent plane to $S$ at $P_1$.
Let $\Pi_1^\perp$ be the vertical plane orthogonal to $\Pi_1$
passing
through $P_1$. In a neighbourhood of $P_1$, the unique curve in $S\cap \Pi_1^\perp$
must be the union of two curves joining at $P_1$, and these curves must belong to
the same half-plane of $\Pi_1^\perp$ with boundary $\pi^{-1}(p)$.
As a consequence, if $p'\in \Omega\cap \Pi_1^\perp$ is in
that half-plane, then $\pi^{-1}(p')$ consists of at least two points; if $p'$ lies in the opposite half-plane, then $\pi^{-1}(p')$ is empty. This means that necessarily $p\in \partial \pi(E)$. Namely, the previous assertion can be strengthened to:

\textit{For all $p\in\textup{int}(E)$ the line $\pi^{-1}(p)$ intersects $S$ transversally at exactly two points $P_1,P_2$.}

The consequence of this is that $f$ and $g$ belong to $W^{1,1}(\textup{int}(\pi(E)))$ and are also smooth in $\textup{int}(\pi(E))$.
Indeed, let $p\in \textup{int}(\pi(E))$, so $f(p)\neq g(p)$, and
\begin{equation}\label{twopoints}
\pi^{-1}(p)\cap S=\{(p,f(p)),(p,g(p))\}.
\end{equation}
Since $S$ is locally the graph of smooth functions around $(p,f(p))$ and $(p,g(p))$,
these functions coincide with $f$ and $g$, respectively. We~can now conclude the proof
of the lemma: let us choose a simple curve $\alpha:[0,1]\to \pi(E)$ with $\alpha(0)\in \partial^D\Om$ and $\alpha(1)=p$ such that \eqref{normals} holds for $\mathcal H^1$ a.e. $p\in \alpha([0,1])$. Since $f\circ \alpha$ and $g\circ \alpha$
are differentiable in $[0,1]$, condition
\eqref{normals} uniquely determines the tangent planes to $\mathcal G_f$ and $\mathcal G_g$, and hence it implies that the derivatives of $f\circ \alpha$ and $g\circ\alpha$ satisfy
\begin{align}
(f\circ \alpha)'(t)+(g\circ \alpha)'(t)=0,\qquad \text{ for a.e. }t\in[0,1].
\end{align}
By continuity of $f$ and $g$ one infers $f\circ\alpha+g\circ\alpha=c$ a.e. on $[0,1]$ (actually everywhere since $f+g$ is continuous), for some constant $c\in\R$. To show that $c=0$ it is sufficient to observe that $f\circ\alpha(0)=\varphi(\alpha(0))$ and $g\circ\alpha(0)=-\varphi(\alpha(0))$. Hence $f(p)=-g(p)$, and the thesis of Lemma \ref{lemma_step5} is achieved.
\end{proof}

We are now in a position to conclude the proof of Theorem \ref{crucial_teo}.

\begin{proof}[Proof of Theorem \ref{crucial_teo}]
Property \ref{crucial1} follows by Lemma \ref{lem:simply_connectedness} and Lemma \ref{lem:trace_on_the_horizontal_plane}.
Properties \ref{crucial2}--\ref{crucial4} follow by Lemma
\ref{lem:graphicality_of_boundary_and_continuity_up_to_the_boundary} and Lemma \ref{lemma_step5}. To see that $\beta_i$ are
$C^\infty$ it is sufficient to observe that, since $S^+$ and $S^-$ are Cartesian surfaces, their intersection
coincides
with the set $S\cap\{x_3=0\}$ which, by standard arguments,
is the image of the zero-set of $\Phi_3$, which is smooth.
\end{proof}
\begin{theorem}\label{teorema-finale}
{ Assume $n=2$ and $\Gamma_j$ not planar for $j=1,2$.}
Then
\begin{align}
2\min_{(s,\zeta)\in \admclassconv}\mathcal F(s,\zeta)= m_2(\Gamma).
\end{align}
\end{theorem}

\skpt
\begin{proof}\vspace*{-.5\baselineskip}
\subsubsection*{Step 1:
$2\min_{(s,\zeta)\in \admclassconv}\mathcal F(s,\zeta)\le m_2(\Gamma)$}

Suppose first $m_2(\Gamma)<m_1(\Gamma_1)+m_1(\Gamma_2)$.
Let $\Phi\in\mathcal{P}_2(\Gamma)$ be a $\mathcal{MY}$ solution to \eqref{catenoid_plateau} and let $S:=\immclosedann$. By Theorem \ref{crucial_teo} the following properties hold:
\begin{itemize}
\item $S\cap(\R^2\times\{0\})=\beta_1\cup\beta_2$ with $\beta_1$ and $\beta_2$ disjoint
embedded analytic curves joining $q_1$ to $p_2$ and $q_2$ to $p_1$, respectively;
\item $S$ is symmetric with respect to $\R^2\times\{0\}$;
\item for $i=1,2$ the closed
region $E_i$ enclosed between $\partial^0_i\Om$ and $\beta_i$ is convex;
\item $S^+=S\cap\{x_3\ge0\}$ is the graph of $\widetilde\psi\in W^{1,1}(U)\cap C^0(\overline U)$, where $U=\Om\setminus( E_1\cup E_2)$ is the open
region enclosed between $\partial^D\Om$ and $\beta_1\cup\beta_2$.
\end{itemize}
Let $(\sigma,\psi)\in\admclassconv$ be given by
\begin{equation*}
\sigma:=(\sigma_1,\sigma_2)\quad\text{and}\quad\psi:=\begin{cases}
0&\text{in } \Om\setminus U,
\\
\widetilde{\psi}&\text{in } U,
\end{cases}
\end{equation*}
where $\sigma_i([0,1])=\beta_i$ for $i=1,2$.
Then clearly $S^+=\mathcal{G}_{\psi\res(\overline{\Om\setminus E(\sigma)})}$ and
\begin{equation*}
\min_{(s,\zeta)\in\admclassconv}\mathcal{F}(s,\zeta)\le \mathcal{F}(\sigma,\psi)=\mathcal{H}^2(S^+)=\frac12m_2(\Gamma).
\end{equation*}

Now, suppose $m_2(\Gamma)=m_1(\Gamma_1)+m_1(\Gamma_2)$.
For $j=1,2$,
let $\Phi_j\in \mathcal{P}_1(\Gamma_j)$ be a solution to \eqref{plateau_1}
and $S_j:=\Phi_j(\overline B_1)$.
Let $D_j$ be the closed convex hull of $\Gamma_j$: clearly $D_1\cap D_2=\emptyset$. By Lemma \ref{lem:plateau_solution} (with $F=\overline\Om$) each $S_j$ satisfies the following properties:
\begin{itemize}
\item $S_j\cap(\R^2\times\{0\})=\beta_j\subset D_j$ is a simple {analytic} curve joining $p_j$ to $q_j$;
\item $S_j$ is symmetric with respect to $\R^2\times\{0\}$;
\item $S_j^+:=S\cap\{x_3\ge0\}$ is the graph of a function $\widetilde \psi_j\in W^{1,1}(U_j)\cap C^0(\overline U_j)$, where $U_j\subset D_j$ is the open region enclosed between $\partial^D_j\Om$ and $\beta_j$;
\item
$\Om\setminus U_j$ is convex.
\end{itemize}
Let $(\sigma,\psi)\in\admclassconv$ be given by
\begin{equation*}
\sigma:=(\sigma_1,\sigma_2)\quad\text{and}\quad\psi:=\begin{cases}
0&\text{in } \Om\setminus (U_1\cup U_2),\\
\widetilde{\psi}_j&\text{in } U_j\, \text{ for }j=1,2,
\end{cases}
\end{equation*}
where $\sigma_1([0,1]):=\overline{p_1q_2}$
and $\sigma_2([0,1]):=\beta_2\cup\,\overline{q_2p_1}\,\cup\beta_1$.
Then $S^+:=S^+_1\cup S^+_2=\mathcal{G}_{\psi\res(\overline{\Om\setminus E(\sigma)})}$ and
\begin{equation*}
\min_{(s,\zeta)\in\admclassconv}\mathcal{F}(s,\zeta)\le \mathcal{F}(\sigma,\psi)=\mathcal{H}^2(S^+)=\frac12(m_1(\Gamma_1)+m_1(\Gamma_2))=\frac12m_2(\Gamma),
\end{equation*}
and the proof of step 1 is concluded.

\subsubsection*{Step 2:
$2\min_{(s,\zeta)\in \admclassconv}\mathcal F(s,\zeta)\ge m_2(\Gamma)$}

Let $(\sigma,\psi)\in\admclassconv$ be a minimizer satisfying
properties
\ref{1.}--\ref{5.} of
Theorem \ref{teo:structure_of_minimizers}.
If $E(\sigma_1)\cup E(\sigma_2)=\emptyset$, by step 1 we can apply Lemma \ref{lem_tec2} and find an injective parametrization
$\Phi\in \mathcal P_2(\Gamma)$ such that
$\Phi_i(\partial \openannulus)=\Gamma$ weakly monotonically,
$\Phi(\overline\Sigma_\mathrm{ann})=\mathcal G_\psi\cup \mathcal G_{-\psi}$, and
\[
2\mathcal{F}(\sigma,\psi)=\int_{\Sigma_\mathrm{ann}}|\partial_{w_1}\Phi\wedge\partial_{w_2}\Phi|dw\geq m_2(\Gamma).
\]
If instead $E(\sigma_1)\cup E(\sigma_2)\ne\emptyset$, similarly we find injective parametrizations
$\Phi_1\in \mathcal P_1(\Gamma_1)$ and $\Phi_2\in \mathcal P_1(\Gamma_2) $
such that $\Phi_j(\partial B_1)=\Gamma_j$ weakly monotonically
for $j=1,2$, $\Phi_1(\overline B_1)\cup\Phi_2(\overline B_1)=\mathcal G_\psi\cup \mathcal G_{-\psi}$, and
\begin{align*}
2\mathcal{F}(\sigma,\psi)&=\int_{B_{1}}|\partial_{w_1}\Phi_1\wedge\partial_{w_2}\Phi_1|dw+\int_{B_{1}}|\partial_{w_1}\Phi_2\wedge\partial_{w_2}\Phi_2|dw\\&\geq m_1(\Gamma_1)+ m_1(\Gamma_2)\ge m_2(\Gamma).
\end{align*}
This concludes the proof.
\end{proof}

Now the proof of Theorem \ref{teo:the_annulus_type_Plateau_problem_n=2} is easily achieved.

\skpt
\begin{proof}[Proof of Theorem \ref{teo:the_annulus_type_Plateau_problem_n=2}]

\ref{1plateau} Let $\Phi\in\mathcal{P}_2(\Gamma)$, $S$, $S^+$, $S^-$ be as in the statement.
By arguing as in the proof of Theorem \ref{teorema-finale} we can find $(\sigma,\psi)\in\admclassconv$ such that $S^\pm=\mathcal{G}_{\pm\psi\res (\overline{\Om\setminus E(\sigma)})}$.
Then by Theorem \ref{teorema-finale} we have

\begin{equation}\label{mesorotta}
\mathcal{F}(\sigma,\psi)=\frac12 m_2(\Gamma)= \min_{(s,\zeta)\in\admclassconv}\mathcal{F}(s,\zeta).
\end{equation}
Hence $(\sigma,\psi)$ is a minimizer for $\mathcal F$ in $\mathcal W$; moreover by the properties of $S$ it also satisfies properties \ref{1.}--\ref{5.} of Theorem \ref{teo:structure_of_minimizers}.

\ref{2plateau}
Let $\Phi_j\in \mathcal{P}_1(\Gamma_j)$, $S_j$ for $j=1,2$, $S^+$, $S^-$ be as in the statement.
Again arguing as in the proof of Theorem \ref{teorema-finale}, we~can find $(\sigma,\psi)\in\admclassconv$ such that $S^\pm=\mathcal{G}_{\pm\psi\res(\overline{\Om\setminus E(\sigma)})}$ and \eqref{mesorotta} holds, so that $(\sigma,\psi)$ is a minimizer of $\mathcal F$ in $\mathcal W$ satisfying properties \ref{1.}--\ref{5.} of Theorem \ref{teo:structure_of_minimizers}.

\ref{3plateau} Let $(\sigma,\psi)\in\admclassconv$ be a minimizer of $\mathcal F$ in $\mathcal W$ satisfying properties \ref{1.}--\ref{5.} of Theorem \ref{teo:structure_of_minimizers}. Let also
\[
S:=\mathcal G_{\psi\res (\overline{\Om\setminus E(\sigma)})}\cup\mathcal G_{-\psi\res (\overline{\Om\setminus E(\sigma)})}.
\]
Suppose first $E(\sigma_1)\cap E(\sigma_2)=\emptyset$. Then
there is $\Phi\in \mathcal{P}_2(\Gamma)$ which is a $\mathcal{MY}$ solution to \eqref{catenoid_plateau} such that $\immclosedann=S$:
indeed, to see this, it is sufficient to apply Lemma \ref{lem_tec2}, since by Theorem \ref{teorema-finale} we have
\begin{equation}\label{cd0}
2\mathcal F(\sigma, \psi)= m_2(\Gamma).
\end{equation}
Now, suppose
$E(\sigma_1)\cap E(\sigma_2)\ne\emptyset$; then with a similar argument we can construct $\Phi_j\in\mathcal P_1(\Gamma_j)$ for $j=1,2$ solutions to \eqref{teo:the_disk_type_Plateau_problem_n=1} such that $\Phi_1(\overline B_1)\cup \Phi_2(\overline B_1)=S$.
The proof is achieved.
\end{proof}

\section{Final remarks and open problems}
\label{sec:final_comments_and_open_problems}

In this section we describe some motivations of the present study,
possible applications and related problems.
Furthermore, we~briefly comment on the hypotheses of our setting
and on possible extensions and generalizations of our results.

\subsubsection*{Connection with the Plateau problem in high codimension}
The main motivation of our study is related to the classical non-parametric Plateau problem in codimension greater than $1$. Specifically, our setting is suited for the description of the
singular part of the $L^1$-relaxation
$\mathcal A(\cdot, U)$
of the Cartesian 2-codimensional area functional
\begin{equation}
\label{eq:area_smooth_cod_2}
\int_U \sqrt{1 + \vert \grad u_1\vert^2 + \vert \grad u_2\vert^2 +
(\det\grad u)^2} ~dx, \qquad u = (u_1, u_2) \in C^1(U; \R^2),
\end{equation}
computed on nonsmooth maps.
The functional
$\mathcal A(\cdot, U)$
computed
out of $C^1(U, \R^2)$
is mostly unknown \cite{AcDa:94,GiMoSu:98_2}, up to
a few exceptions, see \cite{BePa:10,BeElPaSc:19,Scala:19,vortex}.
One of the remarkable exceptions is given by
the vortex map
$u_V:B_\ell(0) \setminus \{0\}\subset\R^2\to\R^2$,
$u_V(x):=\sfrac{x}{|x|}$:
in this case it can be proved \cite{BES1,BES2,BES3} that
\begin{align}\label{value_vortex}
\mathcal A(u_V; B_\ell(0))
= \int_{B_\ell(0)}\sqrt{1+|\nabla u_V|^2}~dx+\inf\mathcal F(\sigma,\psi),
\end{align}
where $\mathcal F(\sigma,\psi)$ is as in \eqref{def:relaxed_functional} with
$\Om=R_{2\ell}=(0,2\ell)\times(-1,1)$ and the
Dirichlet datum $\varphi :\partial R_{2\ell}\to [0,+\infty)$
is given by
\begin{equation}\label{vortex}
\varphi(z_1,z_2):=\begin{cases}
\sqrt{1-z_2^2}&\text{ on }\partial^DR_{2\ell},\\
0& \text{ on }\partial^0R_{2\ell},
\end{cases}
\end{equation}
with $\partial^DR_{2\ell}=(\{0\}\times (-1,1))\cup ([0,2\ell]\times \{-1\})\cup (\{2\ell\}\times (-1,1))$
and $\partial^0R_{2\ell}=(0,2\ell)\times \{1\}$.
Here the infimum is taken over all pairs $(\sigma,\psi)\in \Sigma\times BV(R_{2\ell})$ with $\sigma$ a unique curve in $\overline R_{2\ell}$ joining $(0,1)$ to $(2\ell,1)$
and $\psi=0$ a.e. in $E(\sigma)$.
This setting is similar to the catenoid case,
with the notable difference that the Dirichlet boundary is here extended to include the basis $(0,2\ell)\times \{-1\}$ and the free curve
is just one simple curve $\sigma$
(see Figure \ref{figura1mezzo}).

\begin{figure}
\begin{center}
\includegraphics[width=0.6\textwidth]{rettangolo1.pdf}
\caption{The domain $R_{2\ell}$ (example
of the vortex map $u_V$). The graph of $\varphi$ on $\partial^DR_{2\ell}$ is emphasized (in particular $\varphi=0$ on the lower horizontal side), together with an admissible curve $\sigma$, which in this specific case partially overlaps the Dirichlet boundary. In this example $n=1$.}\label{figura1mezzo}
\end{center}
\end{figure}

To construct a recovery sequence for \eqref{value_vortex}, it is
crucial to analyze the existence
and regularity of minimizers of $\mathcal F$.
In~particular, it is necessary to show that there is at least one sufficiently regular
minimizer $(\sigma,\psi)$.
The shape of the curve $\sigma$ and the graph of $\psi$
are related to the vertical part of a Cartesian $2$-current in
$B_\ell(0)\times \R^2 \subset \R^4$ which arises as a
limit of (the graphs of) a recovery sequence $(v_k)\subset C^1(B_\ell(0);\R^2)$ for $\mathcal A(u_V;B_\ell(0))$.

According to what happens for the catenoid, also in this case we have a dichotomy for the behaviour of minimizers $(\sigma,\psi)$.
When $\ell>0$ is small, the solution $(\sigma,\psi)$ consists of a curve $\sigma$ joining $p$ and $q$
having relative interior contained in $R_{2\ell}$, and
so that $E(\sigma)$ is convex;
at the same
time the graph of $\psi$ on $R_{2\ell}\setminus E(\sigma)$ is a sort of half-catenoid, so that if we double it considering also its symmetric with respect to the plane containing $R_{2\ell}$, it becomes a sort of catenoid spanning two unit circles and constrained to contain the segment $(0,2\ell)\times \{-1\}$. When instead $\ell$ is larger than a certain threshold, the solution reduces to two circles spanning the two unit parallel circles. Notice however that in the setting of \eqref{vortex}
on a part of the Dirichlet boundary we have $\varphi=0$. This leads to
a number of additional difficulties which must be treated
separately with approximation techniques (we
refer to \cite{BES2} for the details).

Another relevant case in which the relaxation
is known, is for the so-called \textit{triple junction function}
$u_T:B_\ell(0)\subset\R^2\to \R^2$, a map taking only
three values, vertices of an equilateral triangle of unit side-length (see \cite{BePa:10,Scala:19}). Also in this case,
in order to compute the singular contribution of $\mathcal A(u_T;B_\ell(0))$, a Cartesian Plateau problem with a partial free boundary must be solved. Following our approach, it is possible to reduce this problem to our setting.
In general,\footnote{We restrict the discussion to the $2$-dimensional case (and to codimension $2$), although this is valid for any dimension and codimension.} given $\Om\subset \R^2$ and $u\in BV(\Om;\R^2)$, the singular contribution of the relaxed area functional $\mathcal A(u;\Om)$ coincides with the mass of vertical parts in the optimal Cartesian current $T_u$ with underlying map $u$ that arises as limit of the graphs $G_{v_k}$ of a recovery sequence $v_k:\Om\to \R^2$. Generally, a few can be said on the structure and properties of these vertical parts. However, for optimal Cartesian currents $T_u$ as above, they enjoy minimality properties under suitable constraints. In the aforementioned known cases (a suitable projection\footnote{These currents live in $\Om\times \R^2$, but stands above $1$-dimensional subsets of $\Omega$, so that, with suitable techniques, they can be identified with integral currents of codimension $1$ (we refer to \cite{BES1,BES2,BES3,BePa:10,Scala:19} for more detail).} of) these vertical parts is exactly the area minimizing solution of the Cartesian Plateau type problem with partial free boundary.

We emphasize that understanding the features of vertical parts of optimal Cartesian currents for the relaxed graph area is crucial in order to detect the behaviour of the area functional. In more general settings and for general maps $u:\Om\subset \R^2\to \R^2$
only partial results are currently available, and specifically, without a
finer analysis of the singularities of these Cartesian currents
only upper and lower bounds can be obtained
(see, e.g., \cite{BSS,ScSc}, where some estimates are given for nonsmooth $\mathbb S^1$-valued maps).

\subsubsection*{Hypotheses} We assume that $\Om$ is convex.
Convexity is crucial to ensure existence of solutions even
in the classical non-parametric setting with no free boundary.
Indeed, there are examples in which $\Om$ is not convex, and a
minimizer of the area functional does not attain the Dirichlet boundary datum.

We also point out that we assume injectivity of the free-boundary
curves $\sigma_i$ (see hypothesis \ref{i'} in the introduction).
This assumption is crucial in order to define the sets
$E(\sigma_i)$ and then to solve the problem in a non-parametric form.
However, if one allows $\sigma_i$ to have self-intersections,
one can look for a disk spanning the curve $\Gamma_\sigma$
in~\eqref{gammasigma} which is not a Jordan curve anymore.
In this case we have to face a singular Plateau problem such as the one
recently studied in \cite{Creutz}
using results of \cite{LW}. Notice that in this case the curves $\sigma_i$ will be also planar and some additional hint to face this problem can be found in \cite{CS}.

\subsubsection*{Further developments} In the present analysis we have assumed that the free boundary curves are included in a plane. Of course, one may ask for domains $\Omega$ which are subset of a manifold (not necessarily a plane), leading to additional difficulties,
since the symmetrization with respect to the plane is strongly used in our arguments.

Furthermore, the correspondence between the Meeks and Yau
solutions is obtained in the special cases $n=1,2$,
although we believe that it holds also for $n\geq3$. However, due to technical difficulties which renders the setting much more involved, we~leave this generalization to future developments.

A further interesting question
is the following. Suppose that
$\partial \Om$ is smooth; then one may ask whether each free
boundary $\partial E(\sigma_j)$
is smooth up to $\partial \Om$, and moreover if there is some special
kind of contact angle condition at $\partial \Om$, due to minimality.
This question should need further investigation in the future.

The problem
considered in this paper seems not directly related to the
partial wetting phenomenon, an interesting behaviour of
soap films pointed out in \cite{Almgren}, see also
\cite{Brakke} and
\cite{2017_Bellettini_Paolini_Pasquarelli}, \cite{BPPS},
where the soap film
(typically in a non Cartesian context) does not
attain the boundary condition, leaving unwetted a part of the
wire $\Gamma$.
However, when the boundary datum $\varphi$ is allowed
to vanish (a case not covered by the results of the present paper),
as in the case of the ``catenoid'' constrained to contain the segment $[0,2\ell] \times \{-1\}$
mentioned above,
it may happen that the
singular solution consisting of
the two half-disks does not wet that segment.

We conclude this section with a couple of additional
examples which are open problems and we consider to be interesting, relating the problem
(and suitable variants) studied in this paper with the relaxation
of the area functional \eqref{eq:area_smooth_cod_2}
in dimension $2$ and codimension $2$.

Let $\widehat u:B_\ell(0)\setminus\{0\}\subset\R^2\to \mathbb S^1$ be the map defined in polar coordinates
\[
\widehat u(\rho,\theta)=e^{2i\theta},
\]
\ie the vortex map of degree $2$. Our conjecture is
that the relaxed area functional $\mathcal A(\widehat u;B_\ell(0)) $ is given by
\begin{align}\label{value_vortex2}
\int_{B_\ell(0)}\sqrt{1+|\nabla \widehat u|^2}~dx+\inf\{\mathcal F_1(\sigma,\psi_1)+\mathcal F_2(\widehat \sigma,\psi_2)\},
\end{align}
where both $\mathcal F_i$, $i=1,2$, are as the functional in \eqref{def:relaxed_functional}, but applied to different domains and variables. Specifically, $\mathcal F_1$ is applied to $\Om=R_{2\ell}$, and $\varphi$, $\sigma=(\sigma_1)$ and $\psi_1=\psi$ are exactly as in the case of $u_V$ (see \eqref{value_vortex} and \eqref{vortex}). Instead, for $\mathcal F_2$, $\Om=R_{2\ell}$ while $\widehat \sigma=(\widehat \sigma_1,\widehat\sigma_2)=( \sigma_1,\widehat\sigma_2)$, and $\varphi$ are as in the example of the catenoid in the introduction. Notice that minimizations of $\mathcal F_1$ and $\mathcal F_2$ are not independent each other, since $\sigma_1=\widehat \sigma_1$.

Another nontrivial example is given by a map $u\in BV(B_\ell(0);\R^2)$ which we assume to jump on three radii of $ B_\ell(0)$ (not necessarily
at $120^o$-degrees angles). Depending on the trace values of $u$ on these radii, we~consider the Plateau problem with partial free boundary as described below: we take as domain $\Om$ a triangle and $\sigma=(\sigma_1,\sigma_2,\sigma_3)$ are three curves in $\Om$ connecting the three pairs of vertices. Let $\varphi$ be a boundary datum on $\partial \Om$ that is null on the three vertices of $\Om$, and denote by $H(\sigma_i)$ the region enclosed between $\sigma_i$ and the side $l_i$ of $\Om$ with the same vertices. We~conjecture that the singular contribution in $\mathcal A(u;B_\ell(0))$ is related to the infimum of the quantity
\begin{multline*}\textstyle
|\Om\setminus (\bigcup_{i=1}^3H(\sigma_i))|\\
+ \sum_{i=1}^3\biggl(\int_\Om\sqrt{1+|\nabla\psi_i|^2}~dx+|D^s\psi_i|(\Om)-|\Om\setminus H(\sigma_i)|+\int_{l_i}|\psi_i-\varphi|~d\mathcal H^1\biggr).
\end{multline*}

\appendix
\section*{Appendix}
\refstepcounter{section}%\label{sec:appendix}

We recall here some classical facts about convex sets and Hausdorff distance.
If~$A,B\subset\R^2$ are nonempty, the symbol $d_H(A,B)$
stands for the Hausdorff distance between $A$ and $B$, that is
\begin{equation*}
\textstyle
d_H(A,B):=\max\{\sup_{a\in A}
\di_B(a),\,\sup_{b\in B}\di_A(b)\},
\end{equation*}
where $\di_F(\cdot)$ is the distance from the nonempty set $F\subseteq \R^2$.
If we restrict $d_H$ to the class of closed sets,
then $d_H$ defines a metric. Moreover:
\begin{enumerate}[label=(H\arabic*),wide]
\item \label{H1}
$\di_A(x)\le \di_B(x)+d_H(A,B)$
for every $x\in\R^2$;
\item \label{H2}
If $\compacts:=\{K\subset\R^2\, \text{nonempty and compact}\}$
then $(\compacts,d_H)$
is a complete metric space;
\item \label{H3}
If $A, B \in \compacts$ are convex then $d_H(A,B)=d_H(\partial A,\partial B)$;
\item \label{H4}
If $A\in \compacts$ is convex,
then there exists a sequence $(A_n)_{n} \subset \compacts$
of convex sets with boundary of class $C^\infty$ such that
$d_H(A_n,A)\to0$ as $n\to +\infty$;
\item \label{H5} Let $(A_n)_{n}$
be a sequence of nonempty closed convex sets in $\R^2$, $A\subset\R^2$ is nonempty and $d_H(A_n,A)\to0$ as $n\to +\infty$. Then $A$ is convex as well;
\item\label{H6} Let $A_n,A\in\mathcal K$ be convex such that $d_H(A_n,A)\to0$ and let $x\in \mathrm{int}(A)$; then $x\in A_n$ for all $n\in \mathbb N$ sufficiently large;
\item \label{H7} Let $A$ and $B$ be nonempty closed subsets of $\R^2$ with $d_H(A,B)=\eps$. Then $A\subset B^+_\eps$ and $B\subset A^+_\eps$
where, for all nonempty $E\subset \R^2$, we~have set
\[
E^+_\eps:=\{x\in\R^2\colon \di_E(x)\le\eps\}.
\]
\end{enumerate}
\begin{remark}
Property \ref{H1} is straightforward, while \ref{H2} is well-known. Also property \ref{H3} is easily obtained (see, e.g., \cite{W}).
Concerning property \ref{H4} we refer to, e.g., \cite[Cor.\,2]{AF}.
To see \ref{H5}, from \ref{H1} we have that $\di_{A_n} \to \di_A$ pointwise,
and therefore since $\di_{A_n}$ is convex, also $\di_A$ is convex, which implies $A$ convex.\footnote{Since $A$ is closed, it
coincides with the sublevel $\{x:d(x,A)\leq 0\}$, which is convex.}
Let us now prove \ref{H6} by contradiction;
assume that there exists a subsequence $(n_k)$ such that $\di_{A_{n_k}}(x)>0$
for all $k\in \mathbb N$; then $x\in \R^2\setminus A_{n_k}$,
$\di_{A_{n_k}}(x)=\di_{\partial A_{n_k}}(x) $, and using~\ref{H1} twice,
\begin{align*}
\di_{\partial A}(x)&\leq \di_{\partial A_{n_k}}(x)+
d_H(\partial A_{n_k},\partial A)=\di_{A_{n_k}}(x)+d_H( A_{n_k},A)\\
&\leq \di_A(x)+ 2d_H(A, A_{n_k})=2d_H(A, A_{n_k})\to0,
\end{align*}
the first equality following from \ref{H3}.
This implies $x\in \partial A$, a contradiction. Finally property \ref{H7} is immediate. Indeed if $a\in A$ then
\begin{equation*}
\mathrm{d}_B(a)\le \sup_{x\in A} \mathrm{d}_B(x)\le d_H(A,B)=\eps,
\end{equation*}
hence $a\in B_\eps^+$ and so $A\subset B_\eps^+$. In a similar way we get $B\subset A_\eps^+$.
\end{remark}

We also recall the following standard result.

\begin{lemma}\label{lm:charact_convex_sets}
Let $K\in\mathcal K$ be convex with nonempty interior.
Then there exists a $1$-periodic map $\widehat \sigma\in {\rm{Lip}}(\mathbb R;\R^2)$, injective on $[0,1)$, such that $\widehat \sigma([0,1])=\partial K$ and
\begin{equation*}
\widehat\sigma(t)=\widehat \sigma(0)+\len(\widehat\sigma)\int_0^t\widehat \gamma(s)\,ds,\quad \widehat\gamma(t)=(\cos(\widehat\theta(t)),\,\sin(\widehat\theta(t)))\quad \text{for all }\, t\in[0,1],
\end{equation*}
with $\widehat \theta\colon\R\to\R$ a non-decreasing function satisfying $\widehat \theta(t+1)-\widehat \theta(t)=2\pi$ for all $t\in\R$, and
$\len(\widehat\sigma):=\int_0^1|\widehat \sigma'(s)|ds$ the length of the curve.
\end{lemma}
Notice that $\widehat \sigma$ is differentiable a.e. in $\mathbb R$ and $\widehat \sigma'(t)=\len(\widehat\sigma)\widehat \gamma(t)$, so that the speed modulus of the curve
$|\hat\sigma'(t)|=\len(\widehat\sigma)$ is constant.

\begin{proof}
We start by approximating $K$ by convex sets with $C^\infty$ boundary. By \ref{H4} for all $n\in\mathbb{N}$ there is a convex compact set $K_n\subset\R^2$ with boundary of class $C^\infty$ and such that
$d_H(K_n,K)\to0$ as $n\to +\infty$.
For any $n\in\mathbb{N}$ we let $\widehat\sigma_n\in C^\infty(\mathbb R;\R^2)$ be a $1$\nobreakdash-periodic function injectively parametrizing $\partial K_n$ on $[0,1)$; therefore $\widehat\sigma_n([0,1])=\partial K_n$, and
\begin{equation*}
\widehat \sigma_n(t)=\widehat \sigma_n(0)+\len(\widehat\sigma_n)\int_0^t\widehat\gamma_n(s)\,ds,\quad \widehat\gamma_n(t)=(\cos(\widehat\theta_n(t)),\,\sin(\widehat\theta_n(t)))
\quad \forall t\in [0,1],
\end{equation*}
where $\widehat \theta_n\!\in\! C^\infty(\R)$ is a non-decreasing function with $\widehat\theta_n(t+1)-\widehat\theta_n(t)\!=\!2\pi$, for all \hbox{$t\!\in\! \mathbb R$}.
In view of (H2), by construction we can find $x_0\in K$, $R>r>0$ such that $B_r(x_0)\subset K_n\subset B_R(x_0)$ for all $n\in \mathbb N$, and therefore
$\mathcal{H}^1(\partial B_r(x_0))\le \len(\widehat\sigma_n)=
\H^1(\mathcal \partial K_n)\le\H^1(\partial B_R(x_0))$; where the last inequality follows since $\partial K_n=\pi_{K_n}(\partial B_R(x_0))$ and from the fact that, since $K_n$ is convex, the projection $\pi_{K_n}$ on $K_n$ does not increase the lengths, thus, up to subsequence, $\len(\widehat\sigma_n)\to \widehat m\in(0,+\infty)$ as $n\to +\infty$. Moreover, up to subsequence, we~might assume $\widehat\sigma_n(0)\to p\in \partial K$.
On the other hand, observing that
\begin{equation*}
\int_t^{t+1}|\widehat\theta'_n(s)|ds=\int_t^{t+1}\widehat\theta_n'(s)ds=2\pi, \text{ for all }t\in \mathbb R,
\end{equation*}
we have that, again up to subsequence, $\widehat\theta_n\weakstar\widehat\theta\in BV_{\mathrm{loc}}(\R)$ and pointwise (by the Helly selection principle), with $\widehat\theta$ a non-decreasing function with $\widehat \theta(t+1)-\widehat \theta(t)=2\pi$ for all $t\in \mathbb R$. We~also have $\widehat \gamma_n\weakstar\widehat \gamma$ in $BV_{\mathrm{loc}}(\mathbb R;\mathbb R^2)$ where $\widehat \gamma(t)=(\cos(\widehat \theta(t)),\,\sin(\widehat \theta(t)))$.\\
We let $\widehat \sigma\in {\rm{Lip}}(\R;\R^2)$ be the $1$-periodic curve, injective on $[0,1)$, defined as
\begin{equation}\label{sigma_boundary}
\widehat \sigma(t):=p+\widehat m\int_0^t\widehat \gamma(s)\,ds\quad \forall t\in\mathbb R.
\end{equation}
Note that $\widehat m=\len(\widehat\sigma)$.
Then clearly $\widehat \sigma_n\to \widehat \sigma$ in $W^{1,1}([0,1];\R^2)$, since
\begin{equation}
\begin{split}
\|\widehat \sigma'_n-\widehat \sigma'\|_{L^1([0,1];\R^2)}&= \int_0^1|\len(\widehat\sigma_n) \widehat \gamma_n(t)-\len(\widehat\sigma)\widehat \gamma(t)|dt\\
&\le |\len(\widehat\sigma_n)-\len(\widehat\sigma)|+\len
(\widehat\sigma)\int_0^1|\widehat \gamma_n(t)-\widehat \gamma(t)|dt\to0.
\end{split}
\label{ti-chiamo-per-nome}
\end{equation}
By the continuous embedding $W^{1,1}([0,1];\R^2)\subset C^0([0,1];\R^2)$ (and by $1$-periodicity, on $\R$) we also get $\widehat \sigma_n\to\widehat \sigma$ uniformly on $[0,1]$.
This, together with property \ref{H3} gives
\begin{equation*}
d_H(\partial K,\widehat \sigma([0,1]))\le d_H(\partial K,\partial K_n)+d_H(\widehat \sigma_n([0,1]),\widehat \sigma([0,1]))\to0,
\end{equation*}
which in turn implies $\widehat\sigma([0,1])=\partial K$. The injectivity of $\widehat \sigma$ on $[0,1)$ follows from expression \eqref{sigma_boundary}, the fact that $\widehat m>0$ and that $K$ is convex with nonempty interior.
\end{proof}

\begin{cor}\label{lm:charact_convex_sets2}
Let $K\in\mathcal K$ be convex with nonempty interior. Let $q,p$ be two distinct points on $\partial K$, and let $\arc{pq}\subset\partial K$ be the relatively open arc with endpoints $q$ and $p$ clockwise oriented.
Then there exists an injective curve $\sigma\in {\rm{Lip}}([0,1];\R^2)$ such that $\sigma((0,1))=\arc{pq}$, $\sigma(0)=q$, $\sigma(1)=p$, and
\begin{equation*}
\sigma(t)=q+\len(\sigma)\int_0^t\gamma(s)\,ds,\quad \gamma(t)=(\cos(\theta(t)),\,\sin(\theta(t)))\quad \text{for all }\, t\in[0,1],
\end{equation*}
with $\theta$ a non-decreasing function satisfying $\theta(1)-\theta(0)\le2\pi$.
\end{cor}

\begin{proof}
Lemma \ref{lm:charact_convex_sets} provides $\widehat \sigma\in \Lip ([0,1];\R^2)$ parametrizing $\partial K$. Then there are two values $t_1,t_2\in [0,1]$, $t_1<t_2$, with $q=\widehat\sigma(t_1)$ and $p=\widehat\sigma(t_2)$ so that the existence of~$\sigma$ follows by reparametrizing the interval $[t_1,t_2]$, and all the properties follow from the corresponding properties of $\widehat \sigma$.
\end{proof}

\backmatter
\bibliographystyle{jepplain+eid}
\bibliography{bellettini-et-al}
\end{document}

